\booksection{Romans}{Πρὸς Ῥωμαίους}

\chapheader{1}
\vs{1}~Παῦλος δοῦλος Χριστοῦ Ἰησοῦ, κλητὸς ἀπόστολος, ἀφωρισμένος εἰς εὐαγγέλιον θεοῦ
\vs{2}~ὃ προεπηγγείλατο διὰ τῶν προφητῶν αὐτοῦ ἐν γραφαῖς ἁγίαις
\vs{3}~περὶ τοῦ υἱοῦ αὐτοῦ, τοῦ γενομένου ἐκ σπέρματος Δαυὶδ κατὰ σάρκα,
\vs{4}~τοῦ ὁρισθέντος υἱοῦ θεοῦ ἐν δυνάμει κατὰ πνεῦμα ἁγιωσύνης ἐξ ἀναστάσεως νεκρῶν, Ἰησοῦ Χριστοῦ τοῦ κυρίου ἡμῶν,
\vs{5}~δι’ οὗ ἐλάβομεν χάριν καὶ ἀποστολὴν εἰς ὑπακοὴν πίστεως ἐν πᾶσιν τοῖς ἔθνεσιν ὑπὲρ τοῦ ὀνόματος αὐτοῦ,
\vs{6}~ἐν οἷς ἐστε καὶ ὑμεῖς κλητοὶ Ἰησοῦ Χριστοῦ,
\vs{7}~πᾶσιν τοῖς οὖσιν ἐν Ῥώμῃ ἀγαπητοῖς θεοῦ, κλητοῖς ἁγίοις· χάρις ὑμῖν καὶ εἰρήνη ἀπὸ θεοῦ πατρὸς ἡμῶν καὶ κυρίου Ἰησοῦ Χριστοῦ.
\vs{8}~Πρῶτον μὲν εὐχαριστῶ τῷ θεῷ μου διὰ Ἰησοῦ Χριστοῦ περὶ πάντων ὑμῶν, ὅτι ἡ πίστις ὑμῶν καταγγέλλεται ἐν ὅλῳ τῷ κόσμῳ.
\vs{9}~μάρτυς γάρ μού ἐστιν ὁ θεός, ᾧ λατρεύω ἐν τῷ πνεύματί μου ἐν τῷ εὐαγγελίῳ τοῦ υἱοῦ αὐτοῦ, ὡς ἀδιαλείπτως μνείαν ὑμῶν ποιοῦμαι
\vs{10}~πάντοτε ἐπὶ τῶν προσευχῶν μου, δεόμενος εἴ πως ἤδη ποτὲ εὐοδωθήσομαι ἐν τῷ θελήματι τοῦ θεοῦ ἐλθεῖν πρὸς ὑμᾶς.
\vs{11}~ἐπιποθῶ γὰρ ἰδεῖν ὑμᾶς, ἵνα τι μεταδῶ χάρισμα ὑμῖν πνευματικὸν εἰς τὸ στηριχθῆναι ὑμᾶς,
\vs{12}~τοῦτο δέ ἐστιν συμπαρακληθῆναι ἐν ὑμῖν διὰ τῆς ἐν ἀλλήλοις πίστεως ὑμῶν τε καὶ ἐμοῦ.
\vs{13}~οὐ θέλω δὲ ὑμᾶς ἀγνοεῖν, \adeA{ἀδελφοί}, ὅτι πολλάκις προεθέμην ἐλθεῖν πρὸς ὑμᾶς, καὶ ἐκωλύθην ἄχρι τοῦ δεῦρο, ἵνα τινὰ καρπὸν σχῶ καὶ ἐν ὑμῖν καθὼς καὶ ἐν τοῖς λοιποῖς ἔθνεσιν.
\vs{14}~Ἕλλησίν τε καὶ βαρβάροις, σοφοῖς τε καὶ ἀνοήτοις ὀφειλέτης εἰμί·
\vs{15}~οὕτως τὸ κατ’ ἐμὲ πρόθυμον καὶ ὑμῖν τοῖς ἐν Ῥώμῃ εὐαγγελίσασθαι.
\vs{16}~Οὐ γὰρ ἐπαισχύνομαι τὸ εὐαγγέλιον, δύναμις γὰρ θεοῦ ἐστιν εἰς σωτηρίαν παντὶ τῷ πιστεύοντι, Ἰουδαίῳ τε πρῶτον καὶ Ἕλληνι·
\vs{17}~δικαιοσύνη γὰρ θεοῦ ἐν αὐτῷ ἀποκαλύπτεται ἐκ πίστεως εἰς πίστιν, καθὼς γέγραπται· Ὁ δὲ δίκαιος ἐκ πίστεως ζήσεται.
\vs{18}~Ἀποκαλύπτεται γὰρ ὀργὴ θεοῦ ἀπ’ οὐρανοῦ ἐπὶ πᾶσαν ἀσέβειαν καὶ ἀδικίαν ἀνθρώπων τῶν τὴν ἀλήθειαν ἐν ἀδικίᾳ κατεχόντων,
\vs{19}~διότι τὸ γνωστὸν τοῦ θεοῦ φανερόν ἐστιν ἐν αὐτοῖς, ὁ θεὸς γὰρ αὐτοῖς ἐφανέρωσεν.
\vs{20}~τὰ γὰρ ἀόρατα αὐτοῦ ἀπὸ κτίσεως κόσμου τοῖς ποιήμασιν νοούμενα καθορᾶται, ἥ τε ἀΐδιος αὐτοῦ δύναμις καὶ θειότης, εἰς τὸ εἶναι αὐτοὺς ἀναπολογήτους,
\vs{21}~διότι γνόντες τὸν θεὸν οὐχ ὡς θεὸν ἐδόξασαν ἢ ηὐχαρίστησαν, ἀλλὰ ἐματαιώθησαν ἐν τοῖς διαλογισμοῖς αὐτῶν καὶ ἐσκοτίσθη ἡ ἀσύνετος αὐτῶν καρδία·
\vs{22}~φάσκοντες εἶναι σοφοὶ ἐμωράνθησαν,
\vs{23}~καὶ ἤλλαξαν τὴν δόξαν τοῦ ἀφθάρτου θεοῦ ἐν ὁμοιώματι εἰκόνος φθαρτοῦ ἀνθρώπου καὶ πετεινῶν καὶ τετραπόδων καὶ ἑρπετῶν.
\vs{24}~Διὸ παρέδωκεν αὐτοὺς ὁ θεὸς ἐν ταῖς ἐπιθυμίαις τῶν καρδιῶν αὐτῶν εἰς ἀκαθαρσίαν τοῦ ἀτιμάζεσθαι τὰ σώματα αὐτῶν ἐν αὐτοῖς,
\vs{25}~οἵτινες μετήλλαξαν τὴν ἀλήθειαν τοῦ θεοῦ ἐν τῷ ψεύδει, καὶ ἐσεβάσθησαν καὶ ἐλάτρευσαν τῇ κτίσει παρὰ τὸν κτίσαντα, ὅς ἐστιν εὐλογητὸς εἰς τοὺς αἰῶνας· ἀμήν.
\vs{26}~Διὰ τοῦτο παρέδωκεν αὐτοὺς ὁ θεὸς εἰς πάθη ἀτιμίας· αἵ τε γὰρ θήλειαι αὐτῶν μετήλλαξαν τὴν φυσικὴν χρῆσιν εἰς τὴν παρὰ φύσιν,
\vs{27}~ὁμοίως τε καὶ οἱ ἄρσενες ἀφέντες τὴν φυσικὴν χρῆσιν τῆς θηλείας ἐξεκαύθησαν ἐν τῇ ὀρέξει αὐτῶν εἰς ἀλλήλους, ἄρσενες ἐν ἄρσεσιν τὴν ἀσχημοσύνην κατεργαζόμενοι καὶ τὴν ἀντιμισθίαν ἣν ἔδει τῆς πλάνης αὐτῶν ἐν ἑαυτοῖς ἀπολαμβάνοντες.
\vs{28}~Καὶ καθὼς οὐκ ἐδοκίμασαν τὸν θεὸν ἔχειν ἐν ἐπιγνώσει, παρέδωκεν αὐτοὺς ὁ θεὸς εἰς ἀδόκιμον νοῦν, ποιεῖν τὰ μὴ καθήκοντα,
\vs{29}~πεπληρωμένους πάσῃ ἀδικίᾳ πονηρίᾳ πλεονεξίᾳ κακίᾳ, μεστοὺς φθόνου φόνου ἔριδος δόλου κακοηθείας, ψιθυριστάς,
\vs{30}~καταλάλους, θεοστυγεῖς, ὑβριστάς, ὑπερηφάνους, ἀλαζόνας, ἐφευρετὰς κακῶν, γονεῦσιν ἀπειθεῖς,
\vs{31}~ἀσυνέτους, ἀσυνθέτους, ἀστόργους, ἀνελεήμονας·
\vs{32}~οἵτινες τὸ δικαίωμα τοῦ θεοῦ ἐπιγνόντες, ὅτι οἱ τὰ τοιαῦτα πράσσοντες ἄξιοι θανάτου εἰσίν, οὐ μόνον αὐτὰ ποιοῦσιν ἀλλὰ καὶ συνευδοκοῦσιν τοῖς πράσσουσιν.

\chapheader{2}
\vs{1}~Διὸ ἀναπολόγητος εἶ, ὦ ἄνθρωπε πᾶς ὁ κρίνων· ἐν ᾧ γὰρ κρίνεις τὸν ἕτερον, σεαυτὸν κατακρίνεις, τὰ γὰρ αὐτὰ πράσσεις ὁ κρίνων·
\vs{2}~οἴδαμεν δὲ ὅτι τὸ κρίμα τοῦ θεοῦ ἐστιν κατὰ ἀλήθειαν ἐπὶ τοὺς τὰ τοιαῦτα πράσσοντας.
\vs{3}~λογίζῃ δὲ τοῦτο, ὦ ἄνθρωπε ὁ κρίνων τοὺς τὰ τοιαῦτα πράσσοντας καὶ ποιῶν αὐτά, ὅτι σὺ ἐκφεύξῃ τὸ κρίμα τοῦ θεοῦ;
\vs{4}~ἢ τοῦ πλούτου τῆς χρηστότητος αὐτοῦ καὶ τῆς ἀνοχῆς καὶ τῆς μακροθυμίας καταφρονεῖς, ἀγνοῶν ὅτι τὸ χρηστὸν τοῦ θεοῦ εἰς μετάνοιάν σε ἄγει;
\vs{5}~κατὰ δὲ τὴν σκληρότητά σου καὶ ἀμετανόητον καρδίαν θησαυρίζεις σεαυτῷ ὀργὴν ἐν ἡμέρᾳ ὀργῆς καὶ ἀποκαλύψεως δικαιοκρισίας τοῦ θεοῦ,
\vs{6}~ὃς ἀποδώσει ἑκάστῳ κατὰ τὰ ἔργα αὐτοῦ·
\vs{7}~τοῖς μὲν καθ’ ὑπομονὴν ἔργου ἀγαθοῦ δόξαν καὶ τιμὴν καὶ ἀφθαρσίαν ζητοῦσιν ζωὴν αἰώνιον·
\vs{8}~τοῖς δὲ ἐξ ἐριθείας καὶ ἀπειθοῦσι τῇ ἀληθείᾳ πειθομένοις δὲ τῇ ἀδικίᾳ ὀργὴ καὶ θυμός,
\vs{9}~θλῖψις καὶ στενοχωρία, ἐπὶ πᾶσαν ψυχὴν ἀνθρώπου τοῦ κατεργαζομένου τὸ κακόν, Ἰουδαίου τε πρῶτον καὶ Ἕλληνος·
\vs{10}~δόξα δὲ καὶ τιμὴ καὶ εἰρήνη παντὶ τῷ ἐργαζομένῳ τὸ ἀγαθόν, Ἰουδαίῳ τε πρῶτον καὶ Ἕλληνι·
\vs{11}~οὐ γάρ ἐστιν προσωπολημψία παρὰ τῷ θεῷ.
\vs{12}~Ὅσοι γὰρ ἀνόμως ἥμαρτον, ἀνόμως καὶ ἀπολοῦνται· καὶ ὅσοι ἐν νόμῳ ἥμαρτον, διὰ νόμου κριθήσονται·
\vs{13}~οὐ γὰρ οἱ ἀκροαταὶ νόμου δίκαιοι παρὰ τῷ θεῷ, ἀλλ’ οἱ ποιηταὶ νόμου δικαιωθήσονται.
\vs{14}~ὅταν γὰρ ἔθνη τὰ μὴ νόμον ἔχοντα φύσει τὰ τοῦ νόμου ποιῶσιν, οὗτοι νόμον μὴ ἔχοντες ἑαυτοῖς εἰσιν νόμος·
\vs{15}~οἵτινες ἐνδείκνυνται τὸ ἔργον τοῦ νόμου γραπτὸν ἐν ταῖς καρδίαις αὐτῶν, συμμαρτυρούσης αὐτῶν τῆς συνειδήσεως καὶ μεταξὺ ἀλλήλων τῶν λογισμῶν κατηγορούντων ἢ καὶ ἀπολογουμένων,
\vs{16}~ἐν ἡμέρᾳ ὅτε κρίνει ὁ θεὸς τὰ κρυπτὰ τῶν ἀνθρώπων κατὰ τὸ εὐαγγέλιόν μου διὰ Χριστοῦ Ἰησοῦ.
\vs{17}~Εἰ δὲ σὺ Ἰουδαῖος ἐπονομάζῃ καὶ ἐπαναπαύῃ νόμῳ καὶ καυχᾶσαι ἐν θεῷ
\vs{18}~καὶ γινώσκεις τὸ θέλημα καὶ δοκιμάζεις τὰ διαφέροντα κατηχούμενος ἐκ τοῦ νόμου,
\vs{19}~πέποιθάς τε σεαυτὸν ὁδηγὸν εἶναι τυφλῶν, φῶς τῶν ἐν σκότει,
\vs{20}~παιδευτὴν ἀφρόνων, διδάσκαλον νηπίων, ἔχοντα τὴν μόρφωσιν τῆς γνώσεως καὶ τῆς ἀληθείας ἐν τῷ νόμῳ—
\vs{21}~ὁ οὖν διδάσκων ἕτερον σεαυτὸν οὐ διδάσκεις; ὁ κηρύσσων μὴ κλέπτειν κλέπτεις;
\vs{22}~ὁ λέγων μὴ μοιχεύειν μοιχεύεις; ὁ βδελυσσόμενος τὰ εἴδωλα ἱεροσυλεῖς;
\vs{23}~ὃς ἐν νόμῳ καυχᾶσαι, διὰ τῆς παραβάσεως τοῦ νόμου τὸν θεὸν ἀτιμάζεις;
\vs{24}~τὸ γὰρ ὄνομα τοῦ θεοῦ δι’ ὑμᾶς βλασφημεῖται ἐν τοῖς ἔθνεσιν, καθὼς γέγραπται.
\vs{25}~Περιτομὴ μὲν γὰρ ὠφελεῖ ἐὰν νόμον πράσσῃς· ἐὰν δὲ παραβάτης νόμου ᾖς, ἡ περιτομή σου ἀκροβυστία γέγονεν.
\vs{26}~ἐὰν οὖν ἡ ἀκροβυστία τὰ δικαιώματα τοῦ νόμου φυλάσσῃ, οὐχ ἡ ἀκροβυστία αὐτοῦ εἰς περιτομὴν λογισθήσεται;
\vs{27}~καὶ κρινεῖ ἡ ἐκ φύσεως ἀκροβυστία τὸν νόμον τελοῦσα σὲ τὸν διὰ γράμματος καὶ περιτομῆς παραβάτην νόμου.
\vs{28}~οὐ γὰρ ὁ ἐν τῷ φανερῷ Ἰουδαῖός ἐστιν, οὐδὲ ἡ ἐν τῷ φανερῷ ἐν σαρκὶ περιτομή·
\vs{29}~ἀλλ’ ὁ ἐν τῷ κρυπτῷ Ἰουδαῖος, καὶ περιτομὴ καρδίας ἐν πνεύματι οὐ γράμματι, οὗ ὁ ἔπαινος οὐκ ἐξ ἀνθρώπων ἀλλ’ ἐκ τοῦ θεοῦ.

\chapheader{3}
\vs{1}~Τί οὖν τὸ περισσὸν τοῦ Ἰουδαίου, ἢ τίς ἡ ὠφέλεια τῆς περιτομῆς;
\vs{2}~πολὺ κατὰ πάντα τρόπον. πρῶτον μὲν γὰρ ὅτι ἐπιστεύθησαν τὰ λόγια τοῦ θεοῦ.
\vs{3}~τί γάρ; εἰ ἠπίστησάν τινες, μὴ ἡ ἀπιστία αὐτῶν τὴν πίστιν τοῦ θεοῦ καταργήσει;
\vs{4}~μὴ γένοιτο· γινέσθω δὲ ὁ θεὸς ἀληθής, πᾶς δὲ ἄνθρωπος ψεύστης, καθὼς γέγραπται· Ὅπως ἂν δικαιωθῇς ἐν τοῖς λόγοις σου καὶ νικήσεις ἐν τῷ κρίνεσθαί σε.
\vs{5}~εἰ δὲ ἡ ἀδικία ἡμῶν θεοῦ δικαιοσύνην συνίστησιν, τί ἐροῦμεν; μὴ ἄδικος ὁ θεὸς ὁ ἐπιφέρων τὴν ὀργήν; κατὰ ἄνθρωπον λέγω.
\vs{6}~μὴ γένοιτο· ἐπεὶ πῶς κρινεῖ ὁ θεὸς τὸν κόσμον;
\vs{7}~εἰ δὲ ἡ ἀλήθεια τοῦ θεοῦ ἐν τῷ ἐμῷ ψεύσματι ἐπερίσσευσεν εἰς τὴν δόξαν αὐτοῦ, τί ἔτι κἀγὼ ὡς ἁμαρτωλὸς κρίνομαι,
\vs{8}~καὶ μὴ καθὼς βλασφημούμεθα καὶ καθώς φασίν τινες ἡμᾶς λέγειν ὅτι Ποιήσωμεν τὰ κακὰ ἵνα ἔλθῃ τὰ ἀγαθά; ὧν τὸ κρίμα ἔνδικόν ἐστιν.
\vs{9}~Τί οὖν; προεχόμεθα; οὐ πάντως, προῃτιασάμεθα γὰρ Ἰουδαίους τε καὶ Ἕλληνας πάντας ὑφ’ ἁμαρτίαν εἶναι,
\vs{10}~καθὼς γέγραπται ὅτι Οὐκ ἔστιν δίκαιος οὐδὲ εἷς,
\vs{11}~οὐκ ἔστιν ὁ συνίων, οὐκ ἔστιν ὁ ἐκζητῶν τὸν θεόν·
\vs{12}~πάντες ἐξέκλιναν, ἅμα ἠχρεώθησαν· οὐκ ἔστιν ποιῶν χρηστότητα, οὐκ ἔστιν ἕως ἑνός.
\vs{13}~τάφος ἀνεῳγμένος ὁ λάρυγξ αὐτῶν, ταῖς γλώσσαις αὐτῶν ἐδολιοῦσαν, ἰὸς ἀσπίδων ὑπὸ τὰ χείλη αὐτῶν,
\vs{14}~ὧν τὸ στόμα ἀρᾶς καὶ πικρίας γέμει·
\vs{15}~ὀξεῖς οἱ πόδες αὐτῶν ἐκχέαι αἷμα,
\vs{16}~σύντριμμα καὶ ταλαιπωρία ἐν ταῖς ὁδοῖς αὐτῶν,
\vs{17}~καὶ ὁδὸν εἰρήνης οὐκ ἔγνωσαν.
\vs{18}~οὐκ ἔστιν φόβος θεοῦ ἀπέναντι τῶν ὀφθαλμῶν αὐτῶν.
\vs{19}~Οἴδαμεν δὲ ὅτι ὅσα ὁ νόμος λέγει τοῖς ἐν τῷ νόμῳ λαλεῖ, ἵνα πᾶν στόμα φραγῇ καὶ ὑπόδικος γένηται πᾶς ὁ κόσμος τῷ θεῷ·
\vs{20}~διότι ἐξ ἔργων νόμου οὐ δικαιωθήσεται πᾶσα σὰρξ ἐνώπιον αὐτοῦ, διὰ γὰρ νόμου ἐπίγνωσις ἁμαρτίας.
\vs{21}~Νυνὶ δὲ χωρὶς νόμου δικαιοσύνη θεοῦ πεφανέρωται, μαρτυρουμένη ὑπὸ τοῦ νόμου καὶ τῶν προφητῶν,
\vs{22}~δικαιοσύνη δὲ θεοῦ διὰ πίστεως Ἰησοῦ Χριστοῦ, εἰς πάντας τοὺς πιστεύοντας, οὐ γάρ ἐστιν διαστολή.
\vs{23}~πάντες γὰρ ἥμαρτον καὶ ὑστεροῦνται τῆς δόξης τοῦ θεοῦ,
\vs{24}~δικαιούμενοι δωρεὰν τῇ αὐτοῦ χάριτι διὰ τῆς ἀπολυτρώσεως τῆς ἐν Χριστῷ Ἰησοῦ·
\vs{25}~ὃν προέθετο ὁ θεὸς ἱλαστήριον διὰ πίστεως ἐν τῷ αὐτοῦ αἵματι εἰς ἔνδειξιν τῆς δικαιοσύνης αὐτοῦ διὰ τὴν πάρεσιν τῶν προγεγονότων ἁμαρτημάτων
\vs{26}~ἐν τῇ ἀνοχῇ τοῦ θεοῦ, πρὸς τὴν ἔνδειξιν τῆς δικαιοσύνης αὐτοῦ ἐν τῷ νῦν καιρῷ, εἰς τὸ εἶναι αὐτὸν δίκαιον καὶ δικαιοῦντα τὸν ἐκ πίστεως Ἰησοῦ.
\vs{27}~Ποῦ οὖν ἡ καύχησις; ἐξεκλείσθη. διὰ ποίου νόμου; τῶν ἔργων; οὐχί, ἀλλὰ διὰ νόμου πίστεως.
\vs{28}~λογιζόμεθα γὰρ δικαιοῦσθαι πίστει ἄνθρωπον χωρὶς ἔργων νόμου.
\vs{29}~ἢ Ἰουδαίων ὁ θεὸς μόνον; οὐχὶ καὶ ἐθνῶν; ναὶ καὶ ἐθνῶν,
\vs{30}~εἴπερ εἷς ὁ θεός, ὃς δικαιώσει περιτομὴν ἐκ πίστεως καὶ ἀκροβυστίαν διὰ τῆς πίστεως.
\vs{31}~νόμον οὖν καταργοῦμεν διὰ τῆς πίστεως; μὴ γένοιτο, ἀλλὰ νόμον ἱστάνομεν.

\chapheader{4}
\vs{1}~Τί οὖν ἐροῦμεν εὑρηκέναι Ἀβραὰμ τὸν προπάτορα ἡμῶν κατὰ σάρκα;
\vs{2}~εἰ γὰρ Ἀβραὰμ ἐξ ἔργων ἐδικαιώθη, ἔχει καύχημα· ἀλλ’ οὐ πρὸς θεόν,
\vs{3}~τί γὰρ ἡ γραφὴ λέγει; Ἐπίστευσεν δὲ Ἀβραὰμ τῷ θεῷ καὶ ἐλογίσθη αὐτῷ εἰς δικαιοσύνην.
\vs{4}~τῷ δὲ ἐργαζομένῳ ὁ μισθὸς οὐ λογίζεται κατὰ χάριν ἀλλὰ κατὰ ὀφείλημα·
\vs{5}~τῷ δὲ μὴ ἐργαζομένῳ, πιστεύοντι δὲ ἐπὶ τὸν δικαιοῦντα τὸν ἀσεβῆ, λογίζεται ἡ πίστις αὐτοῦ εἰς δικαιοσύνην,
\vs{6}~καθάπερ καὶ Δαυὶδ λέγει τὸν μακαρισμὸν τοῦ ἀνθρώπου ᾧ ὁ θεὸς λογίζεται δικαιοσύνην χωρὶς ἔργων·
\vs{7}~Μακάριοι ὧν ἀφέθησαν αἱ ἀνομίαι καὶ ὧν ἐπεκαλύφθησαν αἱ ἁμαρτίαι,
\vs{8}~μακάριος ἀνὴρ οὗ οὐ μὴ λογίσηται κύριος ἁμαρτίαν.
\vs{9}~Ὁ μακαρισμὸς οὖν οὗτος ἐπὶ τὴν περιτομὴν ἢ καὶ ἐπὶ τὴν ἀκροβυστίαν; λέγομεν γάρ· Ἐλογίσθη τῷ Ἀβραὰμ ἡ πίστις εἰς δικαιοσύνην.
\vs{10}~πῶς οὖν ἐλογίσθη; ἐν περιτομῇ ὄντι ἢ ἐν ἀκροβυστίᾳ; οὐκ ἐν περιτομῇ ἀλλ’ ἐν ἀκροβυστίᾳ·
\vs{11}~καὶ σημεῖον ἔλαβεν περιτομῆς, σφραγῖδα τῆς δικαιοσύνης τῆς πίστεως τῆς ἐν τῇ ἀκροβυστίᾳ, εἰς τὸ εἶναι αὐτὸν πατέρα πάντων τῶν πιστευόντων δι’ ἀκροβυστίας, εἰς τὸ λογισθῆναι αὐτοῖς τὴν δικαιοσύνην,
\vs{12}~καὶ πατέρα περιτομῆς τοῖς οὐκ ἐκ περιτομῆς μόνον ἀλλὰ καὶ τοῖς στοιχοῦσιν τοῖς ἴχνεσιν τῆς ἐν ἀκροβυστίᾳ πίστεως τοῦ πατρὸς ἡμῶν Ἀβραάμ.
\vs{13}~Οὐ γὰρ διὰ νόμου ἡ ἐπαγγελία τῷ Ἀβραὰμ ἢ τῷ σπέρματι αὐτοῦ, τὸ κληρονόμον αὐτὸν εἶναι κόσμου, ἀλλὰ διὰ δικαιοσύνης πίστεως·
\vs{14}~εἰ γὰρ οἱ ἐκ νόμου κληρονόμοι, κεκένωται ἡ πίστις καὶ κατήργηται ἡ ἐπαγγελία·
\vs{15}~ὁ γὰρ νόμος ὀργὴν κατεργάζεται, οὗ δὲ οὐκ ἔστιν νόμος, οὐδὲ παράβασις.
\vs{16}~Διὰ τοῦτο ἐκ πίστεως, ἵνα κατὰ χάριν, εἰς τὸ εἶναι βεβαίαν τὴν ἐπαγγελίαν παντὶ τῷ σπέρματι, οὐ τῷ ἐκ τοῦ νόμου μόνον ἀλλὰ καὶ τῷ ἐκ πίστεως Ἀβραάμ (ὅς ἐστιν πατὴρ πάντων ἡμῶν,
\vs{17}~καθὼς γέγραπται ὅτι Πατέρα πολλῶν ἐθνῶν τέθεικά σε), κατέναντι οὗ ἐπίστευσεν θεοῦ τοῦ ζῳοποιοῦντος τοὺς νεκροὺς καὶ καλοῦντος τὰ μὴ ὄντα ὡς ὄντα·
\vs{18}~ὃς παρ’ ἐλπίδα ἐπ’ ἐλπίδι ἐπίστευσεν εἰς τὸ γενέσθαι αὐτὸν πατέρα πολλῶν ἐθνῶν κατὰ τὸ εἰρημένον· Οὕτως ἔσται τὸ σπέρμα σου·
\vs{19}~καὶ μὴ ἀσθενήσας τῇ πίστει κατενόησεν τὸ ἑαυτοῦ σῶμα νενεκρωμένον, ἑκατονταετής που ὑπάρχων, καὶ τὴν νέκρωσιν τῆς μήτρας Σάρρας,
\vs{20}~εἰς δὲ τὴν ἐπαγγελίαν τοῦ θεοῦ οὐ διεκρίθη τῇ ἀπιστίᾳ ἀλλὰ ἐνεδυναμώθη τῇ πίστει, δοὺς δόξαν τῷ θεῷ
\vs{21}~καὶ πληροφορηθεὶς ὅτι ὃ ἐπήγγελται δυνατός ἐστιν καὶ ποιῆσαι.
\vs{22}~διὸ ἐλογίσθη αὐτῷ εἰς δικαιοσύνην.
\vs{23}~Οὐκ ἐγράφη δὲ δι’ αὐτὸν μόνον ὅτι ἐλογίσθη αὐτῷ,
\vs{24}~ἀλλὰ καὶ δι’ ἡμᾶς οἷς μέλλει λογίζεσθαι, τοῖς πιστεύουσιν ἐπὶ τὸν ἐγείραντα Ἰησοῦν τὸν κύριον ἡμῶν ἐκ νεκρῶν,
\vs{25}~ὃς παρεδόθη διὰ τὰ παραπτώματα ἡμῶν καὶ ἠγέρθη διὰ τὴν δικαίωσιν ἡμῶν.

\chapheader{5}
\vs{1}~Δικαιωθέντες οὖν ἐκ πίστεως εἰρήνην ἔχομεν πρὸς τὸν θεὸν διὰ τοῦ κυρίου ἡμῶν Ἰησοῦ Χριστοῦ,
\vs{2}~δι’ οὗ καὶ τὴν προσαγωγὴν ἐσχήκαμεν τῇ πίστει εἰς τὴν χάριν ταύτην ἐν ᾗ ἑστήκαμεν, καὶ καυχώμεθα ἐπ’ ἐλπίδι τῆς δόξης τοῦ θεοῦ·
\vs{3}~οὐ μόνον δέ, ἀλλὰ καὶ καυχώμεθα ἐν ταῖς θλίψεσιν, εἰδότες ὅτι ἡ θλῖψις ὑπομονὴν κατεργάζεται,
\vs{4}~ἡ δὲ ὑπομονὴ δοκιμήν, ἡ δὲ δοκιμὴ ἐλπίδα.
\vs{5}~ἡ δὲ ἐλπὶς οὐ καταισχύνει· ὅτι ἡ ἀγάπη τοῦ θεοῦ ἐκκέχυται ἐν ταῖς καρδίαις ἡμῶν διὰ πνεύματος ἁγίου τοῦ δοθέντος ἡμῖν.
\vs{6}~Ἔτι γὰρ Χριστὸς ὄντων ἡμῶν ἀσθενῶν ἔτι κατὰ καιρὸν ὑπὲρ ἀσεβῶν ἀπέθανεν.
\vs{7}~μόλις γὰρ ὑπὲρ δικαίου τις ἀποθανεῖται· ὑπὲρ γὰρ τοῦ ἀγαθοῦ τάχα τις καὶ τολμᾷ ἀποθανεῖν·
\vs{8}~συνίστησιν δὲ τὴν ἑαυτοῦ ἀγάπην εἰς ἡμᾶς ὁ θεὸς ὅτι ἔτι ἁμαρτωλῶν ὄντων ἡμῶν Χριστὸς ὑπὲρ ἡμῶν ἀπέθανεν.
\vs{9}~πολλῷ οὖν μᾶλλον δικαιωθέντες νῦν ἐν τῷ αἵματι αὐτοῦ σωθησόμεθα δι’ αὐτοῦ ἀπὸ τῆς ὀργῆς.
\vs{10}~εἰ γὰρ ἐχθροὶ ὄντες κατηλλάγημεν τῷ θεῷ διὰ τοῦ θανάτου τοῦ υἱοῦ αὐτοῦ, πολλῷ μᾶλλον καταλλαγέντες σωθησόμεθα ἐν τῇ ζωῇ αὐτοῦ·
\vs{11}~οὐ μόνον δέ, ἀλλὰ καὶ καυχώμενοι ἐν τῷ θεῷ διὰ τοῦ κυρίου ἡμῶν Ἰησοῦ Χριστοῦ, δι’ οὗ νῦν τὴν καταλλαγὴν ἐλάβομεν.
\vs{12}~Διὰ τοῦτο ὥσπερ δι’ ἑνὸς ἀνθρώπου ἡ ἁμαρτία εἰς τὸν κόσμον εἰσῆλθεν καὶ διὰ τῆς ἁμαρτίας ὁ θάνατος, καὶ οὕτως εἰς πάντας ἀνθρώπους ὁ θάνατος διῆλθεν ἐφ’ ᾧ πάντες ἥμαρτον—
\vs{13}~ἄχρι γὰρ νόμου ἁμαρτία ἦν ἐν κόσμῳ, ἁμαρτία δὲ οὐκ ἐλλογεῖται μὴ ὄντος νόμου,
\vs{14}~ἀλλὰ ἐβασίλευσεν ὁ θάνατος ἀπὸ Ἀδὰμ μέχρι Μωϋσέως καὶ ἐπὶ τοὺς μὴ ἁμαρτήσαντας ἐπὶ τῷ ὁμοιώματι τῆς παραβάσεως Ἀδάμ, ὅς ἐστιν τύπος τοῦ μέλλοντος.
\vs{15}~Ἀλλ’ οὐχ ὡς τὸ παράπτωμα, οὕτως καὶ τὸ χάρισμα· εἰ γὰρ τῷ τοῦ ἑνὸς παραπτώματι οἱ πολλοὶ ἀπέθανον, πολλῷ μᾶλλον ἡ χάρις τοῦ θεοῦ καὶ ἡ δωρεὰ ἐν χάριτι τῇ τοῦ ἑνὸς ἀνθρώπου Ἰησοῦ Χριστοῦ εἰς τοὺς πολλοὺς ἐπερίσσευσεν.
\vs{16}~καὶ οὐχ ὡς δι’ ἑνὸς ἁμαρτήσαντος τὸ δώρημα· τὸ μὲν γὰρ κρίμα ἐξ ἑνὸς εἰς κατάκριμα, τὸ δὲ χάρισμα ἐκ πολλῶν παραπτωμάτων εἰς δικαίωμα.
\vs{17}~εἰ γὰρ τῷ τοῦ ἑνὸς παραπτώματι ὁ θάνατος ἐβασίλευσεν διὰ τοῦ ἑνός, πολλῷ μᾶλλον οἱ τὴν περισσείαν τῆς χάριτος καὶ τῆς δωρεᾶς τῆς δικαιοσύνης λαμβάνοντες ἐν ζωῇ βασιλεύσουσιν διὰ τοῦ ἑνὸς Ἰησοῦ Χριστοῦ.
\vs{18}~Ἄρα οὖν ὡς δι’ ἑνὸς παραπτώματος εἰς πάντας ἀνθρώπους εἰς κατάκριμα, οὕτως καὶ δι’ ἑνὸς δικαιώματος εἰς πάντας ἀνθρώπους εἰς δικαίωσιν ζωῆς·
\vs{19}~ὥσπερ γὰρ διὰ τῆς παρακοῆς τοῦ ἑνὸς ἀνθρώπου ἁμαρτωλοὶ κατεστάθησαν οἱ πολλοί, οὕτως καὶ διὰ τῆς ὑπακοῆς τοῦ ἑνὸς δίκαιοι κατασταθήσονται οἱ πολλοί.
\vs{20}~νόμος δὲ παρεισῆλθεν ἵνα πλεονάσῃ τὸ παράπτωμα· οὗ δὲ ἐπλεόνασεν ἡ ἁμαρτία, ὑπερεπερίσσευσεν ἡ χάρις,
\vs{21}~ἵνα ὥσπερ ἐβασίλευσεν ἡ ἁμαρτία ἐν τῷ θανάτῳ, οὕτως καὶ ἡ χάρις βασιλεύσῃ διὰ δικαιοσύνης εἰς ζωὴν αἰώνιον διὰ Ἰησοῦ Χριστοῦ τοῦ κυρίου ἡμῶν.

\chapheader{6}
\vs{1}~Τί οὖν ἐροῦμεν; ἐπιμένωμεν τῇ ἁμαρτίᾳ, ἵνα ἡ χάρις πλεονάσῃ;
\vs{2}~μὴ γένοιτο· οἵτινες ἀπεθάνομεν τῇ ἁμαρτίᾳ, πῶς ἔτι ζήσομεν ἐν αὐτῇ;
\vs{3}~ἢ ἀγνοεῖτε ὅτι ὅσοι ἐβαπτίσθημεν εἰς Χριστὸν Ἰησοῦν εἰς τὸν θάνατον αὐτοῦ ἐβαπτίσθημεν;
\vs{4}~συνετάφημεν οὖν αὐτῷ διὰ τοῦ βαπτίσματος εἰς τὸν θάνατον, ἵνα ὥσπερ ἠγέρθη Χριστὸς ἐκ νεκρῶν διὰ τῆς δόξης τοῦ πατρός, οὕτως καὶ ἡμεῖς ἐν καινότητι ζωῆς περιπατήσωμεν.
\vs{5}~Εἰ γὰρ σύμφυτοι γεγόναμεν τῷ ὁμοιώματι τοῦ θανάτου αὐτοῦ, ἀλλὰ καὶ τῆς ἀναστάσεως ἐσόμεθα·
\vs{6}~τοῦτο γινώσκοντες ὅτι ὁ παλαιὸς ἡμῶν ἄνθρωπος συνεσταυρώθη, ἵνα καταργηθῇ τὸ σῶμα τῆς ἁμαρτίας, τοῦ μηκέτι δουλεύειν ἡμᾶς τῇ ἁμαρτίᾳ,
\vs{7}~ὁ γὰρ ἀποθανὼν δεδικαίωται ἀπὸ τῆς ἁμαρτίας.
\vs{8}~εἰ δὲ ἀπεθάνομεν σὺν Χριστῷ, πιστεύομεν ὅτι καὶ συζήσομεν αὐτῷ·
\vs{9}~εἰδότες ὅτι Χριστὸς ἐγερθεὶς ἐκ νεκρῶν οὐκέτι ἀποθνῄσκει, θάνατος αὐτοῦ οὐκέτι κυριεύει·
\vs{10}~ὃ γὰρ ἀπέθανεν, τῇ ἁμαρτίᾳ ἀπέθανεν ἐφάπαξ· ὃ δὲ ζῇ, ζῇ τῷ θεῷ.
\vs{11}~οὕτως καὶ ὑμεῖς λογίζεσθε ἑαυτοὺς εἶναι νεκροὺς μὲν τῇ ἁμαρτίᾳ ζῶντας δὲ τῷ θεῷ ἐν Χριστῷ Ἰησοῦ.
\vs{12}~Μὴ οὖν βασιλευέτω ἡ ἁμαρτία ἐν τῷ θνητῷ ὑμῶν σώματι εἰς τὸ ὑπακούειν ταῖς ἐπιθυμίαις αὐτοῦ,
\vs{13}~μηδὲ παριστάνετε τὰ μέλη ὑμῶν ὅπλα ἀδικίας τῇ ἁμαρτίᾳ, ἀλλὰ παραστήσατε ἑαυτοὺς τῷ θεῷ ὡσεὶ ἐκ νεκρῶν ζῶντας καὶ τὰ μέλη ὑμῶν ὅπλα δικαιοσύνης τῷ θεῷ.
\vs{14}~ἁμαρτία γὰρ ὑμῶν οὐ κυριεύσει, οὐ γάρ ἐστε ὑπὸ νόμον ἀλλὰ ὑπὸ χάριν.
\vs{15}~Τί οὖν; ἁμαρτήσωμεν ὅτι οὐκ ἐσμὲν ὑπὸ νόμον ἀλλὰ ὑπὸ χάριν; μὴ γένοιτο·
\vs{16}~οὐκ οἴδατε ὅτι ᾧ παριστάνετε ἑαυτοὺς δούλους εἰς ὑπακοήν, δοῦλοί ἐστε ᾧ ὑπακούετε, ἤτοι ἁμαρτίας εἰς θάνατον ἢ ὑπακοῆς εἰς δικαιοσύνην;
\vs{17}~χάρις δὲ τῷ θεῷ ὅτι ἦτε δοῦλοι τῆς ἁμαρτίας ὑπηκούσατε δὲ ἐκ καρδίας εἰς ὃν παρεδόθητε τύπον διδαχῆς,
\vs{18}~ἐλευθερωθέντες δὲ ἀπὸ τῆς ἁμαρτίας ἐδουλώθητε τῇ δικαιοσύνῃ·
\vs{19}~ἀνθρώπινον λέγω διὰ τὴν ἀσθένειαν τῆς σαρκὸς ὑμῶν· ὥσπερ γὰρ παρεστήσατε τὰ μέλη ὑμῶν δοῦλα τῇ ἀκαθαρσίᾳ καὶ τῇ ἀνομίᾳ εἰς τὴν ἀνομίαν, οὕτως νῦν παραστήσατε τὰ μέλη ὑμῶν δοῦλα τῇ δικαιοσύνῃ εἰς ἁγιασμόν.
\vs{20}~Ὅτε γὰρ δοῦλοι ἦτε τῆς ἁμαρτίας, ἐλεύθεροι ἦτε τῇ δικαιοσύνῃ.
\vs{21}~τίνα οὖν καρπὸν εἴχετε τότε ἐφ’ οἷς νῦν ἐπαισχύνεσθε; τὸ γὰρ τέλος ἐκείνων θάνατος·
\vs{22}~νυνὶ δέ, ἐλευθερωθέντες ἀπὸ τῆς ἁμαρτίας δουλωθέντες δὲ τῷ θεῷ, ἔχετε τὸν καρπὸν ὑμῶν εἰς ἁγιασμόν, τὸ δὲ τέλος ζωὴν αἰώνιον.
\vs{23}~τὰ γὰρ ὀψώνια τῆς ἁμαρτίας θάνατος, τὸ δὲ χάρισμα τοῦ θεοῦ ζωὴ αἰώνιος ἐν Χριστῷ Ἰησοῦ τῷ κυρίῳ ἡμῶν.

\chapheader{7}
\vs{1}~Ἢ ἀγνοεῖτε, \adeA{ἀδελφοί}, γινώσκουσιν γὰρ νόμον λαλῶ, ὅτι ὁ νόμος κυριεύει τοῦ ἀνθρώπου ἐφ’ ὅσον χρόνον ζῇ;
\vs{2}~ἡ γὰρ ὕπανδρος γυνὴ τῷ ζῶντι ἀνδρὶ δέδεται νόμῳ· ἐὰν δὲ ἀποθάνῃ ὁ ἀνήρ, κατήργηται ἀπὸ τοῦ νόμου τοῦ ἀνδρός.
\vs{3}~ἄρα οὖν ζῶντος τοῦ ἀνδρὸς μοιχαλὶς χρηματίσει ἐὰν γένηται ἀνδρὶ ἑτέρῳ· ἐὰν δὲ ἀποθάνῃ ὁ ἀνήρ, ἐλευθέρα ἐστὶν ἀπὸ τοῦ νόμου, τοῦ μὴ εἶναι αὐτὴν μοιχαλίδα γενομένην ἀνδρὶ ἑτέρῳ.
\vs{4}~Ὥστε, \adeA{ἀδελφοί} μου, καὶ ὑμεῖς ἐθανατώθητε τῷ νόμῳ διὰ τοῦ σώματος τοῦ Χριστοῦ, εἰς τὸ γενέσθαι ὑμᾶς ἑτέρῳ, τῷ ἐκ νεκρῶν ἐγερθέντι ἵνα καρποφορήσωμεν τῷ θεῷ.
\vs{5}~ὅτε γὰρ ἦμεν ἐν τῇ σαρκί, τὰ παθήματα τῶν ἁμαρτιῶν τὰ διὰ τοῦ νόμου ἐνηργεῖτο ἐν τοῖς μέλεσιν ἡμῶν εἰς τὸ καρποφορῆσαι τῷ θανάτῳ·
\vs{6}~νυνὶ δὲ κατηργήθημεν ἀπὸ τοῦ νόμου, ἀποθανόντες ἐν ᾧ κατειχόμεθα, ὥστε δουλεύειν ἡμᾶς ἐν καινότητι πνεύματος καὶ οὐ παλαιότητι γράμματος.
\vs{7}~Τί οὖν ἐροῦμεν; ὁ νόμος ἁμαρτία; μὴ γένοιτο· ἀλλὰ τὴν ἁμαρτίαν οὐκ ἔγνων εἰ μὴ διὰ νόμου, τήν τε γὰρ ἐπιθυμίαν οὐκ ᾔδειν εἰ μὴ ὁ νόμος ἔλεγεν· Οὐκ ἐπιθυμήσεις·
\vs{8}~ἀφορμὴν δὲ λαβοῦσα ἡ ἁμαρτία διὰ τῆς ἐντολῆς κατειργάσατο ἐν ἐμοὶ πᾶσαν ἐπιθυμίαν, χωρὶς γὰρ νόμου ἁμαρτία νεκρά.
\vs{9}~ἐγὼ δὲ ἔζων χωρὶς νόμου ποτέ· ἐλθούσης δὲ τῆς ἐντολῆς ἡ ἁμαρτία ἀνέζησεν,
\vs{10}~ἐγὼ δὲ ἀπέθανον, καὶ εὑρέθη μοι ἡ ἐντολὴ ἡ εἰς ζωὴν αὕτη εἰς θάνατον·
\vs{11}~ἡ γὰρ ἁμαρτία ἀφορμὴν λαβοῦσα διὰ τῆς ἐντολῆς ἐξηπάτησέν με καὶ δι’ αὐτῆς ἀπέκτεινεν.
\vs{12}~ὥστε ὁ μὲν νόμος ἅγιος, καὶ ἡ ἐντολὴ ἁγία καὶ δικαία καὶ ἀγαθή.
\vs{13}~Τὸ οὖν ἀγαθὸν ἐμοὶ ἐγένετο θάνατος; μὴ γένοιτο· ἀλλὰ ἡ ἁμαρτία, ἵνα φανῇ ἁμαρτία διὰ τοῦ ἀγαθοῦ μοι κατεργαζομένη θάνατον· ἵνα γένηται καθ’ ὑπερβολὴν ἁμαρτωλὸς ἡ ἁμαρτία διὰ τῆς ἐντολῆς.
\vs{14}~Οἴδαμεν γὰρ ὅτι ὁ νόμος πνευματικός ἐστιν· ἐγὼ δὲ σάρκινός εἰμι, πεπραμένος ὑπὸ τὴν ἁμαρτίαν.
\vs{15}~ὃ γὰρ κατεργάζομαι οὐ γινώσκω· οὐ γὰρ ὃ θέλω τοῦτο πράσσω, ἀλλ’ ὃ μισῶ τοῦτο ποιῶ.
\vs{16}~εἰ δὲ ὃ οὐ θέλω τοῦτο ποιῶ, σύμφημι τῷ νόμῳ ὅτι καλός.
\vs{17}~νυνὶ δὲ οὐκέτι ἐγὼ κατεργάζομαι αὐτὸ ἀλλὰ ἡ οἰκοῦσα ἐν ἐμοὶ ἁμαρτία.
\vs{18}~οἶδα γὰρ ὅτι οὐκ οἰκεῖ ἐν ἐμοί, τοῦτ’ ἔστιν ἐν τῇ σαρκί μου, ἀγαθόν· τὸ γὰρ θέλειν παράκειταί μοι, τὸ δὲ κατεργάζεσθαι τὸ καλὸν οὔ·
\vs{19}~οὐ γὰρ ὃ θέλω ποιῶ ἀγαθόν, ἀλλὰ ὃ οὐ θέλω κακὸν τοῦτο πράσσω.
\vs{20}~εἰ δὲ ὃ οὐ θέλω τοῦτο ποιῶ, οὐκέτι ἐγὼ κατεργάζομαι αὐτὸ ἀλλὰ ἡ οἰκοῦσα ἐν ἐμοὶ ἁμαρτία.
\vs{21}~Εὑρίσκω ἄρα τὸν νόμον τῷ θέλοντι ἐμοὶ ποιεῖν τὸ καλὸν ὅτι ἐμοὶ τὸ κακὸν παράκειται·
\vs{22}~συνήδομαι γὰρ τῷ νόμῳ τοῦ θεοῦ κατὰ τὸν ἔσω ἄνθρωπον,
\vs{23}~βλέπω δὲ ἕτερον νόμον ἐν τοῖς μέλεσίν μου ἀντιστρατευόμενον τῷ νόμῳ τοῦ νοός μου καὶ αἰχμαλωτίζοντά με ἐν τῷ νόμῳ τῆς ἁμαρτίας τῷ ὄντι ἐν τοῖς μέλεσίν μου.
\vs{24}~ταλαίπωρος ἐγὼ ἄνθρωπος· τίς με ῥύσεται ἐκ τοῦ σώματος τοῦ θανάτου τούτου;
\vs{25}~χάρις τῷ θεῷ διὰ Ἰησοῦ Χριστοῦ τοῦ κυρίου ἡμῶν. Ἄρα οὖν αὐτὸς ἐγὼ τῷ μὲν νοῒ δουλεύω νόμῳ θεοῦ, τῇ δὲ σαρκὶ νόμῳ ἁμαρτίας.

\chapheader{8}
\vs{1}~Οὐδὲν ἄρα νῦν κατάκριμα τοῖς ἐν Χριστῷ Ἰησοῦ·
\vs{2}~ὁ γὰρ νόμος τοῦ πνεύματος τῆς ζωῆς ἐν Χριστῷ Ἰησοῦ ἠλευθέρωσέν σε ἀπὸ τοῦ νόμου τῆς ἁμαρτίας καὶ τοῦ θανάτου.
\vs{3}~τὸ γὰρ ἀδύνατον τοῦ νόμου, ἐν ᾧ ἠσθένει διὰ τῆς σαρκός, ὁ θεὸς τὸν ἑαυτοῦ υἱὸν πέμψας ἐν ὁμοιώματι σαρκὸς ἁμαρτίας καὶ περὶ ἁμαρτίας κατέκρινε τὴν ἁμαρτίαν ἐν τῇ σαρκί,
\vs{4}~ἵνα τὸ δικαίωμα τοῦ νόμου πληρωθῇ ἐν ἡμῖν τοῖς μὴ κατὰ σάρκα περιπατοῦσιν ἀλλὰ κατὰ πνεῦμα·
\vs{5}~οἱ γὰρ κατὰ σάρκα ὄντες τὰ τῆς σαρκὸς φρονοῦσιν, οἱ δὲ κατὰ πνεῦμα τὰ τοῦ πνεύματος.
\vs{6}~τὸ γὰρ φρόνημα τῆς σαρκὸς θάνατος, τὸ δὲ φρόνημα τοῦ πνεύματος ζωὴ καὶ εἰρήνη·
\vs{7}~διότι τὸ φρόνημα τῆς σαρκὸς ἔχθρα εἰς θεόν, τῷ γὰρ νόμῳ τοῦ θεοῦ οὐχ ὑποτάσσεται, οὐδὲ γὰρ δύναται·
\vs{8}~οἱ δὲ ἐν σαρκὶ ὄντες θεῷ ἀρέσαι οὐ δύνανται.
\vs{9}~Ὑμεῖς δὲ οὐκ ἐστὲ ἐν σαρκὶ ἀλλὰ ἐν πνεύματι, εἴπερ πνεῦμα θεοῦ οἰκεῖ ἐν ὑμῖν. εἰ δέ τις πνεῦμα Χριστοῦ οὐκ ἔχει, οὗτος οὐκ ἔστιν αὐτοῦ.
\vs{10}~εἰ δὲ Χριστὸς ἐν ὑμῖν, τὸ μὲν σῶμα νεκρὸν διὰ ἁμαρτίαν, τὸ δὲ πνεῦμα ζωὴ διὰ δικαιοσύνην.
\vs{11}~εἰ δὲ τὸ πνεῦμα τοῦ ἐγείραντος τὸν Ἰησοῦν ἐκ νεκρῶν οἰκεῖ ἐν ὑμῖν, ὁ ἐγείρας ἐκ νεκρῶν Χριστὸν Ἰησοῦν ζῳοποιήσει καὶ τὰ θνητὰ σώματα ὑμῶν διὰ τὸ ἐνοικοῦν αὐτοῦ πνεῦμα ἐν ὑμῖν.
\vs{12}~Ἄρα οὖν, \adeA{ἀδελφοί}, ὀφειλέται ἐσμέν, οὐ τῇ σαρκὶ τοῦ κατὰ σάρκα ζῆν,
\vs{13}~εἰ γὰρ κατὰ σάρκα ζῆτε μέλλετε ἀποθνῄσκειν, εἰ δὲ πνεύματι τὰς πράξεις τοῦ σώματος θανατοῦτε, ζήσεσθε.
\vs{14}~ὅσοι γὰρ πνεύματι θεοῦ ἄγονται, οὗτοι υἱοί εἰσιν θεοῦ.
\vs{15}~οὐ γὰρ ἐλάβετε πνεῦμα δουλείας πάλιν εἰς φόβον, ἀλλὰ ἐλάβετε πνεῦμα υἱοθεσίας ἐν ᾧ κράζομεν· Αββα ὁ πατήρ·
\vs{16}~αὐτὸ τὸ πνεῦμα συμμαρτυρεῖ τῷ πνεύματι ἡμῶν ὅτι ἐσμὲν τέκνα θεοῦ.
\vs{17}~εἰ δὲ τέκνα, καὶ κληρονόμοι· κληρονόμοι μὲν θεοῦ, συγκληρονόμοι δὲ Χριστοῦ, εἴπερ συμπάσχομεν ἵνα καὶ συνδοξασθῶμεν.
\vs{18}~Λογίζομαι γὰρ ὅτι οὐκ ἄξια τὰ παθήματα τοῦ νῦν καιροῦ πρὸς τὴν μέλλουσαν δόξαν ἀποκαλυφθῆναι εἰς ἡμᾶς.
\vs{19}~ἡ γὰρ ἀποκαραδοκία τῆς κτίσεως τὴν ἀποκάλυψιν τῶν υἱῶν τοῦ θεοῦ ἀπεκδέχεται·
\vs{20}~τῇ γὰρ ματαιότητι ἡ κτίσις ὑπετάγη, οὐχ ἑκοῦσα ἀλλὰ διὰ τὸν ὑποτάξαντα, ἐφ’ ἑλπίδι
\vs{21}~ὅτι καὶ αὐτὴ ἡ κτίσις ἐλευθερωθήσεται ἀπὸ τῆς δουλείας τῆς φθορᾶς εἰς τὴν ἐλευθερίαν τῆς δόξης τῶν τέκνων τοῦ θεοῦ.
\vs{22}~οἴδαμεν γὰρ ὅτι πᾶσα ἡ κτίσις συστενάζει καὶ συνωδίνει ἄχρι τοῦ νῦν·
\vs{23}~οὐ μόνον δέ, ἀλλὰ καὶ αὐτοὶ τὴν ἀπαρχὴν τοῦ πνεύματος ἔχοντες ἡμεῖς καὶ αὐτοὶ ἐν ἑαυτοῖς στενάζομεν, υἱοθεσίαν ἀπεκδεχόμενοι τὴν ἀπολύτρωσιν τοῦ σώματος ἡμῶν.
\vs{24}~τῇ γὰρ ἐλπίδι ἐσώθημεν· ἐλπὶς δὲ βλεπομένη οὐκ ἔστιν ἐλπίς, ὃ γὰρ βλέπει τίς ἐλπίζει;
\vs{25}~εἰ δὲ ὃ οὐ βλέπομεν ἐλπίζομεν, δι’ ὑπομονῆς ἀπεκδεχόμεθα.
\vs{26}~Ὡσαύτως δὲ καὶ τὸ πνεῦμα συναντιλαμβάνεται τῇ ἀσθενείᾳ ἡμῶν· τὸ γὰρ τί προσευξώμεθα καθὸ δεῖ οὐκ οἴδαμεν, ἀλλὰ αὐτὸ τὸ πνεῦμα ὑπερεντυγχάνει στεναγμοῖς ἀλαλήτοις,
\vs{27}~ὁ δὲ ἐραυνῶν τὰς καρδίας οἶδεν τί τὸ φρόνημα τοῦ πνεύματος, ὅτι κατὰ θεὸν ἐντυγχάνει ὑπὲρ ἁγίων.
\vs{28}~Οἴδαμεν δὲ ὅτι τοῖς ἀγαπῶσι τὸν θεὸν πάντα συνεργεῖ εἰς ἀγαθόν, τοῖς κατὰ πρόθεσιν κλητοῖς οὖσιν.
\vs{29}~ὅτι οὓς προέγνω, καὶ προώρισεν συμμόρφους τῆς εἰκόνος τοῦ υἱοῦ αὐτοῦ, εἰς τὸ εἶναι αὐτὸν πρωτότοκον ἐν πολλοῖς \adeA{ἀδελφοῖς}·
\vs{30}~οὓς δὲ προώρισεν, τούτους καὶ ἐκάλεσεν· καὶ οὓς ἐκάλεσεν, τούτους καὶ ἐδικαίωσεν· οὓς δὲ ἐδικαίωσεν, τούτους καὶ ἐδόξασεν.
\vs{31}~Τί οὖν ἐροῦμεν πρὸς ταῦτα; εἰ ὁ θεὸς ὑπὲρ ἡμῶν, τίς καθ’ ἡμῶν;
\vs{32}~ὅς γε τοῦ ἰδίου υἱοῦ οὐκ ἐφείσατο, ἀλλὰ ὑπὲρ ἡμῶν πάντων παρέδωκεν αὐτόν, πῶς οὐχὶ καὶ σὺν αὐτῷ τὰ πάντα ἡμῖν χαρίσεται;
\vs{33}~τίς ἐγκαλέσει κατὰ ἐκλεκτῶν θεοῦ; θεὸς ὁ δικαιῶν·
\vs{34}~τίς ὁ κατακρινῶν; Χριστὸς ὁ ἀποθανών, μᾶλλον δὲ ἐγερθείς, ὅς καί ἐστιν ἐν δεξιᾷ τοῦ θεοῦ, ὃς καὶ ἐντυγχάνει ὑπὲρ ἡμῶν·
\vs{35}~τίς ἡμᾶς χωρίσει ἀπὸ τῆς ἀγάπης τοῦ Χριστοῦ; θλῖψις ἢ στενοχωρία ἢ διωγμὸς ἢ λιμὸς ἢ γυμνότης ἢ κίνδυνος ἢ μάχαιρα;
\vs{36}~καθὼς γέγραπται ὅτι Ἕνεκεν σοῦ θανατούμεθα ὅλην τὴν ἡμέραν, ἐλογίσθημεν ὡς πρόβατα σφαγῆς.
\vs{37}~ἀλλ’ ἐν τούτοις πᾶσιν ὑπερνικῶμεν διὰ τοῦ ἀγαπήσαντος ἡμᾶς.
\vs{38}~πέπεισμαι γὰρ ὅτι οὔτε θάνατος οὔτε ζωὴ οὔτε ἄγγελοι οὔτε ἀρχαὶ οὔτε ἐνεστῶτα οὔτε μέλλοντα οὔτε δυνάμεις
\vs{39}~οὔτε ὕψωμα οὔτε βάθος οὔτε τις κτίσις ἑτέρα δυνήσεται ἡμᾶς χωρίσαι ἀπὸ τῆς ἀγάπης τοῦ θεοῦ τῆς ἐν Χριστῷ Ἰησοῦ τῷ κυρίῳ ἡμῶν.

\chapheader{9}
\vs{1}~Ἀλήθειαν λέγω ἐν Χριστῷ, οὐ ψεύδομαι, συμμαρτυρούσης μοι τῆς συνειδήσεώς μου ἐν πνεύματι ἁγίῳ,
\vs{2}~ὅτι λύπη μοί ἐστιν μεγάλη καὶ ἀδιάλειπτος ὀδύνη τῇ καρδίᾳ μου·
\vs{3}~ηὐχόμην γὰρ ἀνάθεμα εἶναι αὐτὸς ἐγὼ ἀπὸ τοῦ Χριστοῦ ὑπὲρ τῶν \adeD{ἀδελφῶν} μου τῶν συγγενῶν μου κατὰ σάρκα,
\vs{4}~οἵτινές εἰσιν Ἰσραηλῖται, ὧν ἡ υἱοθεσία καὶ ἡ δόξα καὶ αἱ διαθῆκαι καὶ ἡ νομοθεσία καὶ ἡ λατρεία καὶ αἱ ἐπαγγελίαι,
\vs{5}~ὧν οἱ πατέρες, καὶ ἐξ ὧν ὁ χριστὸς τὸ κατὰ σάρκα, ὁ ὢν ἐπὶ πάντων, θεὸς εὐλογητὸς εἰς τοὺς αἰῶνας· ἀμήν.
\vs{6}~Οὐχ οἷον δὲ ὅτι ἐκπέπτωκεν ὁ λόγος τοῦ θεοῦ. οὐ γὰρ πάντες οἱ ἐξ Ἰσραήλ, οὗτοι Ἰσραήλ·
\vs{7}~οὐδ’ ὅτι εἰσὶν σπέρμα Ἀβραάμ, πάντες τέκνα, ἀλλ’· Ἐν Ἰσαὰκ κληθήσεταί σοι σπέρμα.
\vs{8}~τοῦτ’ ἔστιν, οὐ τὰ τέκνα τῆς σαρκὸς ταῦτα τέκνα τοῦ θεοῦ, ἀλλὰ τὰ τέκνα τῆς ἐπαγγελίας λογίζεται εἰς σπέρμα·
\vs{9}~ἐπαγγελίας γὰρ ὁ λόγος οὗτος· Κατὰ τὸν καιρὸν τοῦτον ἐλεύσομαι καὶ ἔσται τῇ Σάρρᾳ υἱός.
\vs{10}~οὐ μόνον δέ, ἀλλὰ καὶ Ῥεβέκκα ἐξ ἑνὸς κοίτην ἔχουσα, Ἰσαὰκ τοῦ πατρὸς ἡμῶν·
\vs{11}~μήπω γὰρ γεννηθέντων μηδὲ πραξάντων τι ἀγαθὸν ἢ φαῦλον, ἵνα ἡ κατ’ ἐκλογὴν πρόθεσις τοῦ θεοῦ μένῃ,
\vs{12}~οὐκ ἐξ ἔργων ἀλλ’ ἐκ τοῦ καλοῦντος, ἐρρέθη αὐτῇ ὅτι Ὁ μείζων δουλεύσει τῷ ἐλάσσονι·
\vs{13}~καθὼς γέγραπται· Τὸν Ἰακὼβ ἠγάπησα, τὸν δὲ Ἠσαῦ ἐμίσησα.
\vs{14}~Τί οὖν ἐροῦμεν; μὴ ἀδικία παρὰ τῷ θεῷ; μὴ γένοιτο·
\vs{15}~τῷ Μωϋσεῖ γὰρ λέγει· Ἐλεήσω ὃν ἂν ἐλεῶ, καὶ οἰκτιρήσω ὃν ἂν οἰκτίρω.
\vs{16}~ἄρα οὖν οὐ τοῦ θέλοντος οὐδὲ τοῦ τρέχοντος ἀλλὰ τοῦ ἐλεῶντος θεοῦ.
\vs{17}~λέγει γὰρ ἡ γραφὴ τῷ Φαραὼ ὅτι Εἰς αὐτὸ τοῦτο ἐξήγειρά σε ὅπως ἐνδείξωμαι ἐν σοὶ τὴν δύναμίν μου, καὶ ὅπως διαγγελῇ τὸ ὄνομά μου ἐν πάσῃ τῇ γῇ.
\vs{18}~ἄρα οὖν ὃν θέλει ἐλεεῖ, ὃν δὲ θέλει σκληρύνει.
\vs{19}~Ἐρεῖς μοι οὖν· Τί οὖν ἔτι μέμφεται; τῷ γὰρ βουλήματι αὐτοῦ τίς ἀνθέστηκεν;
\vs{20}~ὦ ἄνθρωπε, μενοῦνγε σὺ τίς εἶ ὁ ἀνταποκρινόμενος τῷ θεῷ; μὴ ἐρεῖ τὸ πλάσμα τῷ πλάσαντι Τί με ἐποίησας οὕτως;
\vs{21}~ἢ οὐκ ἔχει ἐξουσίαν ὁ κεραμεὺς τοῦ πηλοῦ ἐκ τοῦ αὐτοῦ φυράματος ποιῆσαι ὃ μὲν εἰς τιμὴν σκεῦος ὃ δὲ εἰς ἀτιμίαν;
\vs{22}~εἰ δὲ θέλων ὁ θεὸς ἐνδείξασθαι τὴν ὀργὴν καὶ γνωρίσαι τὸ δυνατὸν αὐτοῦ ἤνεγκεν ἐν πολλῇ μακροθυμίᾳ σκεύη ὀργῆς κατηρτισμένα εἰς ἀπώλειαν,
\vs{23}~καὶ ἵνα γνωρίσῃ τὸν πλοῦτον τῆς δόξης αὐτοῦ ἐπὶ σκεύη ἐλέους, ἃ προητοίμασεν εἰς δόξαν,
\vs{24}~οὓς καὶ ἐκάλεσεν ἡμᾶς οὐ μόνον ἐξ Ἰουδαίων ἀλλὰ καὶ ἐξ ἐθνῶν;—
\vs{25}~ὡς καὶ ἐν τῷ Ὡσηὲ λέγει· Καλέσω τὸν οὐ λαόν μου λαόν μου καὶ τὴν οὐκ ἠγαπημένην ἠγαπημένην·
\vs{26}~καὶ ἔσται ἐν τῷ τόπῳ οὗ ἐρρέθη αὐτοῖς· Οὐ λαός μου ὑμεῖς, ἐκεῖ κληθήσονται υἱοὶ θεοῦ ζῶντος.
\vs{27}~Ἠσαΐας δὲ κράζει ὑπὲρ τοῦ Ἰσραήλ· Ἐὰν ᾖ ὁ ἀριθμὸς τῶν υἱῶν Ἰσραὴλ ὡς ἡ ἄμμος τῆς θαλάσσης, τὸ ὑπόλειμμα σωθήσεται·
\vs{28}~λόγον γὰρ συντελῶν καὶ συντέμνων ποιήσει κύριος ἐπὶ τῆς γῆς.
\vs{29}~καὶ καθὼς προείρηκεν Ἠσαΐας· Εἰ μὴ κύριος Σαβαὼθ ἐγκατέλιπεν ἡμῖν σπέρμα, ὡς Σόδομα ἂν ἐγενήθημεν καὶ ὡς Γόμορρα ἂν ὡμοιώθημεν.
\vs{30}~Τί οὖν ἐροῦμεν; ὅτι ἔθνη τὰ μὴ διώκοντα δικαιοσύνην κατέλαβεν δικαιοσύνην, δικαιοσύνην δὲ τὴν ἐκ πίστεως·
\vs{31}~Ἰσραὴλ δὲ διώκων νόμον δικαιοσύνης εἰς νόμον οὐκ ἔφθασεν.
\vs{32}~διὰ τί; ὅτι οὐκ ἐκ πίστεως ἀλλ’ ὡς ἐξ ἔργων· προσέκοψαν τῷ λίθῳ τοῦ προσκόμματος,
\vs{33}~καθὼς γέγραπται· Ἰδοὺ τίθημι ἐν Σιὼν λίθον προσκόμματος καὶ πέτραν σκανδάλου, καὶ ὁ πιστεύων ἐπ’ αὐτῷ οὐ καταισχυνθήσεται.

\chapheader{10}
\vs{1}~\adeA{Ἀδελφοί}, ἡ μὲν εὐδοκία τῆς ἐμῆς καρδίας καὶ ἡ δέησις πρὸς τὸν θεὸν ὑπὲρ αὐτῶν εἰς σωτηρίαν.
\vs{2}~μαρτυρῶ γὰρ αὐτοῖς ὅτι ζῆλον θεοῦ ἔχουσιν· ἀλλ’ οὐ κατ’ ἐπίγνωσιν,
\vs{3}~ἀγνοοῦντες γὰρ τὴν τοῦ θεοῦ δικαιοσύνην, καὶ τὴν ἰδίαν ζητοῦντες στῆσαι, τῇ δικαιοσύνῃ τοῦ θεοῦ οὐχ ὑπετάγησαν·
\vs{4}~τέλος γὰρ νόμου Χριστὸς εἰς δικαιοσύνην παντὶ τῷ πιστεύοντι.
\vs{5}~Μωϋσῆς γὰρ γράφει ὅτι τὴν δικαιοσύνην τὴν ἐκ τοῦ νόμου ὁ ποιήσας ἄνθρωπος ζήσεται ἐν αὐτῇ.
\vs{6}~ἡ δὲ ἐκ πίστεως δικαιοσύνη οὕτως λέγει· Μὴ εἴπῃς ἐν τῇ καρδίᾳ σου· Τίς ἀναβήσεται εἰς τὸν οὐρανόν; τοῦτ’ ἔστιν Χριστὸν καταγαγεῖν·
\vs{7}~ἤ· Τίς καταβήσεται εἰς τὴν ἄβυσσον; τοῦτ’ ἔστιν Χριστὸν ἐκ νεκρῶν ἀναγαγεῖν.
\vs{8}~ἀλλὰ τί λέγει; Ἐγγύς σου τὸ ῥῆμά ἐστιν, ἐν τῷ στόματί σου καὶ ἐν τῇ καρδίᾳ σου, τοῦτ’ ἔστιν τὸ ῥῆμα τῆς πίστεως ὃ κηρύσσομεν.
\vs{9}~ὅτι ἐὰν ὁμολογήσῃς ἐν τῷ στόματί σου κύριον Ἰησοῦν, καὶ πιστεύσῃς ἐν τῇ καρδίᾳ σου ὅτι ὁ θεὸς αὐτὸν ἤγειρεν ἐκ νεκρῶν, σωθήσῃ·
\vs{10}~καρδίᾳ γὰρ πιστεύεται εἰς δικαιοσύνην, στόματι δὲ ὁμολογεῖται εἰς σωτηρίαν·
\vs{11}~λέγει γὰρ ἡ γραφή· Πᾶς ὁ πιστεύων ἐπ’ αὐτῷ οὐ καταισχυνθήσεται.
\vs{12}~οὐ γάρ ἐστιν διαστολὴ Ἰουδαίου τε καὶ Ἕλληνος, ὁ γὰρ αὐτὸς κύριος πάντων, πλουτῶν εἰς πάντας τοὺς ἐπικαλουμένους αὐτόν·
\vs{13}~Πᾶς γὰρ ὃς ἂν ἐπικαλέσηται τὸ ὄνομα κυρίου σωθήσεται.
\vs{14}~Πῶς οὖν ἐπικαλέσωνται εἰς ὃν οὐκ ἐπίστευσαν; πῶς δὲ πιστεύσωσιν οὗ οὐκ ἤκουσαν; πῶς δὲ ἀκούσωσιν χωρὶς κηρύσσοντος;
\vs{15}~πῶς δὲ κηρύξωσιν ἐὰν μὴ ἀποσταλῶσιν; καθὼς γέγραπται· Ὡς ὡραῖοι οἱ πόδες τῶν εὐαγγελιζομένων τὰ ἀγαθά.
\vs{16}~ἀλλ’ οὐ πάντες ὑπήκουσαν τῷ εὐαγγελίῳ· Ἠσαΐας γὰρ λέγει· Κύριε, τίς ἐπίστευσεν τῇ ἀκοῇ ἡμῶν;
\vs{17}~ἄρα ἡ πίστις ἐξ ἀκοῆς, ἡ δὲ ἀκοὴ διὰ ῥήματος Χριστοῦ.
\vs{18}~Ἀλλὰ λέγω, μὴ οὐκ ἤκουσαν; μενοῦνγε· Εἰς πᾶσαν τὴν γῆν ἐξῆλθεν ὁ φθόγγος αὐτῶν, καὶ εἰς τὰ πέρατα τῆς οἰκουμένης τὰ ῥήματα αὐτῶν.
\vs{19}~ἀλλὰ λέγω, μὴ Ἰσραὴλ οὐκ ἔγνω; πρῶτος Μωϋσῆς λέγει· Ἐγὼ παραζηλώσω ὑμᾶς ἐπ’ οὐκ ἔθνει, ἐπ’ ἔθνει ἀσυνέτῳ παροργιῶ ὑμᾶς.
\vs{20}~Ἠσαΐας δὲ ἀποτολμᾷ καὶ λέγει· Εὑρέθην ἐν τοῖς ἐμὲ μὴ ζητοῦσιν, ἐμφανὴς ἐγενόμην τοῖς ἐμὲ μὴ ἐπερωτῶσιν.
\vs{21}~πρὸς δὲ τὸν Ἰσραὴλ λέγει· Ὅλην τὴν ἡμέραν ἐξεπέτασα τὰς χεῖράς μου πρὸς λαὸν ἀπειθοῦντα καὶ ἀντιλέγοντα.

\chapheader{11}
\vs{1}~Λέγω οὖν, μὴ ἀπώσατο ὁ θεὸς τὸν λαὸν αὐτοῦ; μὴ γένοιτο· καὶ γὰρ ἐγὼ Ἰσραηλίτης εἰμί, ἐκ σπέρματος Ἀβραάμ, φυλῆς Βενιαμίν.
\vs{2}~οὐκ ἀπώσατο ὁ θεὸς τὸν λαὸν αὐτοῦ ὃν προέγνω. ἢ οὐκ οἴδατε ἐν Ἠλίᾳ τί λέγει ἡ γραφή, ὡς ἐντυγχάνει τῷ θεῷ κατὰ τοῦ Ἰσραήλ;
\vs{3}~Κύριε, τοὺς προφήτας σου ἀπέκτειναν, τὰ θυσιαστήριά σου κατέσκαψαν, κἀγὼ ὑπελείφθην μόνος, καὶ ζητοῦσιν τὴν ψυχήν μου.
\vs{4}~ἀλλὰ τί λέγει αὐτῷ ὁ χρηματισμός; Κατέλιπον ἐμαυτῷ ἑπτακισχιλίους ἄνδρας, οἵτινες οὐκ ἔκαμψαν γόνυ τῇ Βάαλ.
\vs{5}~οὕτως οὖν καὶ ἐν τῷ νῦν καιρῷ λεῖμμα κατ’ ἐκλογὴν χάριτος γέγονεν·
\vs{6}~εἰ δὲ χάριτι, οὐκέτι ἐξ ἔργων, ἐπεὶ ἡ χάρις οὐκέτι γίνεται χάρις.
\vs{7}~τί οὖν; ὃ ἐπιζητεῖ Ἰσραήλ, τοῦτο οὐκ ἐπέτυχεν, ἡ δὲ ἐκλογὴ ἐπέτυχεν· οἱ δὲ λοιποὶ ἐπωρώθησαν,
\vs{8}~καθὼς γέγραπται· Ἔδωκεν αὐτοῖς ὁ θεὸς πνεῦμα κατανύξεως, ὀφθαλμοὺς τοῦ μὴ βλέπειν καὶ ὦτα τοῦ μὴ ἀκούειν, ἕως τῆς σήμερον ἡμέρας.
\vs{9}~καὶ Δαυὶδ λέγει· Γενηθήτω ἡ τράπεζα αὐτῶν εἰς παγίδα καὶ εἰς θήραν καὶ εἰς σκάνδαλον καὶ εἰς ἀνταπόδομα αὐτοῖς,
\vs{10}~σκοτισθήτωσαν οἱ ὀφθαλμοὶ αὐτῶν τοῦ μὴ βλέπειν, καὶ τὸν νῶτον αὐτῶν διὰ παντὸς σύγκαμψον.
\vs{11}~Λέγω οὖν, μὴ ἔπταισαν ἵνα πέσωσιν; μὴ γένοιτο· ἀλλὰ τῷ αὐτῶν παραπτώματι ἡ σωτηρία τοῖς ἔθνεσιν, εἰς τὸ παραζηλῶσαι αὐτούς.
\vs{12}~εἰ δὲ τὸ παράπτωμα αὐτῶν πλοῦτος κόσμου καὶ τὸ ἥττημα αὐτῶν πλοῦτος ἐθνῶν, πόσῳ μᾶλλον τὸ πλήρωμα αὐτῶν.
\vs{13}~Ὑμῖν δὲ λέγω τοῖς ἔθνεσιν. ἐφ’ ὅσον μὲν οὖν εἰμι ἐγὼ ἐθνῶν ἀπόστολος, τὴν διακονίαν μου δοξάζω,
\vs{14}~εἴ πως παραζηλώσω μου τὴν σάρκα καὶ σώσω τινὰς ἐξ αὐτῶν.
\vs{15}~εἰ γὰρ ἡ ἀποβολὴ αὐτῶν καταλλαγὴ κόσμου, τίς ἡ πρόσλημψις εἰ μὴ ζωὴ ἐκ νεκρῶν;
\vs{16}~εἰ δὲ ἡ ἀπαρχὴ ἁγία, καὶ τὸ φύραμα· καὶ εἰ ἡ ῥίζα ἁγία, καὶ οἱ κλάδοι.
\vs{17}~Εἰ δέ τινες τῶν κλάδων ἐξεκλάσθησαν, σὺ δὲ ἀγριέλαιος ὢν ἐνεκεντρίσθης ἐν αὐτοῖς καὶ συγκοινωνὸς τῆς ῥίζης τῆς πιότητος τῆς ἐλαίας ἐγένου,
\vs{18}~μὴ κατακαυχῶ τῶν κλάδων· εἰ δὲ κατακαυχᾶσαι, οὐ σὺ τὴν ῥίζαν βαστάζεις ἀλλὰ ἡ ῥίζα σέ.
\vs{19}~ἐρεῖς οὖν· Ἐξεκλάσθησαν κλάδοι ἵνα ἐγὼ ἐγκεντρισθῶ.
\vs{20}~καλῶς· τῇ ἀπιστίᾳ ἐξεκλάσθησαν, σὺ δὲ τῇ πίστει ἕστηκας. μὴ ὑψηλὰ φρόνει, ἀλλὰ φοβοῦ·
\vs{21}~εἰ γὰρ ὁ θεὸς τῶν κατὰ φύσιν κλάδων οὐκ ἐφείσατο, οὐδὲ σοῦ φείσεται.
\vs{22}~ἴδε οὖν χρηστότητα καὶ ἀποτομίαν θεοῦ· ἐπὶ μὲν τοὺς πεσόντας ἀποτομία, ἐπὶ δὲ σὲ χρηστότης θεοῦ, ἐὰν ἐπιμένῃς τῇ χρηστότητι, ἐπεὶ καὶ σὺ ἐκκοπήσῃ.
\vs{23}~κἀκεῖνοι δέ, ἐὰν μὴ ἐπιμένωσι τῇ ἀπιστίᾳ, ἐγκεντρισθήσονται· δυνατὸς γάρ ἐστιν ὁ θεὸς πάλιν ἐγκεντρίσαι αὐτούς.
\vs{24}~εἰ γὰρ σὺ ἐκ τῆς κατὰ φύσιν ἐξεκόπης ἀγριελαίου καὶ παρὰ φύσιν ἐνεκεντρίσθης εἰς καλλιέλαιον, πόσῳ μᾶλλον οὗτοι οἱ κατὰ φύσιν ἐγκεντρισθήσονται τῇ ἰδίᾳ ἐλαίᾳ.
\vs{25}~Οὐ γὰρ θέλω ὑμᾶς ἀγνοεῖν, \adeA{ἀδελφοί}, τὸ μυστήριον τοῦτο, ἵνα μὴ ἦτε ἑαυτοῖς φρόνιμοι, ὅτι πώρωσις ἀπὸ μέρους τῷ Ἰσραὴλ γέγονεν ἄχρι οὗ τὸ πλήρωμα τῶν ἐθνῶν εἰσέλθῃ,
\vs{26}~καὶ οὕτως πᾶς Ἰσραὴλ σωθήσεται· καθὼς γέγραπται· Ἥξει ἐκ Σιὼν ὁ ῥυόμενος, ἀποστρέψει ἀσεβείας ἀπὸ Ἰακώβ.
\vs{27}~καὶ αὕτη αὐτοῖς ἡ παρ’ ἐμοῦ διαθήκη, ὅταν ἀφέλωμαι τὰς ἁμαρτίας αὐτῶν.
\vs{28}~κατὰ μὲν τὸ εὐαγγέλιον ἐχθροὶ δι’ ὑμᾶς, κατὰ δὲ τὴν ἐκλογὴν ἀγαπητοὶ διὰ τοὺς πατέρας·
\vs{29}~ἀμεταμέλητα γὰρ τὰ χαρίσματα καὶ ἡ κλῆσις τοῦ θεοῦ.
\vs{30}~ὥσπερ γὰρ ὑμεῖς ποτε ἠπειθήσατε τῷ θεῷ, νῦν δὲ ἠλεήθητε τῇ τούτων ἀπειθείᾳ,
\vs{31}~οὕτως καὶ οὗτοι νῦν ἠπείθησαν τῷ ὑμετέρῳ ἐλέει ἵνα καὶ αὐτοὶ νῦν ἐλεηθῶσιν·
\vs{32}~συνέκλεισεν γὰρ ὁ θεὸς τοὺς πάντας εἰς ἀπείθειαν ἵνα τοὺς πάντας ἐλεήσῃ.
\vs{33}~Ὦ βάθος πλούτου καὶ σοφίας καὶ γνώσεως θεοῦ· ὡς ἀνεξεραύνητα τὰ κρίματα αὐτοῦ καὶ ἀνεξιχνίαστοι αἱ ὁδοὶ αὐτοῦ.
\vs{34}~Τίς γὰρ ἔγνω νοῦν κυρίου; ἢ τίς σύμβουλος αὐτοῦ ἐγένετο;
\vs{35}~ἢ τίς προέδωκεν αὐτῷ, καὶ ἀνταποδοθήσεται αὐτῷ;
\vs{36}~ὅτι ἐξ αὐτοῦ καὶ δι’ αὐτοῦ καὶ εἰς αὐτὸν τὰ πάντα· αὐτῷ ἡ δόξα εἰς τοὺς αἰῶνας, ἀμήν.

\chapheader{12}
\vs{1}~Παρακαλῶ οὖν ὑμᾶς, \adeA{ἀδελφοί}, διὰ τῶν οἰκτιρμῶν τοῦ θεοῦ παραστῆσαι τὰ σώματα ὑμῶν θυσίαν ζῶσαν ἁγίαν εὐάρεστον τῷ θεῷ, τὴν λογικὴν λατρείαν ὑμῶν·
\vs{2}~καὶ μὴ συσχηματίζεσθε τῷ αἰῶνι τούτῳ, ἀλλὰ μεταμορφοῦσθε τῇ ἀνακαινώσει τοῦ νοός, εἰς τὸ δοκιμάζειν ὑμᾶς τί τὸ θέλημα τοῦ θεοῦ, τὸ ἀγαθὸν καὶ εὐάρεστον καὶ τέλειον.
\vs{3}~Λέγω γὰρ διὰ τῆς χάριτος τῆς δοθείσης μοι παντὶ τῷ ὄντι ἐν ὑμῖν μὴ ὑπερφρονεῖν παρ’ ὃ δεῖ φρονεῖν, ἀλλὰ φρονεῖν εἰς τὸ σωφρονεῖν, ἑκάστῳ ὡς ὁ θεὸς ἐμέρισεν μέτρον πίστεως.
\vs{4}~καθάπερ γὰρ ἐν ἑνὶ σώματι πολλὰ μέλη ἔχομεν, τὰ δὲ μέλη πάντα οὐ τὴν αὐτὴν ἔχει πρᾶξιν,
\vs{5}~οὕτως οἱ πολλοὶ ἓν σῶμά ἐσμεν ἐν Χριστῷ, τὸ δὲ καθ’ εἷς ἀλλήλων μέλη.
\vs{6}~ἔχοντες δὲ χαρίσματα κατὰ τὴν χάριν τὴν δοθεῖσαν ἡμῖν διάφορα, εἴτε προφητείαν κατὰ τὴν ἀναλογίαν τῆς πίστεως,
\vs{7}~εἴτε διακονίαν ἐν τῇ διακονίᾳ, εἴτε ὁ διδάσκων ἐν τῇ διδασκαλίᾳ,
\vs{8}~εἴτε ὁ παρακαλῶν ἐν τῇ παρακλήσει, ὁ μεταδιδοὺς ἐν ἁπλότητι, ὁ προϊστάμενος ἐν σπουδῇ, ὁ ἐλεῶν ἐν ἱλαρότητι.
\vs{9}~Ἡ ἀγάπη ἀνυπόκριτος. ἀποστυγοῦντες τὸ πονηρόν, κολλώμενοι τῷ ἀγαθῷ·
\vs{10}~τῇ \adeA{φιλαδελφίᾳ} εἰς ἀλλήλους φιλόστοργοι, τῇ τιμῇ ἀλλήλους προηγούμενοι,
\vs{11}~τῇ σπουδῇ μὴ ὀκνηροί, τῷ πνεύματι ζέοντες, τῷ κυρίῳ δουλεύοντες,
\vs{12}~τῇ ἐλπίδι χαίροντες, τῇ θλίψει ὑπομένοντες, τῇ προσευχῇ προσκαρτεροῦντες,
\vs{13}~ταῖς χρείαις τῶν ἁγίων κοινωνοῦντες, τὴν φιλοξενίαν διώκοντες.
\vs{14}~εὐλογεῖτε τοὺς διώκοντας, εὐλογεῖτε καὶ μὴ καταρᾶσθε.
\vs{15}~χαίρειν μετὰ χαιρόντων, κλαίειν μετὰ κλαιόντων.
\vs{16}~τὸ αὐτὸ εἰς ἀλλήλους φρονοῦντες, μὴ τὰ ὑψηλὰ φρονοῦντες ἀλλὰ τοῖς ταπεινοῖς συναπαγόμενοι. μὴ γίνεσθε φρόνιμοι παρ’ ἑαυτοῖς.
\vs{17}~μηδενὶ κακὸν ἀντὶ κακοῦ ἀποδιδόντες· προνοούμενοι καλὰ ἐνώπιον πάντων ἀνθρώπων·
\vs{18}~εἰ δυνατόν, τὸ ἐξ ὑμῶν μετὰ πάντων ἀνθρώπων εἰρηνεύοντες·
\vs{19}~μὴ ἑαυτοὺς ἐκδικοῦντες, ἀγαπητοί, ἀλλὰ δότε τόπον τῇ ὀργῇ, γέγραπται γάρ· Ἐμοὶ ἐκδίκησις, ἐγὼ ἀνταποδώσω, λέγει κύριος.
\vs{20}~ἀλλὰ ἐὰν πεινᾷ ὁ ἐχθρός σου, ψώμιζε αὐτόν· ἐὰν διψᾷ, πότιζε αὐτόν· τοῦτο γὰρ ποιῶν ἄνθρακας πυρὸς σωρεύσεις ἐπὶ τὴν κεφαλὴν αὐτοῦ.
\vs{21}~μὴ νικῶ ὑπὸ τοῦ κακοῦ, ἀλλὰ νίκα ἐν τῷ ἀγαθῷ τὸ κακόν.

\chapheader{13}
\vs{1}~Πᾶσα ψυχὴ ἐξουσίαις ὑπερεχούσαις ὑποτασσέσθω, οὐ γὰρ ἔστιν ἐξουσία εἰ μὴ ὑπὸ θεοῦ, αἱ δὲ οὖσαι ὑπὸ θεοῦ τεταγμέναι εἰσίν.
\vs{2}~ὥστε ὁ ἀντιτασσόμενος τῇ ἐξουσίᾳ τῇ τοῦ θεοῦ διαταγῇ ἀνθέστηκεν, οἱ δὲ ἀνθεστηκότες ἑαυτοῖς κρίμα λήμψονται.
\vs{3}~οἱ γὰρ ἄρχοντες οὐκ εἰσὶν φόβος τῷ ἀγαθῷ ἔργῳ ἀλλὰ τῷ κακῷ. θέλεις δὲ μὴ φοβεῖσθαι τὴν ἐξουσίαν; τὸ ἀγαθὸν ποίει, καὶ ἕξεις ἔπαινον ἐξ αὐτῆς·
\vs{4}~θεοῦ γὰρ διάκονός ἐστιν σοὶ εἰς τὸ ἀγαθόν. ἐὰν δὲ τὸ κακὸν ποιῇς, φοβοῦ· οὐ γὰρ εἰκῇ τὴν μάχαιραν φορεῖ· θεοῦ γὰρ διάκονός ἐστιν, ἔκδικος εἰς ὀργὴν τῷ τὸ κακὸν πράσσοντι.
\vs{5}~διὸ ἀνάγκη ὑποτάσσεσθαι, οὐ μόνον διὰ τὴν ὀργὴν ἀλλὰ καὶ διὰ τὴν συνείδησιν,
\vs{6}~διὰ τοῦτο γὰρ καὶ φόρους τελεῖτε, λειτουργοὶ γὰρ θεοῦ εἰσιν εἰς αὐτὸ τοῦτο προσκαρτεροῦντες.
\vs{7}~ἀπόδοτε πᾶσι τὰς ὀφειλάς, τῷ τὸν φόρον τὸν φόρον, τῷ τὸ τέλος τὸ τέλος, τῷ τὸν φόβον τὸν φόβον, τῷ τὴν τιμὴν τὴν τιμήν.
\vs{8}~Μηδενὶ μηδὲν ὀφείλετε, εἰ μὴ τὸ ἀλλήλους ἀγαπᾶν· ὁ γὰρ ἀγαπῶν τὸν ἕτερον νόμον πεπλήρωκεν.
\vs{9}~τὸ γάρ· Οὐ μοιχεύσεις, Οὐ φονεύσεις, Οὐ κλέψεις, Οὐκ ἐπιθυμήσεις, καὶ εἴ τις ἑτέρα ἐντολή, ἐν τῷ λόγῳ τούτῳ ἀνακεφαλαιοῦται, ἐν τῷ· Ἀγαπήσεις τὸν πλησίον σου ὡς σεαυτόν.
\vs{10}~ἡ ἀγάπη τῷ πλησίον κακὸν οὐκ ἐργάζεται· πλήρωμα οὖν νόμου ἡ ἀγάπη.
\vs{11}~Καὶ τοῦτο εἰδότες τὸν καιρόν, ὅτι ὥρα ἤδη ὑμᾶς ἐξ ὕπνου ἐγερθῆναι, νῦν γὰρ ἐγγύτερον ἡμῶν ἡ σωτηρία ἢ ὅτε ἐπιστεύσαμεν.
\vs{12}~ἡ νὺξ προέκοψεν, ἡ δὲ ἡμέρα ἤγγικεν. ἀποβαλώμεθα οὖν τὰ ἔργα τοῦ σκότους, ἐνδυσώμεθα δὲ τὰ ὅπλα τοῦ φωτός.
\vs{13}~ὡς ἐν ἡμέρᾳ εὐσχημόνως περιπατήσωμεν, μὴ κώμοις καὶ μέθαις, μὴ κοίταις καὶ ἀσελγείαις, μὴ ἔριδι καὶ ζήλῳ,
\vs{14}~ἀλλὰ ἐνδύσασθε τὸν κύριον Ἰησοῦν Χριστόν, καὶ τῆς σαρκὸς πρόνοιαν μὴ ποιεῖσθε εἰς ἐπιθυμίας.

\chapheader{14}
\vs{1}~Τὸν δὲ ἀσθενοῦντα τῇ πίστει προσλαμβάνεσθε, μὴ εἰς διακρίσεις διαλογισμῶν.
\vs{2}~ὃς μὲν πιστεύει φαγεῖν πάντα, ὁ δὲ ἀσθενῶν λάχανα ἐσθίει.
\vs{3}~ὁ ἐσθίων τὸν μὴ ἐσθίοντα μὴ ἐξουθενείτω, ὁ δὲ μὴ ἐσθίων τὸν ἐσθίοντα μὴ κρινέτω, ὁ θεὸς γὰρ αὐτὸν προσελάβετο.
\vs{4}~σὺ τίς εἶ ὁ κρίνων ἀλλότριον οἰκέτην; τῷ ἰδίῳ κυρίῳ στήκει ἢ πίπτει· σταθήσεται δέ, δυνατεῖ γὰρ ὁ κύριος στῆσαι αὐτόν.
\vs{5}~Ὃς μὲν κρίνει ἡμέραν παρ’ ἡμέραν, ὃς δὲ κρίνει πᾶσαν ἡμέραν· ἕκαστος ἐν τῷ ἰδίῳ νοῒ πληροφορείσθω·
\vs{6}~ὁ φρονῶν τὴν ἡμέραν κυρίῳ φρονεῖ. καὶ ὁ ἐσθίων κυρίῳ ἐσθίει, εὐχαριστεῖ γὰρ τῷ θεῷ· καὶ ὁ μὴ ἐσθίων κυρίῳ οὐκ ἐσθίει, καὶ εὐχαριστεῖ τῷ θεῷ.
\vs{7}~Οὐδεὶς γὰρ ἡμῶν ἑαυτῷ ζῇ, καὶ οὐδεὶς ἑαυτῷ ἀποθνῄσκει·
\vs{8}~ἐάν τε γὰρ ζῶμεν, τῷ κυρίῳ ζῶμεν, ἐάν τε ἀποθνῄσκωμεν, τῷ κυρίῳ ἀποθνῄσκομεν. ἐάν τε οὖν ζῶμεν ἐάν τε ἀποθνῄσκωμεν, τοῦ κυρίου ἐσμέν.
\vs{9}~εἰς τοῦτο γὰρ Χριστὸς ἀπέθανεν καὶ ἔζησεν ἵνα καὶ νεκρῶν καὶ ζώντων κυριεύσῃ.
\vs{10}~Σὺ δὲ τί κρίνεις τὸν \adeA{ἀδελφόν} σου; ἢ καὶ σὺ τί ἐξουθενεῖς τὸν \adeA{ἀδελφόν} σου; πάντες γὰρ παραστησόμεθα τῷ βήματι τοῦ θεοῦ,
\vs{11}~γέγραπται γάρ· Ζῶ ἐγώ, λέγει κύριος, ὅτι ἐμοὶ κάμψει πᾶν γόνυ, καὶ πᾶσα γλῶσσα ἐξομολογήσεται τῷ θεῷ.
\vs{12}~ἄρα ἕκαστος ἡμῶν περὶ ἑαυτοῦ λόγον δώσει.
\vs{13}~Μηκέτι οὖν ἀλλήλους κρίνωμεν· ἀλλὰ τοῦτο κρίνατε μᾶλλον, τὸ μὴ τιθέναι πρόσκομμα τῷ \adeA{ἀδελφῷ} ἢ σκάνδαλον.
\vs{14}~οἶδα καὶ πέπεισμαι ἐν κυρίῳ Ἰησοῦ ὅτι οὐδὲν κοινὸν δι’ ἑαυτοῦ· εἰ μὴ τῷ λογιζομένῳ τι κοινὸν εἶναι, ἐκείνῳ κοινόν.
\vs{15}~εἰ γὰρ διὰ βρῶμα ὁ \adeA{ἀδελφός} σου λυπεῖται, οὐκέτι κατὰ ἀγάπην περιπατεῖς. μὴ τῷ βρώματί σου ἐκεῖνον ἀπόλλυε ὑπὲρ οὗ Χριστὸς ἀπέθανεν.
\vs{16}~μὴ βλασφημείσθω οὖν ὑμῶν τὸ ἀγαθόν.
\vs{17}~οὐ γάρ ἐστιν ἡ βασιλεία τοῦ θεοῦ βρῶσις καὶ πόσις, ἀλλὰ δικαιοσύνη καὶ εἰρήνη καὶ χαρὰ ἐν πνεύματι ἁγίῳ·
\vs{18}~ὁ γὰρ ἐν τούτῳ δουλεύων τῷ Χριστῷ εὐάρεστος τῷ θεῷ καὶ δόκιμος τοῖς ἀνθρώποις.
\vs{19}~ἄρα οὖν τὰ τῆς εἰρήνης διώκωμεν καὶ τὰ τῆς οἰκοδομῆς τῆς εἰς ἀλλήλους.
\vs{20}~μὴ ἕνεκεν βρώματος κατάλυε τὸ ἔργον τοῦ θεοῦ. πάντα μὲν καθαρά, ἀλλὰ κακὸν τῷ ἀνθρώπῳ τῷ διὰ προσκόμματος ἐσθίοντι.
\vs{21}~καλὸν τὸ μὴ φαγεῖν κρέα μηδὲ πιεῖν οἶνον μηδὲ ἐν ᾧ ὁ \adeA{ἀδελφός} σου προσκόπτει ἢ σκανδαλίζεται ἢ ἀσθενεῖ·
\vs{22}~σὺ πίστιν ἣν ἔχεις κατὰ σεαυτὸν ἔχε ἐνώπιον τοῦ θεοῦ. μακάριος ὁ μὴ κρίνων ἑαυτὸν ἐν ᾧ δοκιμάζει·
\vs{23}~ὁ δὲ διακρινόμενος ἐὰν φάγῃ κατακέκριται, ὅτι οὐκ ἐκ πίστεως· πᾶν δὲ ὃ οὐκ ἐκ πίστεως ἁμαρτία ἐστίν.

\chapheader{15}
\vs{1}~Ὀφείλομεν δὲ ἡμεῖς οἱ δυνατοὶ τὰ ἀσθενήματα τῶν ἀδυνάτων βαστάζειν, καὶ μὴ ἑαυτοῖς ἀρέσκειν.
\vs{2}~ἕκαστος ἡμῶν τῷ πλησίον ἀρεσκέτω εἰς τὸ ἀγαθὸν πρὸς οἰκοδομήν·
\vs{3}~καὶ γὰρ ὁ Χριστὸς οὐχ ἑαυτῷ ἤρεσεν· ἀλλὰ καθὼς γέγραπται· Οἱ ὀνειδισμοὶ τῶν ὀνειδιζόντων σε ἐπέπεσαν ἐπ’ ἐμέ.
\vs{4}~ὅσα γὰρ προεγράφη, εἰς τὴν ἡμετέραν διδασκαλίαν ἐγράφη, ἵνα διὰ τῆς ὑπομονῆς καὶ διὰ τῆς παρακλήσεως τῶν γραφῶν τὴν ἐλπίδα ἔχωμεν.
\vs{5}~ὁ δὲ θεὸς τῆς ὑπομονῆς καὶ τῆς παρακλήσεως δῴη ὑμῖν τὸ αὐτὸ φρονεῖν ἐν ἀλλήλοις κατὰ Χριστὸν Ἰησοῦν,
\vs{6}~ἵνα ὁμοθυμαδὸν ἐν ἑνὶ στόματι δοξάζητε τὸν θεὸν καὶ πατέρα τοῦ κυρίου ἡμῶν Ἰησοῦ Χριστοῦ.
\vs{7}~Διὸ προσλαμβάνεσθε ἀλλήλους, καθὼς καὶ ὁ Χριστὸς προσελάβετο ὑμᾶς, εἰς δόξαν τοῦ θεοῦ.
\vs{8}~λέγω γὰρ Χριστὸν διάκονον γεγενῆσθαι περιτομῆς ὑπὲρ ἀληθείας θεοῦ, εἰς τὸ βεβαιῶσαι τὰς ἐπαγγελίας τῶν πατέρων,
\vs{9}~τὰ δὲ ἔθνη ὑπὲρ ἐλέους δοξάσαι τὸν θεόν· καθὼς γέγραπται· Διὰ τοῦτο ἐξομολογήσομαί σοι ἐν ἔθνεσι, καὶ τῷ ὀνόματί σου ψαλῶ.
\vs{10}~καὶ πάλιν λέγει· Εὐφράνθητε, ἔθνη, μετὰ τοῦ λαοῦ αὐτοῦ.
\vs{11}~καὶ πάλιν· Αἰνεῖτε, πάντα τὰ ἔθνη, τὸν κύριον, καὶ ἐπαινεσάτωσαν αὐτὸν πάντες οἱ λαοί.
\vs{12}~καὶ πάλιν Ἠσαΐας λέγει· Ἔσται ἡ ῥίζα τοῦ Ἰεσσαί, καὶ ὁ ἀνιστάμενος ἄρχειν ἐθνῶν· ἐπ’ αὐτῷ ἔθνη ἐλπιοῦσιν.
\vs{13}~ὁ δὲ θεὸς τῆς ἐλπίδος πληρώσαι ὑμᾶς πάσης χαρᾶς καὶ εἰρήνης ἐν τῷ πιστεύειν, εἰς τὸ περισσεύειν ὑμᾶς ἐν τῇ ἐλπίδι ἐν δυνάμει πνεύματος ἁγίου.
\vs{14}~Πέπεισμαι δέ, \adeA{ἀδελφοί} μου, καὶ αὐτὸς ἐγὼ περὶ ὑμῶν, ὅτι καὶ αὐτοὶ μεστοί ἐστε ἀγαθωσύνης, πεπληρωμένοι πάσης γνώσεως, δυνάμενοι καὶ ἀλλήλους νουθετεῖν.
\vs{15}~τολμηρότερον δὲ ἔγραψα ὑμῖν ἀπὸ μέρους, ὡς ἐπαναμιμνῄσκων ὑμᾶς, διὰ τὴν χάριν τὴν δοθεῖσάν μοι ὑπὸ τοῦ θεοῦ
\vs{16}~εἰς τὸ εἶναί με λειτουργὸν Χριστοῦ Ἰησοῦ εἰς τὰ ἔθνη, ἱερουργοῦντα τὸ εὐαγγέλιον τοῦ θεοῦ, ἵνα γένηται ἡ προσφορὰ τῶν ἐθνῶν εὐπρόσδεκτος, ἡγιασμένη ἐν πνεύματι ἁγίῳ.
\vs{17}~ἔχω οὖν τὴν καύχησιν ἐν Χριστῷ Ἰησοῦ τὰ πρὸς τὸν θεόν·
\vs{18}~οὐ γὰρ τολμήσω τι λαλεῖν ὧν οὐ κατειργάσατο Χριστὸς δι’ ἐμοῦ εἰς ὑπακοὴν ἐθνῶν, λόγῳ καὶ ἔργῳ,
\vs{19}~ἐν δυνάμει σημείων καὶ τεράτων, ἐν δυνάμει πνεύματος· ὥστε με ἀπὸ Ἰερουσαλὴμ καὶ κύκλῳ μέχρι τοῦ Ἰλλυρικοῦ πεπληρωκέναι τὸ εὐαγγέλιον τοῦ Χριστοῦ,
\vs{20}~οὕτως δὲ φιλοτιμούμενον εὐαγγελίζεσθαι οὐχ ὅπου ὠνομάσθη Χριστός, ἵνα μὴ ἐπ’ ἀλλότριον θεμέλιον οἰκοδομῶ,
\vs{21}~ἀλλὰ καθὼς γέγραπται· Οἷς οὐκ ἀνηγγέλη περὶ αὐτοῦ ὄψονται, καὶ οἳ οὐκ ἀκηκόασιν συνήσουσιν.
\vs{22}~Διὸ καὶ ἐνεκοπτόμην τὰ πολλὰ τοῦ ἐλθεῖν πρὸς ὑμᾶς·
\vs{23}~νυνὶ δὲ μηκέτι τόπον ἔχων ἐν τοῖς κλίμασι τούτοις, ἐπιποθίαν δὲ ἔχων τοῦ ἐλθεῖν πρὸς ὑμᾶς ἀπὸ ἱκανῶν ἐτῶν,
\vs{24}~ὡς ἂν πορεύωμαι εἰς τὴν Σπανίαν, ἐλπίζω γὰρ διαπορευόμενος θεάσασθαι ὑμᾶς καὶ ὑφ’ ὑμῶν προπεμφθῆναι ἐκεῖ ἐὰν ὑμῶν πρῶτον ἀπὸ μέρους ἐμπλησθῶ—
\vs{25}~νυνὶ δὲ πορεύομαι εἰς Ἰερουσαλὴμ διακονῶν τοῖς ἁγίοις.
\vs{26}~εὐδόκησαν γὰρ Μακεδονία καὶ Ἀχαΐα κοινωνίαν τινὰ ποιήσασθαι εἰς τοὺς πτωχοὺς τῶν ἁγίων τῶν ἐν Ἰερουσαλήμ.
\vs{27}~εὐδόκησαν γάρ, καὶ ὀφειλέται εἰσὶν αὐτῶν· εἰ γὰρ τοῖς πνευματικοῖς αὐτῶν ἐκοινώνησαν τὰ ἔθνη, ὀφείλουσιν καὶ ἐν τοῖς σαρκικοῖς λειτουργῆσαι αὐτοῖς.
\vs{28}~τοῦτο οὖν ἐπιτελέσας, καὶ σφραγισάμενος αὐτοῖς τὸν καρπὸν τοῦτον, ἀπελεύσομαι δι’ ὑμῶν εἰς Σπανίαν·
\vs{29}~οἶδα δὲ ὅτι ἐρχόμενος πρὸς ὑμᾶς ἐν πληρώματι εὐλογίας Χριστοῦ ἐλεύσομαι.
\vs{30}~Παρακαλῶ δὲ ὑμᾶς, \adeA{ἀδελφοί}, διὰ τοῦ κυρίου ἡμῶν Ἰησοῦ Χριστοῦ καὶ διὰ τῆς ἀγάπης τοῦ πνεύματος συναγωνίσασθαί μοι ἐν ταῖς προσευχαῖς ὑπὲρ ἐμοῦ πρὸς τὸν θεόν,
\vs{31}~ἵνα ῥυσθῶ ἀπὸ τῶν ἀπειθούντων ἐν τῇ Ἰουδαίᾳ καὶ ἡ διακονία μου ἡ εἰς Ἰερουσαλὴμ εὐπρόσδεκτος τοῖς ἁγίοις γένηται,
\vs{32}~ἵνα ἐν χαρᾷ ἐλθὼν πρὸς ὑμᾶς διὰ θελήματος θεοῦ συναναπαύσωμαι ὑμῖν.
\vs{33}~ὁ δὲ θεὸς τῆς εἰρήνης μετὰ πάντων ὑμῶν· ἀμήν.

\chapheader{16}
\vs{1}~Συνίστημι δὲ ὑμῖν Φοίβην τὴν \adeA{ἀδελφὴν} ἡμῶν, οὖσαν καὶ διάκονον τῆς ἐκκλησίας τῆς ἐν Κεγχρεαῖς,
\vs{2}~ἵνα αὐτὴν προσδέξησθε ἐν κυρίῳ ἀξίως τῶν ἁγίων, καὶ παραστῆτε αὐτῇ ἐν ᾧ ἂν ὑμῶν χρῄζῃ πράγματι, καὶ γὰρ αὐτὴ προστάτις πολλῶν ἐγενήθη καὶ ἐμοῦ αὐτοῦ.
\vs{3}~Ἀσπάσασθε Πρίσκαν καὶ Ἀκύλαν τοὺς συνεργούς μου ἐν Χριστῷ Ἰησοῦ,
\vs{4}~οἵτινες ὑπὲρ τῆς ψυχῆς μου τὸν ἑαυτῶν τράχηλον ὑπέθηκαν, οἷς οὐκ ἐγὼ μόνος εὐχαριστῶ ἀλλὰ καὶ πᾶσαι αἱ ἐκκλησίαι τῶν ἐθνῶν,
\vs{5}~καὶ τὴν κατ’ οἶκον αὐτῶν ἐκκλησίαν. ἀσπάσασθε Ἐπαίνετον τὸν ἀγαπητόν μου, ὅς ἐστιν ἀπαρχὴ τῆς Ἀσίας εἰς Χριστόν.
\vs{6}~ἀσπάσασθε Μαριάμ, ἥτις πολλὰ ἐκοπίασεν εἰς ὑμᾶς.
\vs{7}~ἀσπάσασθε Ἀνδρόνικον καὶ Ἰουνίαν τοὺς συγγενεῖς μου καὶ συναιχμαλώτους μου, οἵτινές εἰσιν ἐπίσημοι ἐν τοῖς ἀποστόλοις, οἳ καὶ πρὸ ἐμοῦ γέγοναν ἐν Χριστῷ.
\vs{8}~ἀσπάσασθε Ἀμπλιᾶτον τὸν ἀγαπητόν μου ἐν κυρίῳ.
\vs{9}~ἀσπάσασθε Οὐρβανὸν τὸν συνεργὸν ἡμῶν ἐν Χριστῷ καὶ Στάχυν τὸν ἀγαπητόν μου.
\vs{10}~ἀσπάσασθε Ἀπελλῆν τὸν δόκιμον ἐν Χριστῷ. ἀσπάσασθε τοὺς ἐκ τῶν Ἀριστοβούλου.
\vs{11}~ἀσπάσασθε Ἡρῳδίωνα τὸν συγγενῆ μου. ἀσπάσασθε τοὺς ἐκ τῶν Ναρκίσσου τοὺς ὄντας ἐν κυρίῳ.
\vs{12}~ἀσπάσασθε Τρύφαιναν καὶ Τρυφῶσαν τὰς κοπιώσας ἐν κυρίῳ. ἀσπάσασθε Περσίδα τὴν ἀγαπητήν, ἥτις πολλὰ ἐκοπίασεν ἐν κυρίῳ.
\vs{13}~ἀσπάσασθε Ῥοῦφον τὸν ἐκλεκτὸν ἐν κυρίῳ καὶ τὴν μητέρα αὐτοῦ καὶ ἐμοῦ.
\vs{14}~ἀσπάσασθε Ἀσύγκριτον, Φλέγοντα, Ἑρμῆν, Πατροβᾶν, Ἑρμᾶν καὶ τοὺς σὺν αὐτοῖς \adeA{ἀδελφούς}.
\vs{15}~ἀσπάσασθε Φιλόλογον καὶ Ἰουλίαν, Νηρέα καὶ τὴν \adeC{ἀδελφὴν} αὐτοῦ, καὶ Ὀλυμπᾶν καὶ τοὺς σὺν αὐτοῖς πάντας ἁγίους.
\vs{16}~Ἀσπάσασθε ἀλλήλους ἐν φιλήματι ἁγίῳ. Ἀσπάζονται ὑμᾶς αἱ ἐκκλησίαι πᾶσαι τοῦ Χριστοῦ.
\vs{17}~Παρακαλῶ δὲ ὑμᾶς, \adeA{ἀδελφοί}, σκοπεῖν τοὺς τὰς διχοστασίας καὶ τὰ σκάνδαλα παρὰ τὴν διδαχὴν ἣν ὑμεῖς ἐμάθετε ποιοῦντας, καὶ ἐκκλίνετε ἀπ’ αὐτῶν·
\vs{18}~οἱ γὰρ τοιοῦτοι τῷ κυρίῳ ἡμῶν Χριστῷ οὐ δουλεύουσιν ἀλλὰ τῇ ἑαυτῶν κοιλίᾳ, καὶ διὰ τῆς χρηστολογίας καὶ εὐλογίας ἐξαπατῶσι τὰς καρδίας τῶν ἀκάκων.
\vs{19}~ἡ γὰρ ὑμῶν ὑπακοὴ εἰς πάντας ἀφίκετο· ἐφ’ ὑμῖν οὖν χαίρω, θέλω δὲ ὑμᾶς σοφοὺς εἶναι εἰς τὸ ἀγαθόν, ἀκεραίους δὲ εἰς τὸ κακόν.
\vs{20}~ὁ δὲ θεὸς τῆς εἰρήνης συντρίψει τὸν Σατανᾶν ὑπὸ τοὺς πόδας ὑμῶν ἐν τάχει. ἡ χάρις τοῦ κυρίου ἡμῶν Ἰησοῦ Χριστοῦ μεθ’ ὑμῶν.
\vs{21}~Ἀσπάζεται ὑμᾶς Τιμόθεος ὁ συνεργός μου, καὶ Λούκιος καὶ Ἰάσων καὶ Σωσίπατρος οἱ συγγενεῖς μου.
\vs{22}~ἀσπάζομαι ὑμᾶς ἐγὼ Τέρτιος ὁ γράψας τὴν ἐπιστολὴν ἐν κυρίῳ.
\vs{23}~ἀσπάζεται ὑμᾶς Γάϊος ὁ ξένος μου καὶ ὅλης τῆς ἐκκλησίας. ἀσπάζεται ὑμᾶς Ἔραστος ὁ οἰκονόμος τῆς πόλεως καὶ Κούαρτος ὁ \adeA{ἀδελφός}.
\vs{24}~Ἡ χάρις τοῦ κυρίου ἡμῶν Ἰησοῦ Χριστοῦ μετὰ πάντων ὑμῶν. Ἀμήν.

\booksection{1 Corinthians}{Πρὸς Κορινθίους α΄}

\chapheader{1}
\vs{1}~Παῦλος κλητὸς ἀπόστολος Χριστοῦ Ἰησοῦ διὰ θελήματος θεοῦ καὶ Σωσθένης ὁ \adeA{ἀδελφὸς}
\vs{2}~τῇ ἐκκλησίᾳ τοῦ θεοῦ, ἡγιασμένοις ἐν Χριστῷ Ἰησοῦ, τῇ οὔσῃ ἐν Κορίνθῳ, κλητοῖς ἁγίοις, σὺν πᾶσιν τοῖς ἐπικαλουμένοις τὸ ὄνομα τοῦ κυρίου ἡμῶν Ἰησοῦ Χριστοῦ ἐν παντὶ τόπῳ αὐτῶν καὶ ἡμῶν·
\vs{3}~χάρις ὑμῖν καὶ εἰρήνη ἀπὸ θεοῦ πατρὸς ἡμῶν καὶ κυρίου Ἰησοῦ Χριστοῦ.
\vs{4}~Εὐχαριστῶ τῷ θεῷ μου πάντοτε περὶ ὑμῶν ἐπὶ τῇ χάριτι τοῦ θεοῦ τῇ δοθείσῃ ὑμῖν ἐν Χριστῷ Ἰησοῦ,
\vs{5}~ὅτι ἐν παντὶ ἐπλουτίσθητε ἐν αὐτῷ, ἐν παντὶ λόγῳ καὶ πάσῃ γνώσει,
\vs{6}~καθὼς τὸ μαρτύριον τοῦ Χριστοῦ ἐβεβαιώθη ἐν ὑμῖν,
\vs{7}~ὥστε ὑμᾶς μὴ ὑστερεῖσθαι ἐν μηδενὶ χαρίσματι, ἀπεκδεχομένους τὴν ἀποκάλυψιν τοῦ κυρίου ἡμῶν Ἰησοῦ Χριστοῦ·
\vs{8}~ὃς καὶ βεβαιώσει ὑμᾶς ἕως τέλους ἀνεγκλήτους ἐν τῇ ἡμέρᾳ τοῦ κυρίου ἡμῶν Ἰησοῦ Χριστοῦ.
\vs{9}~πιστὸς ὁ θεὸς δι’ οὗ ἐκλήθητε εἰς κοινωνίαν τοῦ υἱοῦ αὐτοῦ Ἰησοῦ Χριστοῦ τοῦ κυρίου ἡμῶν.
\vs{10}~Παρακαλῶ δὲ ὑμᾶς, \adeA{ἀδελφοί}, διὰ τοῦ ὀνόματος τοῦ κυρίου ἡμῶν Ἰησοῦ Χριστοῦ ἵνα τὸ αὐτὸ λέγητε πάντες, καὶ μὴ ᾖ ἐν ὑμῖν σχίσματα, ἦτε δὲ κατηρτισμένοι ἐν τῷ αὐτῷ νοῒ καὶ ἐν τῇ αὐτῇ γνώμῃ.
\vs{11}~ἐδηλώθη γάρ μοι περὶ ὑμῶν, \adeA{ἀδελφοί} μου, ὑπὸ τῶν Χλόης ὅτι ἔριδες ἐν ὑμῖν εἰσιν.
\vs{12}~λέγω δὲ τοῦτο ὅτι ἕκαστος ὑμῶν λέγει· Ἐγὼ μέν εἰμι Παύλου, Ἐγὼ δὲ Ἀπολλῶ, Ἐγὼ δὲ Κηφᾶ, Ἐγὼ δὲ Χριστοῦ.
\vs{13}~μεμέρισται ὁ Χριστός; μὴ Παῦλος ἐσταυρώθη ὑπὲρ ὑμῶν, ἢ εἰς τὸ ὄνομα Παύλου ἐβαπτίσθητε;
\vs{14}~εὐχαριστῶ ὅτι οὐδένα ὑμῶν ἐβάπτισα εἰ μὴ Κρίσπον καὶ Γάϊον,
\vs{15}~ἵνα μή τις εἴπῃ ὅτι εἰς τὸ ἐμὸν ὄνομα ἐβαπτίσθητε·
\vs{16}~ἐβάπτισα δὲ καὶ τὸν Στεφανᾶ οἶκον· λοιπὸν οὐκ οἶδα εἴ τινα ἄλλον ἐβάπτισα.
\vs{17}~οὐ γὰρ ἀπέστειλέν με Χριστὸς βαπτίζειν ἀλλὰ εὐαγγελίζεσθαι, οὐκ ἐν σοφίᾳ λόγου, ἵνα μὴ κενωθῇ ὁ σταυρὸς τοῦ Χριστοῦ.
\vs{18}~Ὁ λόγος γὰρ ὁ τοῦ σταυροῦ τοῖς μὲν ἀπολλυμένοις μωρία ἐστίν, τοῖς δὲ σῳζομένοις ἡμῖν δύναμις θεοῦ ἐστιν.
\vs{19}~γέγραπται γάρ· Ἀπολῶ τὴν σοφίαν τῶν σοφῶν, καὶ τὴν σύνεσιν τῶν συνετῶν ἀθετήσω.
\vs{20}~ποῦ σοφός; ποῦ γραμματεύς; ποῦ συζητητὴς τοῦ αἰῶνος τούτου; οὐχὶ ἐμώρανεν ὁ θεὸς τὴν σοφίαν τοῦ κόσμου;
\vs{21}~ἐπειδὴ γὰρ ἐν τῇ σοφίᾳ τοῦ θεοῦ οὐκ ἔγνω ὁ κόσμος διὰ τῆς σοφίας τὸν θεόν, εὐδόκησεν ὁ θεὸς διὰ τῆς μωρίας τοῦ κηρύγματος σῶσαι τοὺς πιστεύοντας.
\vs{22}~ἐπειδὴ καὶ Ἰουδαῖοι σημεῖα αἰτοῦσιν καὶ Ἕλληνες σοφίαν ζητοῦσιν·
\vs{23}~ἡμεῖς δὲ κηρύσσομεν Χριστὸν ἐσταυρωμένον, Ἰουδαίοις μὲν σκάνδαλον ἔθνεσιν δὲ μωρίαν,
\vs{24}~αὐτοῖς δὲ τοῖς κλητοῖς, Ἰουδαίοις τε καὶ Ἕλλησιν, Χριστὸν θεοῦ δύναμιν καὶ θεοῦ σοφίαν.
\vs{25}~ὅτι τὸ μωρὸν τοῦ θεοῦ σοφώτερον τῶν ἀνθρώπων ἐστίν, καὶ τὸ ἀσθενὲς τοῦ θεοῦ ἰσχυρότερον τῶν ἀνθρώπων.
\vs{26}~Βλέπετε γὰρ τὴν κλῆσιν ὑμῶν, \adeA{ἀδελφοί}, ὅτι οὐ πολλοὶ σοφοὶ κατὰ σάρκα, οὐ πολλοὶ δυνατοί, οὐ πολλοὶ εὐγενεῖς·
\vs{27}~ἀλλὰ τὰ μωρὰ τοῦ κόσμου ἐξελέξατο ὁ θεός, ἵνα καταισχύνῃ τοὺς σοφούς, καὶ τὰ ἀσθενῆ τοῦ κόσμου ἐξελέξατο ὁ θεός, ἵνα καταισχύνῃ τὰ ἰσχυρά,
\vs{28}~καὶ τὰ ἀγενῆ τοῦ κόσμου καὶ τὰ ἐξουθενημένα ἐξελέξατο ὁ θεός, τὰ μὴ ὄντα, ἵνα τὰ ὄντα καταργήσῃ,
\vs{29}~ὅπως μὴ καυχήσηται πᾶσα σὰρξ ἐνώπιον τοῦ θεοῦ.
\vs{30}~ἐξ αὐτοῦ δὲ ὑμεῖς ἐστε ἐν Χριστῷ Ἰησοῦ, ὃς ἐγενήθη σοφία ἡμῖν ἀπὸ θεοῦ, δικαιοσύνη τε καὶ ἁγιασμὸς καὶ ἀπολύτρωσις,
\vs{31}~ἵνα καθὼς γέγραπται· Ὁ καυχώμενος ἐν κυρίῳ καυχάσθω.

\chapheader{2}
\vs{1}~Κἀγὼ ἐλθὼν πρὸς ὑμᾶς, \adeA{ἀδελφοί}, ἦλθον οὐ καθ’ ὑπεροχὴν λόγου ἢ σοφίας καταγγέλλων ὑμῖν τὸ μαρτύριον τοῦ θεοῦ.
\vs{2}~οὐ γὰρ ἔκρινά τι εἰδέναι ἐν ὑμῖν εἰ μὴ Ἰησοῦν Χριστὸν καὶ τοῦτον ἐσταυρωμένον·
\vs{3}~κἀγὼ ἐν ἀσθενείᾳ καὶ ἐν φόβῳ καὶ ἐν τρόμῳ πολλῷ ἐγενόμην πρὸς ὑμᾶς,
\vs{4}~καὶ ὁ λόγος μου καὶ τὸ κήρυγμά μου οὐκ ἐν πειθοῖ σοφίας ἀλλ’ ἐν ἀποδείξει πνεύματος καὶ δυνάμεως,
\vs{5}~ἵνα ἡ πίστις ὑμῶν μὴ ᾖ ἐν σοφίᾳ ἀνθρώπων ἀλλ’ ἐν δυνάμει θεοῦ.
\vs{6}~Σοφίαν δὲ λαλοῦμεν ἐν τοῖς τελείοις, σοφίαν δὲ οὐ τοῦ αἰῶνος τούτου οὐδὲ τῶν ἀρχόντων τοῦ αἰῶνος τούτου τῶν καταργουμένων·
\vs{7}~ἀλλὰ λαλοῦμεν θεοῦ σοφίαν ἐν μυστηρίῳ, τὴν ἀποκεκρυμμένην, ἣν προώρισεν ὁ θεὸς πρὸ τῶν αἰώνων εἰς δόξαν ἡμῶν·
\vs{8}~ἣν οὐδεὶς τῶν ἀρχόντων τοῦ αἰῶνος τούτου ἔγνωκεν, εἰ γὰρ ἔγνωσαν, οὐκ ἂν τὸν κύριον τῆς δόξης ἐσταύρωσαν·
\vs{9}~ἀλλὰ καθὼς γέγραπται· Ἃ ὀφθαλμὸς οὐκ εἶδεν καὶ οὖς οὐκ ἤκουσεν καὶ ἐπὶ καρδίαν ἀνθρώπου οὐκ ἀνέβη, ὅσα ἡτοίμασεν ὁ θεὸς τοῖς ἀγαπῶσιν αὐτόν.
\vs{10}~ἡμῖν γὰρ ἀπεκάλυψεν ὁ θεὸς διὰ τοῦ πνεύματος, τὸ γὰρ πνεῦμα πάντα ἐραυνᾷ, καὶ τὰ βάθη τοῦ θεοῦ.
\vs{11}~τίς γὰρ οἶδεν ἀνθρώπων τὰ τοῦ ἀνθρώπου εἰ μὴ τὸ πνεῦμα τοῦ ἀνθρώπου τὸ ἐν αὐτῷ; οὕτως καὶ τὰ τοῦ θεοῦ οὐδεὶς ἔγνωκεν εἰ μὴ τὸ πνεῦμα τοῦ θεοῦ.
\vs{12}~ἡμεῖς δὲ οὐ τὸ πνεῦμα τοῦ κόσμου ἐλάβομεν ἀλλὰ τὸ πνεῦμα τὸ ἐκ τοῦ θεοῦ, ἵνα εἰδῶμεν τὰ ὑπὸ τοῦ θεοῦ χαρισθέντα ἡμῖν·
\vs{13}~ἃ καὶ λαλοῦμεν οὐκ ἐν διδακτοῖς ἀνθρωπίνης σοφίας λόγοις, ἀλλ’ ἐν διδακτοῖς πνεύματος, πνευματικοῖς πνευματικὰ συγκρίνοντες.
\vs{14}~Ψυχικὸς δὲ ἄνθρωπος οὐ δέχεται τὰ τοῦ πνεύματος τοῦ θεοῦ, μωρία γὰρ αὐτῷ ἐστίν, καὶ οὐ δύναται γνῶναι, ὅτι πνευματικῶς ἀνακρίνεται·
\vs{15}~ὁ δὲ πνευματικὸς ἀνακρίνει τὰ πάντα, αὐτὸς δὲ ὑπ’ οὐδενὸς ἀνακρίνεται.
\vs{16}~τίς γὰρ ἔγνω νοῦν κυρίου, ὃς συμβιβάσει αὐτόν; ἡμεῖς δὲ νοῦν Χριστοῦ ἔχομεν.

\chapheader{3}
\vs{1}~Κἀγώ, \adeA{ἀδελφοί}, οὐκ ἠδυνήθην λαλῆσαι ὑμῖν ὡς πνευματικοῖς ἀλλ’ ὡς σαρκίνοις, ὡς νηπίοις ἐν Χριστῷ.
\vs{2}~γάλα ὑμᾶς ἐπότισα, οὐ βρῶμα, οὔπω γὰρ ἐδύνασθε. ἀλλ’ οὐδὲ ἔτι νῦν δύνασθε,
\vs{3}~ἔτι γὰρ σαρκικοί ἐστε. ὅπου γὰρ ἐν ὑμῖν ζῆλος καὶ ἔρις, οὐχὶ σαρκικοί ἐστε καὶ κατὰ ἄνθρωπον περιπατεῖτε;
\vs{4}~ὅταν γὰρ λέγῃ τις· Ἐγὼ μέν εἰμι Παύλου, ἕτερος δέ· Ἐγὼ Ἀπολλῶ, οὐκ ἄνθρωποί ἐστε;
\vs{5}~Τί οὖν ἐστιν Ἀπολλῶς; τί δέ ἐστιν Παῦλος; διάκονοι δι’ ὧν ἐπιστεύσατε, καὶ ἑκάστῳ ὡς ὁ κύριος ἔδωκεν.
\vs{6}~ἐγὼ ἐφύτευσα, Ἀπολλῶς ἐπότισεν, ἀλλὰ ὁ θεὸς ηὔξανεν·
\vs{7}~ὥστε οὔτε ὁ φυτεύων ἐστίν τι οὔτε ὁ ποτίζων, ἀλλ’ ὁ αὐξάνων θεός.
\vs{8}~ὁ φυτεύων δὲ καὶ ὁ ποτίζων ἕν εἰσιν, ἕκαστος δὲ τὸν ἴδιον μισθὸν λήμψεται κατὰ τὸν ἴδιον κόπον,
\vs{9}~θεοῦ γάρ ἐσμεν συνεργοί· θεοῦ γεώργιον, θεοῦ οἰκοδομή ἐστε.
\vs{10}~Κατὰ τὴν χάριν τοῦ θεοῦ τὴν δοθεῖσάν μοι ὡς σοφὸς ἀρχιτέκτων θεμέλιον ἔθηκα, ἄλλος δὲ ἐποικοδομεῖ. ἕκαστος δὲ βλεπέτω πῶς ἐποικοδομεῖ·
\vs{11}~θεμέλιον γὰρ ἄλλον οὐδεὶς δύναται θεῖναι παρὰ τὸν κείμενον, ὅς ἐστιν Ἰησοῦς Χριστός·
\vs{12}~εἰ δέ τις ἐποικοδομεῖ ἐπὶ τὸν θεμέλιον χρυσόν, ἄργυρον, λίθους τιμίους, ξύλα, χόρτον, καλάμην,
\vs{13}~ἑκάστου τὸ ἔργον φανερὸν γενήσεται, ἡ γὰρ ἡμέρα δηλώσει· ὅτι ἐν πυρὶ ἀποκαλύπτεται, καὶ ἑκάστου τὸ ἔργον ὁποῖόν ἐστιν τὸ πῦρ αὐτὸ δοκιμάσει.
\vs{14}~εἴ τινος τὸ ἔργον μενεῖ ὃ ἐποικοδόμησεν, μισθὸν λήμψεται·
\vs{15}~εἴ τινος τὸ ἔργον κατακαήσεται, ζημιωθήσεται, αὐτὸς δὲ σωθήσεται, οὕτως δὲ ὡς διὰ πυρός.
\vs{16}~Οὐκ οἴδατε ὅτι ναὸς θεοῦ ἐστε καὶ τὸ πνεῦμα τοῦ θεοῦ οἰκεῖ ἐν ὑμῖν;
\vs{17}~εἴ τις τὸν ναὸν τοῦ θεοῦ φθείρει, φθερεῖ τοῦτον ὁ θεός· ὁ γὰρ ναὸς τοῦ θεοῦ ἅγιός ἐστιν, οἵτινές ἐστε ὑμεῖς.
\vs{18}~Μηδεὶς ἑαυτὸν ἐξαπατάτω· εἴ τις δοκεῖ σοφὸς εἶναι ἐν ὑμῖν ἐν τῷ αἰῶνι τούτῳ, μωρὸς γενέσθω, ἵνα γένηται σοφός,
\vs{19}~ἡ γὰρ σοφία τοῦ κόσμου τούτου μωρία παρὰ τῷ θεῷ ἐστιν· γέγραπται γάρ· Ὁ δρασσόμενος τοὺς σοφοὺς ἐν τῇ πανουργίᾳ αὐτῶν·
\vs{20}~καὶ πάλιν· Κύριος γινώσκει τοὺς διαλογισμοὺς τῶν σοφῶν ὅτι εἰσὶν μάταιοι.
\vs{21}~ὥστε μηδεὶς καυχάσθω ἐν ἀνθρώποις· πάντα γὰρ ὑμῶν ἐστιν,
\vs{22}~εἴτε Παῦλος εἴτε Ἀπολλῶς εἴτε Κηφᾶς εἴτε κόσμος εἴτε ζωὴ εἴτε θάνατος εἴτε ἐνεστῶτα εἴτε μέλλοντα, πάντα ὑμῶν,
\vs{23}~ὑμεῖς δὲ Χριστοῦ, Χριστὸς δὲ θεοῦ.

\chapheader{4}
\vs{1}~Οὕτως ἡμᾶς λογιζέσθω ἄνθρωπος ὡς ὑπηρέτας Χριστοῦ καὶ οἰκονόμους μυστηρίων θεοῦ.
\vs{2}~ὧδε λοιπὸν ζητεῖται ἐν τοῖς οἰκονόμοις ἵνα πιστός τις εὑρεθῇ.
\vs{3}~ἐμοὶ δὲ εἰς ἐλάχιστόν ἐστιν, ἵνα ὑφ’ ὑμῶν ἀνακριθῶ ἢ ὑπὸ ἀνθρωπίνης ἡμέρας· ἀλλ’ οὐδὲ ἐμαυτὸν ἀνακρίνω·
\vs{4}~οὐδὲν γὰρ ἐμαυτῷ σύνοιδα, ἀλλ’ οὐκ ἐν τούτῳ δεδικαίωμαι, ὁ δὲ ἀνακρίνων με κύριός ἐστιν.
\vs{5}~ὥστε μὴ πρὸ καιροῦ τι κρίνετε, ἕως ἂν ἔλθῃ ὁ κύριος, ὃς καὶ φωτίσει τὰ κρυπτὰ τοῦ σκότους καὶ φανερώσει τὰς βουλὰς τῶν καρδιῶν, καὶ τότε ὁ ἔπαινος γενήσεται ἑκάστῳ ἀπὸ τοῦ θεοῦ.
\vs{6}~Ταῦτα δέ, \adeA{ἀδελφοί}, μετεσχημάτισα εἰς ἐμαυτὸν καὶ Ἀπολλῶν δι’ ὑμᾶς, ἵνα ἐν ἡμῖν μάθητε τό· Μὴ ὑπὲρ ἃ γέγραπται, ἵνα μὴ εἷς ὑπὲρ τοῦ ἑνὸς φυσιοῦσθε κατὰ τοῦ ἑτέρου.
\vs{7}~τίς γάρ σε διακρίνει; τί δὲ ἔχεις ὃ οὐκ ἔλαβες; εἰ δὲ καὶ ἔλαβες, τί καυχᾶσαι ὡς μὴ λαβών;
\vs{8}~Ἤδη κεκορεσμένοι ἐστέ, ἤδη ἐπλουτήσατε, χωρὶς ἡμῶν ἐβασιλεύσατε· καὶ ὄφελόν γε ἐβασιλεύσατε, ἵνα καὶ ἡμεῖς ὑμῖν συμβασιλεύσωμεν.
\vs{9}~δοκῶ γάρ, ὁ θεὸς ἡμᾶς τοὺς ἀποστόλους ἐσχάτους ἀπέδειξεν ὡς ἐπιθανατίους, ὅτι θέατρον ἐγενήθημεν τῷ κόσμῳ καὶ ἀγγέλοις καὶ ἀνθρώποις.
\vs{10}~ἡμεῖς μωροὶ διὰ Χριστόν, ὑμεῖς δὲ φρόνιμοι ἐν Χριστῷ· ἡμεῖς ἀσθενεῖς, ὑμεῖς δὲ ἰσχυροί· ὑμεῖς ἔνδοξοι, ἡμεῖς δὲ ἄτιμοι.
\vs{11}~ἄχρι τῆς ἄρτι ὥρας καὶ πεινῶμεν καὶ διψῶμεν καὶ γυμνιτεύομεν καὶ κολαφιζόμεθα καὶ ἀστατοῦμεν
\vs{12}~καὶ κοπιῶμεν ἐργαζόμενοι ταῖς ἰδίαις χερσίν· λοιδορούμενοι εὐλογοῦμεν, διωκόμενοι ἀνεχόμεθα,
\vs{13}~δυσφημούμενοι παρακαλοῦμεν· ὡς περικαθάρματα τοῦ κόσμου ἐγενήθημεν, πάντων περίψημα ἕως ἄρτι.
\vs{14}~Οὐκ ἐντρέπων ὑμᾶς γράφω ταῦτα, ἀλλ’ ὡς τέκνα μου ἀγαπητὰ νουθετῶν·
\vs{15}~ἐὰν γὰρ μυρίους παιδαγωγοὺς ἔχητε ἐν Χριστῷ, ἀλλ’ οὐ πολλοὺς πατέρας, ἐν γὰρ Χριστῷ Ἰησοῦ διὰ τοῦ εὐαγγελίου ἐγὼ ὑμᾶς ἐγέννησα.
\vs{16}~παρακαλῶ οὖν ὑμᾶς, μιμηταί μου γίνεσθε.
\vs{17}~διὰ τοῦτο ἔπεμψα ὑμῖν Τιμόθεον, ὅς ἐστίν μου τέκνον ἀγαπητὸν καὶ πιστὸν ἐν κυρίῳ, ὃς ὑμᾶς ἀναμνήσει τὰς ὁδούς μου τὰς ἐν Χριστῷ Ἰησοῦ, καθὼς πανταχοῦ ἐν πάσῃ ἐκκλησίᾳ διδάσκω.
\vs{18}~ὡς μὴ ἐρχομένου δέ μου πρὸς ὑμᾶς ἐφυσιώθησάν τινες·
\vs{19}~ἐλεύσομαι δὲ ταχέως πρὸς ὑμᾶς, ἐὰν ὁ κύριος θελήσῃ, καὶ γνώσομαι οὐ τὸν λόγον τῶν πεφυσιωμένων ἀλλὰ τὴν δύναμιν,
\vs{20}~οὐ γὰρ ἐν λόγῳ ἡ βασιλεία τοῦ θεοῦ ἀλλ’ ἐν δυνάμει.
\vs{21}~τί θέλετε; ἐν ῥάβδῳ ἔλθω πρὸς ὑμᾶς, ἢ ἐν ἀγάπῃ πνεύματί τε πραΰτητος;

\chapheader{5}
\vs{1}~Ὅλως ἀκούεται ἐν ὑμῖν πορνεία, καὶ τοιαύτη πορνεία ἥτις οὐδὲ ἐν τοῖς ἔθνεσιν, ὥστε γυναῖκά τινα τοῦ πατρὸς ἔχειν.
\vs{2}~καὶ ὑμεῖς πεφυσιωμένοι ἐστέ, καὶ οὐχὶ μᾶλλον ἐπενθήσατε, ἵνα ἀρθῇ ἐκ μέσου ὑμῶν ὁ τὸ ἔργον τοῦτο ποιήσας;
\vs{3}~Ἐγὼ μὲν γάρ, ἀπὼν τῷ σώματι παρὼν δὲ τῷ πνεύματι, ἤδη κέκρικα ὡς παρὼν τὸν οὕτως τοῦτο κατεργασάμενον
\vs{4}~ἐν τῷ ὀνόματι τοῦ κυρίου ἡμῶν Ἰησοῦ, συναχθέντων ὑμῶν καὶ τοῦ ἐμοῦ πνεύματος σὺν τῇ δυνάμει τοῦ κυρίου ἡμῶν Ἰησοῦ,
\vs{5}~παραδοῦναι τὸν τοιοῦτον τῷ Σατανᾷ εἰς ὄλεθρον τῆς σαρκός, ἵνα τὸ πνεῦμα σωθῇ ἐν τῇ ἡμέρᾳ τοῦ κυρίου.
\vs{6}~Οὐ καλὸν τὸ καύχημα ὑμῶν. οὐκ οἴδατε ὅτι μικρὰ ζύμη ὅλον τὸ φύραμα ζυμοῖ;
\vs{7}~ἐκκαθάρατε τὴν παλαιὰν ζύμην, ἵνα ἦτε νέον φύραμα, καθώς ἐστε ἄζυμοι. καὶ γὰρ τὸ πάσχα ἡμῶν ἐτύθη Χριστός·
\vs{8}~ὥστε ἑορτάζωμεν, μὴ ἐν ζύμῃ παλαιᾷ μηδὲ ἐν ζύμῃ κακίας καὶ πονηρίας, ἀλλ’ ἐν ἀζύμοις εἰλικρινείας καὶ ἀληθείας.
\vs{9}~Ἔγραψα ὑμῖν ἐν τῇ ἐπιστολῇ μὴ συναναμίγνυσθαι πόρνοις,
\vs{10}~οὐ πάντως τοῖς πόρνοις τοῦ κόσμου τούτου ἢ τοῖς πλεονέκταις καὶ ἅρπαξιν ἢ εἰδωλολάτραις, ἐπεὶ ὠφείλετε ἄρα ἐκ τοῦ κόσμου ἐξελθεῖν.
\vs{11}~νῦν δὲ ἔγραψα ὑμῖν μὴ συναναμίγνυσθαι ἐάν τις \adeA{ἀδελφὸς} ὀνομαζόμενος ᾖ πόρνος ἢ πλεονέκτης ἢ εἰδωλολάτρης ἢ λοίδορος ἢ μέθυσος ἢ ἅρπαξ, τῷ τοιούτῳ μηδὲ συνεσθίειν.
\vs{12}~τί γάρ μοι τοὺς ἔξω κρίνειν; οὐχὶ τοὺς ἔσω ὑμεῖς κρίνετε,
\vs{13}~τοὺς δὲ ἔξω ὁ θεὸς κρίνει; ἐξάρατε τὸν πονηρὸν ἐξ ὑμῶν αὐτῶν.

\chapheader{6}
\vs{1}~Τολμᾷ τις ὑμῶν πρᾶγμα ἔχων πρὸς τὸν ἕτερον κρίνεσθαι ἐπὶ τῶν ἀδίκων, καὶ οὐχὶ ἐπὶ τῶν ἁγίων;
\vs{2}~ἢ οὐκ οἴδατε ὅτι οἱ ἅγιοι τὸν κόσμον κρινοῦσιν; καὶ εἰ ἐν ὑμῖν κρίνεται ὁ κόσμος, ἀνάξιοί ἐστε κριτηρίων ἐλαχίστων;
\vs{3}~οὐκ οἴδατε ὅτι ἀγγέλους κρινοῦμεν, μήτιγε βιωτικά;
\vs{4}~βιωτικὰ μὲν οὖν κριτήρια ἐὰν ἔχητε, τοὺς ἐξουθενημένους ἐν τῇ ἐκκλησίᾳ, τούτους καθίζετε;
\vs{5}~πρὸς ἐντροπὴν ὑμῖν λέγω. οὕτως οὐκ ἔνι ἐν ὑμῖν οὐδεὶς σοφὸς ὃς δυνήσεται διακρῖναι ἀνὰ μέσον τοῦ \adeA{ἀδελφοῦ} αὐτοῦ,
\vs{6}~ἀλλὰ \adeA{ἀδελφὸς} μετὰ \adeA{ἀδελφοῦ} κρίνεται, καὶ τοῦτο ἐπὶ ἀπίστων;
\vs{7}~ἤδη μὲν οὖν ὅλως ἥττημα ὑμῖν ἐστιν ὅτι κρίματα ἔχετε μεθ’ ἑαυτῶν· διὰ τί οὐχὶ μᾶλλον ἀδικεῖσθε; διὰ τί οὐχὶ μᾶλλον ἀποστερεῖσθε;
\vs{8}~ἀλλὰ ὑμεῖς ἀδικεῖτε καὶ ἀποστερεῖτε, καὶ τοῦτο \adeA{ἀδελφούς}.
\vs{9}~Ἢ οὐκ οἴδατε ὅτι ἄδικοι θεοῦ βασιλείαν οὐ κληρονομήσουσιν; μὴ πλανᾶσθε· οὔτε πόρνοι οὔτε εἰδωλολάτραι οὔτε μοιχοὶ οὔτε μαλακοὶ οὔτε ἀρσενοκοῖται
\vs{10}~οὔτε κλέπται οὔτε πλεονέκται, οὐ μέθυσοι, οὐ λοίδοροι, οὐχ ἅρπαγες βασιλείαν θεοῦ κληρονομήσουσιν.
\vs{11}~καὶ ταῦτά τινες ἦτε· ἀλλὰ ἀπελούσασθε, ἀλλὰ ἡγιάσθητε, ἀλλὰ ἐδικαιώθητε ἐν τῷ ὀνόματι τοῦ κυρίου Ἰησοῦ καὶ ἐν τῷ πνεύματι τοῦ θεοῦ ἡμῶν.
\vs{12}~Πάντα μοι ἔξεστιν· ἀλλ’ οὐ πάντα συμφέρει. πάντα μοι ἔξεστιν· ἀλλ’ οὐκ ἐγὼ ἐξουσιασθήσομαι ὑπό τινος.
\vs{13}~τὰ βρώματα τῇ κοιλίᾳ, καὶ ἡ κοιλία τοῖς βρώμασιν· ὁ δὲ θεὸς καὶ ταύτην καὶ ταῦτα καταργήσει. τὸ δὲ σῶμα οὐ τῇ πορνείᾳ ἀλλὰ τῷ κυρίῳ, καὶ ὁ κύριος τῷ σώματι·
\vs{14}~ὁ δὲ θεὸς καὶ τὸν κύριον ἤγειρεν καὶ ἡμᾶς ἐξεγερεῖ διὰ τῆς δυνάμεως αὐτοῦ.
\vs{15}~οὐκ οἴδατε ὅτι τὰ σώματα ὑμῶν μέλη Χριστοῦ ἐστιν; ἄρας οὖν τὰ μέλη τοῦ Χριστοῦ ποιήσω πόρνης μέλη; μὴ γένοιτο.
\vs{16}~ἢ οὐκ οἴδατε ὅτι ὁ κολλώμενος τῇ πόρνῃ ἓν σῶμά ἐστιν; Ἔσονται γάρ, φησίν, οἱ δύο εἰς σάρκα μίαν.
\vs{17}~ὁ δὲ κολλώμενος τῷ κυρίῳ ἓν πνεῦμά ἐστιν.
\vs{18}~φεύγετε τὴν πορνείαν· πᾶν ἁμάρτημα ὃ ἐὰν ποιήσῃ ἄνθρωπος ἐκτὸς τοῦ σώματός ἐστιν, ὁ δὲ πορνεύων εἰς τὸ ἴδιον σῶμα ἁμαρτάνει.
\vs{19}~ἢ οὐκ οἴδατε ὅτι τὸ σῶμα ὑμῶν ναὸς τοῦ ἐν ὑμῖν ἁγίου πνεύματός ἐστιν, οὗ ἔχετε ἀπὸ θεοῦ; καὶ οὐκ ἐστὲ ἑαυτῶν,
\vs{20}~ἠγοράσθητε γὰρ τιμῆς· δοξάσατε δὴ τὸν θεὸν ἐν τῷ σώματι ὑμῶν.

\chapheader{7}
\vs{1}~Περὶ δὲ ὧν ἐγράψατε, καλὸν ἀνθρώπῳ γυναικὸς μὴ ἅπτεσθαι·
\vs{2}~διὰ δὲ τὰς πορνείας ἕκαστος τὴν ἑαυτοῦ γυναῖκα ἐχέτω, καὶ ἑκάστη τὸν ἴδιον ἄνδρα ἐχέτω.
\vs{3}~τῇ γυναικὶ ὁ ἀνὴρ τὴν ὀφειλὴν ἀποδιδότω, ὁμοίως δὲ καὶ ἡ γυνὴ τῷ ἀνδρί.
\vs{4}~ἡ γυνὴ τοῦ ἰδίου σώματος οὐκ ἐξουσιάζει ἀλλὰ ὁ ἀνήρ· ὁμοίως δὲ καὶ ὁ ἀνὴρ τοῦ ἰδίου σώματος οὐκ ἐξουσιάζει ἀλλὰ ἡ γυνή.
\vs{5}~μὴ ἀποστερεῖτε ἀλλήλους, εἰ μήτι ἂν ἐκ συμφώνου πρὸς καιρὸν ἵνα σχολάσητε τῇ προσευχῇ καὶ πάλιν ἐπὶ τὸ αὐτὸ ἦτε, ἵνα μὴ πειράζῃ ὑμᾶς ὁ Σατανᾶς διὰ τὴν ἀκρασίαν ὑμῶν.
\vs{6}~τοῦτο δὲ λέγω κατὰ συγγνώμην, οὐ κατ’ ἐπιταγήν.
\vs{7}~θέλω δὲ πάντας ἀνθρώπους εἶναι ὡς καὶ ἐμαυτόν· ἀλλὰ ἕκαστος ἴδιον ἔχει χάρισμα ἐκ θεοῦ, ὁ μὲν οὕτως, ὁ δὲ οὕτως.
\vs{8}~Λέγω δὲ τοῖς ἀγάμοις καὶ ταῖς χήραις, καλὸν αὐτοῖς ἐὰν μείνωσιν ὡς κἀγώ·
\vs{9}~εἰ δὲ οὐκ ἐγκρατεύονται, γαμησάτωσαν, κρεῖττον γάρ ἐστιν γαμῆσαι ἢ πυροῦσθαι.
\vs{10}~Τοῖς δὲ γεγαμηκόσιν παραγγέλλω, οὐκ ἐγὼ ἀλλὰ ὁ κύριος, γυναῖκα ἀπὸ ἀνδρὸς μὴ χωρισθῆναι—
\vs{11}~ἐὰν δὲ καὶ χωρισθῇ, μενέτω ἄγαμος ἢ τῷ ἀνδρὶ καταλλαγήτω— καὶ ἄνδρα γυναῖκα μὴ ἀφιέναι.
\vs{12}~Τοῖς δὲ λοιποῖς λέγω ἐγώ, οὐχ ὁ κύριος· εἴ τις \adeA{ἀδελφὸς} γυναῖκα ἔχει ἄπιστον, καὶ αὕτη συνευδοκεῖ οἰκεῖν μετ’ αὐτοῦ, μὴ ἀφιέτω αὐτήν·
\vs{13}~καὶ γυνὴ εἴ τις ἔχει ἄνδρα ἄπιστον, καὶ οὗτος συνευδοκεῖ οἰκεῖν μετ’ αὐτῆς, μὴ ἀφιέτω τὸν ἄνδρα.
\vs{14}~ἡγίασται γὰρ ὁ ἀνὴρ ὁ ἄπιστος ἐν τῇ γυναικί, καὶ ἡγίασται ἡ γυνὴ ἡ ἄπιστος ἐν τῷ \adeA{ἀδελφῷ}· ἐπεὶ ἄρα τὰ τέκνα ὑμῶν ἀκάθαρτά ἐστιν, νῦν δὲ ἅγιά ἐστιν.
\vs{15}~εἰ δὲ ὁ ἄπιστος χωρίζεται, χωριζέσθω· οὐ δεδούλωται ὁ \adeA{ἀδελφὸς} ἢ ἡ \adeA{ἀδελφὴ} ἐν τοῖς τοιούτοις, ἐν δὲ εἰρήνῃ κέκληκεν ἡμᾶς ὁ θεός.
\vs{16}~τί γὰρ οἶδας, γύναι, εἰ τὸν ἄνδρα σώσεις; ἢ τί οἶδας, ἄνερ, εἰ τὴν γυναῖκα σώσεις;
\vs{17}~Εἰ μὴ ἑκάστῳ ὡς ἐμέρισεν ὁ κύριος, ἕκαστον ὡς κέκληκεν ὁ θεός, οὕτως περιπατείτω· καὶ οὕτως ἐν ταῖς ἐκκλησίαις πάσαις διατάσσομαι.
\vs{18}~περιτετμημένος τις ἐκλήθη; μὴ ἐπισπάσθω· ἐν ἀκροβυστίᾳ κέκληταί τις; μὴ περιτεμνέσθω.
\vs{19}~ἡ περιτομὴ οὐδέν ἐστιν, καὶ ἡ ἀκροβυστία οὐδέν ἐστιν, ἀλλὰ τήρησις ἐντολῶν θεοῦ.
\vs{20}~ἕκαστος ἐν τῇ κλήσει ᾗ ἐκλήθη ἐν ταύτῃ μενέτω.
\vs{21}~Δοῦλος ἐκλήθης; μή σοι μελέτω· ἀλλ’ εἰ καὶ δύνασαι ἐλεύθερος γενέσθαι, μᾶλλον χρῆσαι.
\vs{22}~ὁ γὰρ ἐν κυρίῳ κληθεὶς δοῦλος ἀπελεύθερος κυρίου ἐστίν· ὁμοίως ὁ ἐλεύθερος κληθεὶς δοῦλός ἐστιν Χριστοῦ.
\vs{23}~τιμῆς ἠγοράσθητε· μὴ γίνεσθε δοῦλοι ἀνθρώπων.
\vs{24}~ἕκαστος ἐν ᾧ ἐκλήθη, \adeA{ἀδελφοί}, ἐν τούτῳ μενέτω παρὰ θεῷ.
\vs{25}~Περὶ δὲ τῶν παρθένων ἐπιταγὴν κυρίου οὐκ ἔχω, γνώμην δὲ δίδωμι ὡς ἠλεημένος ὑπὸ κυρίου πιστὸς εἶναι.
\vs{26}~νομίζω οὖν τοῦτο καλὸν ὑπάρχειν διὰ τὴν ἐνεστῶσαν ἀνάγκην, ὅτι καλὸν ἀνθρώπῳ τὸ οὕτως εἶναι.
\vs{27}~δέδεσαι γυναικί; μὴ ζήτει λύσιν· λέλυσαι ἀπὸ γυναικός; μὴ ζήτει γυναῖκα·
\vs{28}~ἐὰν δὲ καὶ γαμήσῃς, οὐχ ἥμαρτες. καὶ ἐὰν γήμῃ ἡ παρθένος, οὐχ ἥμαρτεν. θλῖψιν δὲ τῇ σαρκὶ ἕξουσιν οἱ τοιοῦτοι, ἐγὼ δὲ ὑμῶν φείδομαι.
\vs{29}~τοῦτο δέ φημι, \adeA{ἀδελφοί}, ὁ καιρὸς συνεσταλμένος ἐστίν· τὸ λοιπὸν ἵνα καὶ οἱ ἔχοντες γυναῖκας ὡς μὴ ἔχοντες ὦσιν,
\vs{30}~καὶ οἱ κλαίοντες ὡς μὴ κλαίοντες, καὶ οἱ χαίροντες ὡς μὴ χαίροντες, καὶ οἱ ἀγοράζοντες ὡς μὴ κατέχοντες,
\vs{31}~καὶ οἱ χρώμενοι τὸν κόσμον ὡς μὴ καταχρώμενοι· παράγει γὰρ τὸ σχῆμα τοῦ κόσμου τούτου.
\vs{32}~Θέλω δὲ ὑμᾶς ἀμερίμνους εἶναι. ὁ ἄγαμος μεριμνᾷ τὰ τοῦ κυρίου, πῶς ἀρέσῃ τῷ κυρίῳ·
\vs{33}~ὁ δὲ γαμήσας μεριμνᾷ τὰ τοῦ κόσμου, πῶς ἀρέσῃ τῇ γυναικί,
\vs{34}~καὶ μεμέρισται. καὶ ἡ γυνὴ ἡ ἄγαμος καὶ ἡ παρθένος μεριμνᾷ τὰ τοῦ κυρίου, ἵνα ᾖ ἁγία καὶ τῷ σώματι καὶ τῷ πνεύματι· ἡ δὲ γαμήσασα μεριμνᾷ τὰ τοῦ κόσμου, πῶς ἀρέσῃ τῷ ἀνδρί.
\vs{35}~τοῦτο δὲ πρὸς τὸ ὑμῶν αὐτῶν σύμφορον λέγω, οὐχ ἵνα βρόχον ὑμῖν ἐπιβάλω, ἀλλὰ πρὸς τὸ εὔσχημον καὶ εὐπάρεδρον τῷ κυρίῳ ἀπερισπάστως.
\vs{36}~Εἰ δέ τις ἀσχημονεῖν ἐπὶ τὴν παρθένον αὐτοῦ νομίζει ἐὰν ᾖ ὑπέρακμος, καὶ οὕτως ὀφείλει γίνεσθαι, ὃ θέλει ποιείτω· οὐχ ἁμαρτάνει· γαμείτωσαν.
\vs{37}~ὃς δὲ ἕστηκεν ἐν τῇ καρδίᾳ αὐτοῦ ἑδραῖος μὴ ἔχων ἀνάγκην, ἐξουσίαν δὲ ἔχει περὶ τοῦ ἰδίου θελήματος, καὶ τοῦτο κέκρικεν ἐν τῇ ἰδίᾳ καρδίᾳ, τηρεῖν τὴν ἑαυτοῦ παρθένον, καλῶς ποιήσει·
\vs{38}~ὥστε καὶ ὁ γαμίζων τὴν παρθένον ἑαυτοῦ καλῶς ποιεῖ, καὶ ὁ μὴ γαμίζων κρεῖσσον ποιήσει.
\vs{39}~Γυνὴ δέδεται ἐφ’ ὅσον χρόνον ζῇ ὁ ἀνὴρ αὐτῆς· ἐὰν δὲ κοιμηθῇ ὁ ἀνήρ, ἐλευθέρα ἐστὶν ᾧ θέλει γαμηθῆναι, μόνον ἐν κυρίῳ·
\vs{40}~μακαριωτέρα δέ ἐστιν ἐὰν οὕτως μείνῃ, κατὰ τὴν ἐμὴν γνώμην, δοκῶ δὲ κἀγὼ πνεῦμα θεοῦ ἔχειν.

\chapheader{8}
\vs{1}~Περὶ δὲ τῶν εἰδωλοθύτων, οἴδαμεν ὅτι πάντες γνῶσιν ἔχομεν. ἡ γνῶσις φυσιοῖ, ἡ δὲ ἀγάπη οἰκοδομεῖ.
\vs{2}~εἴ τις δοκεῖ ἐγνωκέναι τι, οὔπω ἔγνω καθὼς δεῖ γνῶναι·
\vs{3}~εἰ δέ τις ἀγαπᾷ τὸν θεόν, οὗτος ἔγνωσται ὑπ’ αὐτοῦ.
\vs{4}~Περὶ τῆς βρώσεως οὖν τῶν εἰδωλοθύτων οἴδαμεν ὅτι οὐδὲν εἴδωλον ἐν κόσμῳ, καὶ ὅτι οὐδεὶς θεὸς εἰ μὴ εἷς.
\vs{5}~καὶ γὰρ εἴπερ εἰσὶν λεγόμενοι θεοὶ εἴτε ἐν οὐρανῷ εἴτε ἐπὶ γῆς, ὥσπερ εἰσὶν θεοὶ πολλοὶ καὶ κύριοι πολλοί,
\vs{6}~ἀλλ’ ἡμῖν εἷς θεὸς ὁ πατήρ, ἐξ οὗ τὰ πάντα καὶ ἡμεῖς εἰς αὐτόν, καὶ εἷς κύριος Ἰησοῦς Χριστός, δι’ οὗ τὰ πάντα καὶ ἡμεῖς δι’ αὐτοῦ.
\vs{7}~Ἀλλ’ οὐκ ἐν πᾶσιν ἡ γνῶσις· τινὲς δὲ τῇ συνηθείᾳ ἕως ἄρτι τοῦ εἰδώλου ὡς εἰδωλόθυτον ἐσθίουσιν, καὶ ἡ συνείδησις αὐτῶν ἀσθενὴς οὖσα μολύνεται.
\vs{8}~βρῶμα δὲ ἡμᾶς οὐ παραστήσει τῷ θεῷ· οὔτε γὰρ ἐὰν φάγωμεν, περισσεύομεν, οὔτε ἐὰν μὴ φάγωμεν, ὑστερούμεθα.
\vs{9}~βλέπετε δὲ μή πως ἡ ἐξουσία ὑμῶν αὕτη πρόσκομμα γένηται τοῖς ἀσθενέσιν.
\vs{10}~ἐὰν γάρ τις ἴδῃ σὲ τὸν ἔχοντα γνῶσιν ἐν εἰδωλείῳ κατακείμενον, οὐχὶ ἡ συνείδησις αὐτοῦ ἀσθενοῦς ὄντος οἰκοδομηθήσεται εἰς τὸ τὰ εἰδωλόθυτα ἐσθίειν;
\vs{11}~ἀπόλλυται γὰρ ὁ ἀσθενῶν ἐν τῇ σῇ γνώσει, ὁ \adeA{ἀδελφὸς} δι’ ὃν Χριστὸς ἀπέθανεν.
\vs{12}~οὕτως δὲ ἁμαρτάνοντες εἰς τοὺς \adeA{ἀδελφοὺς} καὶ τύπτοντες αὐτῶν τὴν συνείδησιν ἀσθενοῦσαν εἰς Χριστὸν ἁμαρτάνετε.
\vs{13}~διόπερ εἰ βρῶμα σκανδαλίζει τὸν \adeA{ἀδελφόν} μου, οὐ μὴ φάγω κρέα εἰς τὸν αἰῶνα, ἵνα μὴ τὸν \adeA{ἀδελφόν} μου σκανδαλίσω.

\chapheader{9}
\vs{1}~Οὐκ εἰμὶ ἐλεύθερος; οὐκ εἰμὶ ἀπόστολος; οὐχὶ Ἰησοῦν τὸν κύριον ἡμῶν ἑόρακα; οὐ τὸ ἔργον μου ὑμεῖς ἐστε ἐν κυρίῳ;
\vs{2}~εἰ ἄλλοις οὐκ εἰμὶ ἀπόστολος, ἀλλά γε ὑμῖν εἰμι, ἡ γὰρ σφραγίς μου τῆς ἀποστολῆς ὑμεῖς ἐστε ἐν κυρίῳ.
\vs{3}~Ἡ ἐμὴ ἀπολογία τοῖς ἐμὲ ἀνακρίνουσίν ἐστιν αὕτη.
\vs{4}~μὴ οὐκ ἔχομεν ἐξουσίαν φαγεῖν καὶ πεῖν;
\vs{5}~μὴ οὐκ ἔχομεν ἐξουσίαν \adeA{ἀδελφὴν} γυναῖκα περιάγειν, ὡς καὶ οἱ λοιποὶ ἀπόστολοι καὶ οἱ \adeC{ἀδελφοὶ} τοῦ κυρίου καὶ Κηφᾶς;
\vs{6}~ἢ μόνος ἐγὼ καὶ Βαρναβᾶς οὐκ ἔχομεν ἐξουσίαν μὴ ἐργάζεσθαι;
\vs{7}~τίς στρατεύεται ἰδίοις ὀψωνίοις ποτέ; τίς φυτεύει ἀμπελῶνα καὶ τὸν καρπὸν αὐτοῦ οὐκ ἐσθίει; τίς ποιμαίνει ποίμνην καὶ ἐκ τοῦ γάλακτος τῆς ποίμνης οὐκ ἐσθίει;
\vs{8}~Μὴ κατὰ ἄνθρωπον ταῦτα λαλῶ ἢ καὶ ὁ νόμος ταῦτα οὐ λέγει;
\vs{9}~ἐν γὰρ τῷ Μωϋσέως νόμῳ γέγραπται· Οὐ κημώσεις βοῦν ἀλοῶντα. μὴ τῶν βοῶν μέλει τῷ θεῷ,
\vs{10}~ἢ δι’ ἡμᾶς πάντως λέγει; δι’ ἡμᾶς γὰρ ἐγράφη, ὅτι ὀφείλει ἐπ’ ἐλπίδι ὁ ἀροτριῶν ἀροτριᾶν, καὶ ὁ ἀλοῶν ἐπ’ ἐλπίδι τοῦ μετέχειν.
\vs{11}~εἰ ἡμεῖς ὑμῖν τὰ πνευματικὰ ἐσπείραμεν, μέγα εἰ ἡμεῖς ὑμῶν τὰ σαρκικὰ θερίσομεν;
\vs{12}~εἰ ἄλλοι τῆς ὑμῶν ἐξουσίας μετέχουσιν, οὐ μᾶλλον ἡμεῖς; Ἀλλ’ οὐκ ἐχρησάμεθα τῇ ἐξουσίᾳ ταύτῃ, ἀλλὰ πάντα στέγομεν ἵνα μή τινα ἐγκοπὴν δῶμεν τῷ εὐαγγελίῳ τοῦ Χριστοῦ.
\vs{13}~οὐκ οἴδατε ὅτι οἱ τὰ ἱερὰ ἐργαζόμενοι τὰ ἐκ τοῦ ἱεροῦ ἐσθίουσιν, οἱ τῷ θυσιαστηρίῳ παρεδρεύοντες τῷ θυσιαστηρίῳ συμμερίζονται;
\vs{14}~οὕτως καὶ ὁ κύριος διέταξεν τοῖς τὸ εὐαγγέλιον καταγγέλλουσιν ἐκ τοῦ εὐαγγελίου ζῆν.
\vs{15}~Ἐγὼ δὲ οὐ κέχρημαι οὐδενὶ τούτων. οὐκ ἔγραψα δὲ ταῦτα ἵνα οὕτως γένηται ἐν ἐμοί, καλὸν γάρ μοι μᾶλλον ἀποθανεῖν ἤ— τὸ καύχημά μου οὐδεὶς κενώσει.
\vs{16}~ἐὰν γὰρ εὐαγγελίζωμαι, οὐκ ἔστιν μοι καύχημα, ἀνάγκη γάρ μοι ἐπίκειται· οὐαὶ γάρ μοί ἐστιν ἐὰν μὴ εὐαγγελίσωμαι.
\vs{17}~εἰ γὰρ ἑκὼν τοῦτο πράσσω, μισθὸν ἔχω· εἰ δὲ ἄκων, οἰκονομίαν πεπίστευμαι.
\vs{18}~τίς οὖν μού ἐστιν ὁ μισθός; ἵνα εὐαγγελιζόμενος ἀδάπανον θήσω τὸ εὐαγγέλιον, εἰς τὸ μὴ καταχρήσασθαι τῇ ἐξουσίᾳ μου ἐν τῷ εὐαγγελίῳ.
\vs{19}~Ἐλεύθερος γὰρ ὢν ἐκ πάντων πᾶσιν ἐμαυτὸν ἐδούλωσα, ἵνα τοὺς πλείονας κερδήσω·
\vs{20}~καὶ ἐγενόμην τοῖς Ἰουδαίοις ὡς Ἰουδαῖος, ἵνα Ἰουδαίους κερδήσω· τοῖς ὑπὸ νόμον ὡς ὑπὸ νόμον, μὴ ὢν αὐτὸς ὑπὸ νόμον, ἵνα τοὺς ὑπὸ νόμον κερδήσω·
\vs{21}~τοῖς ἀνόμοις ὡς ἄνομος, μὴ ὢν ἄνομος θεοῦ ἀλλ’ ἔννομος Χριστοῦ, ἵνα κερδάνω τοὺς ἀνόμους·
\vs{22}~ἐγενόμην τοῖς ἀσθενέσιν ἀσθενής, ἵνα τοὺς ἀσθενεῖς κερδήσω· τοῖς πᾶσιν γέγονα πάντα, ἵνα πάντως τινὰς σώσω.
\vs{23}~πάντα δὲ ποιῶ διὰ τὸ εὐαγγέλιον, ἵνα συγκοινωνὸς αὐτοῦ γένωμαι.
\vs{24}~Οὐκ οἴδατε ὅτι οἱ ἐν σταδίῳ τρέχοντες πάντες μὲν τρέχουσιν, εἷς δὲ λαμβάνει τὸ βραβεῖον; οὕτως τρέχετε ἵνα καταλάβητε.
\vs{25}~πᾶς δὲ ὁ ἀγωνιζόμενος πάντα ἐγκρατεύεται, ἐκεῖνοι μὲν οὖν ἵνα φθαρτὸν στέφανον λάβωσιν, ἡμεῖς δὲ ἄφθαρτον.
\vs{26}~ἐγὼ τοίνυν οὕτως τρέχω ὡς οὐκ ἀδήλως, οὕτως πυκτεύω ὡς οὐκ ἀέρα δέρων·
\vs{27}~ἀλλὰ ὑπωπιάζω μου τὸ σῶμα καὶ δουλαγωγῶ, μή πως ἄλλοις κηρύξας αὐτὸς ἀδόκιμος γένωμαι.

\chapheader{10}
\vs{1}~Οὐ θέλω γὰρ ὑμᾶς ἀγνοεῖν, \adeA{ἀδελφοί}, ὅτι οἱ πατέρες ἡμῶν πάντες ὑπὸ τὴν νεφέλην ἦσαν καὶ πάντες διὰ τῆς θαλάσσης διῆλθον,
\vs{2}~καὶ πάντες εἰς τὸν Μωϋσῆν ἐβαπτίσαντο ἐν τῇ νεφέλῃ καὶ ἐν τῇ θαλάσσῃ,
\vs{3}~καὶ πάντες τὸ αὐτὸ πνευματικὸν βρῶμα ἔφαγον
\vs{4}~καὶ πάντες τὸ αὐτὸ πνευματικὸν ἔπιον πόμα, ἔπινον γὰρ ἐκ πνευματικῆς ἀκολουθούσης πέτρας, ἡ πέτρα δὲ ἦν ὁ Χριστός·
\vs{5}~ἀλλ’ οὐκ ἐν τοῖς πλείοσιν αὐτῶν ηὐδόκησεν ὁ θεός, κατεστρώθησαν γὰρ ἐν τῇ ἐρήμῳ.
\vs{6}~Ταῦτα δὲ τύποι ἡμῶν ἐγενήθησαν, εἰς τὸ μὴ εἶναι ἡμᾶς ἐπιθυμητὰς κακῶν, καθὼς κἀκεῖνοι ἐπεθύμησαν.
\vs{7}~μηδὲ εἰδωλολάτραι γίνεσθε, καθώς τινες αὐτῶν· ὥσπερ γέγραπται· Ἐκάθισεν ὁ λαὸς φαγεῖν καὶ πεῖν, καὶ ἀνέστησαν παίζειν.
\vs{8}~μηδὲ πορνεύωμεν, καθώς τινες αὐτῶν ἐπόρνευσαν, καὶ ἔπεσαν μιᾷ ἡμέρᾳ εἴκοσι τρεῖς χιλιάδες.
\vs{9}~μηδὲ ἐκπειράζωμεν τὸν Χριστόν, καθώς τινες αὐτῶν ἐπείρασαν, καὶ ὑπὸ τῶν ὄφεων ἀπώλλυντο.
\vs{10}~μηδὲ γογγύζετε, καθάπερ τινὲς αὐτῶν ἐγόγγυσαν, καὶ ἀπώλοντο ὑπὸ τοῦ ὀλοθρευτοῦ.
\vs{11}~ταῦτα δὲ τυπικῶς συνέβαινεν ἐκείνοις, ἐγράφη δὲ πρὸς νουθεσίαν ἡμῶν, εἰς οὓς τὰ τέλη τῶν αἰώνων κατήντηκεν.
\vs{12}~ὥστε ὁ δοκῶν ἑστάναι βλεπέτω μὴ πέσῃ,
\vs{13}~πειρασμὸς ὑμᾶς οὐκ εἴληφεν εἰ μὴ ἀνθρώπινος· πιστὸς δὲ ὁ θεός, ὃς οὐκ ἐάσει ὑμᾶς πειρασθῆναι ὑπὲρ ὃ δύνασθε, ἀλλὰ ποιήσει σὺν τῷ πειρασμῷ καὶ τὴν ἔκβασιν τοῦ δύνασθαι ὑπενεγκεῖν.
\vs{14}~Διόπερ, ἀγαπητοί μου, φεύγετε ἀπὸ τῆς εἰδωλολατρίας.
\vs{15}~ὡς φρονίμοις λέγω· κρίνατε ὑμεῖς ὅ φημι.
\vs{16}~τὸ ποτήριον τῆς εὐλογίας ὃ εὐλογοῦμεν, οὐχὶ κοινωνία ἐστὶν τοῦ αἵματος τοῦ Χριστοῦ; τὸν ἄρτον ὃν κλῶμεν, οὐχὶ κοινωνία τοῦ σώματος τοῦ Χριστοῦ ἐστιν;
\vs{17}~ὅτι εἷς ἄρτος, ἓν σῶμα οἱ πολλοί ἐσμεν, οἱ γὰρ πάντες ἐκ τοῦ ἑνὸς ἄρτου μετέχομεν.
\vs{18}~βλέπετε τὸν Ἰσραὴλ κατὰ σάρκα· οὐχ οἱ ἐσθίοντες τὰς θυσίας κοινωνοὶ τοῦ θυσιαστηρίου εἰσίν;
\vs{19}~τί οὖν φημι; ὅτι εἰδωλόθυτόν τί ἐστιν, ἢ ὅτι εἴδωλόν τί ἐστιν;
\vs{20}~ἀλλ’ ὅτι ἃ θύουσιν, δαιμονίοις καὶ οὐ θεῷ θύουσιν, οὐ θέλω δὲ ὑμᾶς κοινωνοὺς τῶν δαιμονίων γίνεσθαι.
\vs{21}~οὐ δύνασθε ποτήριον κυρίου πίνειν καὶ ποτήριον δαιμονίων· οὐ δύνασθε τραπέζης κυρίου μετέχειν καὶ τραπέζης δαιμονίων.
\vs{22}~ἢ παραζηλοῦμεν τὸν κύριον; μὴ ἰσχυρότεροι αὐτοῦ ἐσμεν;
\vs{23}~Πάντα ἔξεστιν· ἀλλ’ οὐ πάντα συμφέρει. πάντα ἔξεστιν· ἀλλ’ οὐ πάντα οἰκοδομεῖ.
\vs{24}~μηδεὶς τὸ ἑαυτοῦ ζητείτω ἀλλὰ τὸ τοῦ ἑτέρου.
\vs{25}~πᾶν τὸ ἐν μακέλλῳ πωλούμενον ἐσθίετε μηδὲν ἀνακρίνοντες διὰ τὴν συνείδησιν,
\vs{26}~τοῦ κυρίου γὰρ ἡ γῆ καὶ τὸ πλήρωμα αὐτῆς.
\vs{27}~εἴ τις καλεῖ ὑμᾶς τῶν ἀπίστων καὶ θέλετε πορεύεσθαι, πᾶν τὸ παρατιθέμενον ὑμῖν ἐσθίετε μηδὲν ἀνακρίνοντες διὰ τὴν συνείδησιν·
\vs{28}~ἐὰν δέ τις ὑμῖν εἴπῃ· Τοῦτο ἱερόθυτόν ἐστιν, μὴ ἐσθίετε δι’ ἐκεῖνον τὸν μηνύσαντα καὶ τὴν συνείδησιν·
\vs{29}~συνείδησιν δὲ λέγω οὐχὶ τὴν ἑαυτοῦ ἀλλὰ τὴν τοῦ ἑτέρου· ἱνατί γὰρ ἡ ἐλευθερία μου κρίνεται ὑπὸ ἄλλης συνειδήσεως;
\vs{30}~εἰ ἐγὼ χάριτι μετέχω, τί βλασφημοῦμαι ὑπὲρ οὗ ἐγὼ εὐχαριστῶ;
\vs{31}~Εἴτε οὖν ἐσθίετε εἴτε πίνετε εἴτε τι ποιεῖτε, πάντα εἰς δόξαν θεοῦ ποιεῖτε.
\vs{32}~ἀπρόσκοποι καὶ Ἰουδαίοις γίνεσθε καὶ Ἕλλησιν καὶ τῇ ἐκκλησίᾳ τοῦ θεοῦ,
\vs{33}~καθὼς κἀγὼ πάντα πᾶσιν ἀρέσκω, μὴ ζητῶν τὸ ἐμαυτοῦ σύμφορον ἀλλὰ τὸ τῶν πολλῶν, ἵνα σωθῶσιν.

\chapheader{11}
\vs{1}~μιμηταί μου γίνεσθε, καθὼς κἀγὼ Χριστοῦ.
\vs{2}~Ἐπαινῶ δὲ ὑμᾶς ὅτι πάντα μου μέμνησθε καὶ καθὼς παρέδωκα ὑμῖν τὰς παραδόσεις κατέχετε.
\vs{3}~θέλω δὲ ὑμᾶς εἰδέναι ὅτι παντὸς ἀνδρὸς ἡ κεφαλὴ ὁ Χριστός ἐστιν, κεφαλὴ δὲ γυναικὸς ὁ ἀνήρ, κεφαλὴ δὲ τοῦ Χριστοῦ ὁ θεός.
\vs{4}~πᾶς ἀνὴρ προσευχόμενος ἢ προφητεύων κατὰ κεφαλῆς ἔχων καταισχύνει τὴν κεφαλὴν αὐτοῦ·
\vs{5}~πᾶσα δὲ γυνὴ προσευχομένη ἢ προφητεύουσα ἀκατακαλύπτῳ τῇ κεφαλῇ καταισχύνει τὴν κεφαλὴν αὐτῆς, ἓν γάρ ἐστιν καὶ τὸ αὐτὸ τῇ ἐξυρημένῃ.
\vs{6}~εἰ γὰρ οὐ κατακαλύπτεται γυνή, καὶ κειράσθω· εἰ δὲ αἰσχρὸν γυναικὶ τὸ κείρασθαι ἢ ξυρᾶσθαι, κατακαλυπτέσθω.
\vs{7}~ἀνὴρ μὲν γὰρ οὐκ ὀφείλει κατακαλύπτεσθαι τὴν κεφαλήν, εἰκὼν καὶ δόξα θεοῦ ὑπάρχων· ἡ γυνὴ δὲ δόξα ἀνδρός ἐστιν.
\vs{8}~οὐ γάρ ἐστιν ἀνὴρ ἐκ γυναικός, ἀλλὰ γυνὴ ἐξ ἀνδρός·
\vs{9}~καὶ γὰρ οὐκ ἐκτίσθη ἀνὴρ διὰ τὴν γυναῖκα, ἀλλὰ γυνὴ διὰ τὸν ἄνδρα.
\vs{10}~διὰ τοῦτο ὀφείλει ἡ γυνὴ ἐξουσίαν ἔχειν ἐπὶ τῆς κεφαλῆς διὰ τοὺς ἀγγέλους.
\vs{11}~πλὴν οὔτε γυνὴ χωρὶς ἀνδρὸς οὔτε ἀνὴρ χωρὶς γυναικὸς ἐν κυρίῳ·
\vs{12}~ὥσπερ γὰρ ἡ γυνὴ ἐκ τοῦ ἀνδρός, οὕτως καὶ ὁ ἀνὴρ διὰ τῆς γυναικός· τὰ δὲ πάντα ἐκ τοῦ θεοῦ.
\vs{13}~ἐν ὑμῖν αὐτοῖς κρίνατε· πρέπον ἐστὶν γυναῖκα ἀκατακάλυπτον τῷ θεῷ προσεύχεσθαι;
\vs{14}~οὐδὲ ἡ φύσις αὐτὴ διδάσκει ὑμᾶς ὅτι ἀνὴρ μὲν ἐὰν κομᾷ, ἀτιμία αὐτῷ ἐστιν,
\vs{15}~γυνὴ δὲ ἐὰν κομᾷ, δόξα αὐτῇ ἐστιν; ὅτι ἡ κόμη ἀντὶ περιβολαίου δέδοται.
\vs{16}~εἰ δέ τις δοκεῖ φιλόνεικος εἶναι, ἡμεῖς τοιαύτην συνήθειαν οὐκ ἔχομεν, οὐδὲ αἱ ἐκκλησίαι τοῦ θεοῦ.
\vs{17}~Τοῦτο δὲ παραγγέλλων οὐκ ἐπαινῶ ὅτι οὐκ εἰς τὸ κρεῖσσον ἀλλὰ εἰς τὸ ἧσσον συνέρχεσθε.
\vs{18}~πρῶτον μὲν γὰρ συνερχομένων ὑμῶν ἐν ἐκκλησίᾳ ἀκούω σχίσματα ἐν ὑμῖν ὑπάρχειν, καὶ μέρος τι πιστεύω.
\vs{19}~δεῖ γὰρ καὶ αἱρέσεις ἐν ὑμῖν εἶναι, ἵνα καὶ οἱ δόκιμοι φανεροὶ γένωνται ἐν ὑμῖν.
\vs{20}~συνερχομένων οὖν ὑμῶν ἐπὶ τὸ αὐτὸ οὐκ ἔστιν κυριακὸν δεῖπνον φαγεῖν,
\vs{21}~ἕκαστος γὰρ τὸ ἴδιον δεῖπνον προλαμβάνει ἐν τῷ φαγεῖν, καὶ ὃς μὲν πεινᾷ, ὃς δὲ μεθύει.
\vs{22}~μὴ γὰρ οἰκίας οὐκ ἔχετε εἰς τὸ ἐσθίειν καὶ πίνειν; ἢ τῆς ἐκκλησίας τοῦ θεοῦ καταφρονεῖτε, καὶ καταισχύνετε τοὺς μὴ ἔχοντας; τί εἴπω ὑμῖν; ἐπαινέσω ὑμᾶς; ἐν τούτῳ οὐκ ἐπαινῶ.
\vs{23}~Ἐγὼ γὰρ παρέλαβον ἀπὸ τοῦ κυρίου, ὃ καὶ παρέδωκα ὑμῖν, ὅτι ὁ κύριος Ἰησοῦς ἐν τῇ νυκτὶ ᾗ παρεδίδετο ἔλαβεν ἄρτον
\vs{24}~καὶ εὐχαριστήσας ἔκλασεν καὶ εἶπεν· Τοῦτό μού ἐστιν τὸ σῶμα τὸ ὑπὲρ ὑμῶν· τοῦτο ποιεῖτε εἰς τὴν ἐμὴν ἀνάμνησιν.
\vs{25}~ὡσαύτως καὶ τὸ ποτήριον μετὰ τὸ δειπνῆσαι, λέγων· Τοῦτο τὸ ποτήριον ἡ καινὴ διαθήκη ἐστὶν ἐν τῷ ἐμῷ αἵματι· τοῦτο ποιεῖτε, ὁσάκις ἐὰν πίνητε, εἰς τὴν ἐμὴν ἀνάμνησιν.
\vs{26}~ὁσάκις γὰρ ἐὰν ἐσθίητε τὸν ἄρτον τοῦτον καὶ τὸ ποτήριον πίνητε, τὸν θάνατον τοῦ κυρίου καταγγέλλετε, ἄχρι οὗ ἔλθῃ.
\vs{27}~Ὥστε ὃς ἂν ἐσθίῃ τὸν ἄρτον ἢ πίνῃ τὸ ποτήριον τοῦ κυρίου ἀναξίως, ἔνοχος ἔσται τοῦ σώματος καὶ τοῦ αἵματος τοῦ κυρίου.
\vs{28}~δοκιμαζέτω δὲ ἄνθρωπος ἑαυτόν, καὶ οὕτως ἐκ τοῦ ἄρτου ἐσθιέτω καὶ ἐκ τοῦ ποτηρίου πινέτω·
\vs{29}~ὁ γὰρ ἐσθίων καὶ πίνων κρίμα ἑαυτῷ ἐσθίει καὶ πίνει μὴ διακρίνων τὸ σῶμα.
\vs{30}~διὰ τοῦτο ἐν ὑμῖν πολλοὶ ἀσθενεῖς καὶ ἄρρωστοι καὶ κοιμῶνται ἱκανοί.
\vs{31}~εἰ δὲ ἑαυτοὺς διεκρίνομεν, οὐκ ἂν ἐκρινόμεθα·
\vs{32}~κρινόμενοι δὲ ὑπὸ κυρίου παιδευόμεθα, ἵνα μὴ σὺν τῷ κόσμῳ κατακριθῶμεν.
\vs{33}~Ὥστε, \adeA{ἀδελφοί} μου, συνερχόμενοι εἰς τὸ φαγεῖν ἀλλήλους ἐκδέχεσθε.
\vs{34}~εἴ τις πεινᾷ, ἐν οἴκῳ ἐσθιέτω, ἵνα μὴ εἰς κρίμα συνέρχησθε. Τὰ δὲ λοιπὰ ὡς ἂν ἔλθω διατάξομαι.

\chapheader{12}
\vs{1}~Περὶ δὲ τῶν πνευματικῶν, \adeA{ἀδελφοί}, οὐ θέλω ὑμᾶς ἀγνοεῖν.
\vs{2}~οἴδατε ὅτι ὅτε ἔθνη ἦτε πρὸς τὰ εἴδωλα τὰ ἄφωνα ὡς ἂν ἤγεσθε ἀπαγόμενοι.
\vs{3}~διὸ γνωρίζω ὑμῖν ὅτι οὐδεὶς ἐν πνεύματι θεοῦ λαλῶν λέγει· Ἀνάθεμα Ἰησοῦς, καὶ οὐδεὶς δύναται εἰπεῖν· Κύριος Ἰησοῦς εἰ μὴ ἐν πνεύματι ἁγίῳ.
\vs{4}~Διαιρέσεις δὲ χαρισμάτων εἰσίν, τὸ δὲ αὐτὸ πνεῦμα·
\vs{5}~καὶ διαιρέσεις διακονιῶν εἰσιν, καὶ ὁ αὐτὸς κύριος·
\vs{6}~καὶ διαιρέσεις ἐνεργημάτων εἰσίν, ὁ δὲ αὐτὸς θεός, ὁ ἐνεργῶν τὰ πάντα ἐν πᾶσιν.
\vs{7}~ἑκάστῳ δὲ δίδοται ἡ φανέρωσις τοῦ πνεύματος πρὸς τὸ συμφέρον.
\vs{8}~ᾧ μὲν γὰρ διὰ τοῦ πνεύματος δίδοται λόγος σοφίας, ἄλλῳ δὲ λόγος γνώσεως κατὰ τὸ αὐτὸ πνεῦμα,
\vs{9}~ἑτέρῳ πίστις ἐν τῷ αὐτῷ πνεύματι, ἄλλῳ χαρίσματα ἰαμάτων ἐν τῷ ἑνὶ πνεύματι,
\vs{10}~ἄλλῳ ἐνεργήματα δυνάμεων, ἄλλῳ προφητεία, 1ἄλλῳ διακρίσεις πνευμάτων, ἑτέρῳ γένη γλωσσῶν, 2ἄλλῳ ἑρμηνεία γλωσσῶν·
\vs{11}~πάντα δὲ ταῦτα ἐνεργεῖ τὸ ἓν καὶ τὸ αὐτὸ πνεῦμα, διαιροῦν ἰδίᾳ ἑκάστῳ καθὼς βούλεται.
\vs{12}~Καθάπερ γὰρ τὸ σῶμα ἕν ἐστιν καὶ μέλη πολλὰ ἔχει, πάντα δὲ τὰ μέλη τοῦ σώματος πολλὰ ὄντα ἕν ἐστιν σῶμα, οὕτως καὶ ὁ Χριστός·
\vs{13}~καὶ γὰρ ἐν ἑνὶ πνεύματι ἡμεῖς πάντες εἰς ἓν σῶμα ἐβαπτίσθημεν, εἴτε Ἰουδαῖοι εἴτε Ἕλληνες, εἴτε δοῦλοι εἴτε ἐλεύθεροι, καὶ πάντες ἓν πνεῦμα ἐποτίσθημεν.
\vs{14}~Καὶ γὰρ τὸ σῶμα οὐκ ἔστιν ἓν μέλος ἀλλὰ πολλά.
\vs{15}~ἐὰν εἴπῃ ὁ πούς· Ὅτι οὐκ εἰμὶ χείρ, οὐκ εἰμὶ ἐκ τοῦ σώματος, οὐ παρὰ τοῦτο οὐκ ἔστιν ἐκ τοῦ σώματος;
\vs{16}~καὶ ἐὰν εἴπῃ τὸ οὖς· Ὅτι οὐκ εἰμὶ ὀφθαλμός, οὐκ εἰμὶ ἐκ τοῦ σώματος, οὐ παρὰ τοῦτο οὐκ ἔστιν ἐκ τοῦ σώματος·
\vs{17}~εἰ ὅλον τὸ σῶμα ὀφθαλμός, ποῦ ἡ ἀκοή; εἰ ὅλον ἀκοή, ποῦ ἡ ὄσφρησις;
\vs{18}~νυνὶ δὲ ὁ θεὸς ἔθετο τὰ μέλη, ἓν ἕκαστον αὐτῶν, ἐν τῷ σώματι καθὼς ἠθέλησεν.
\vs{19}~εἰ δὲ ἦν τὰ πάντα ἓν μέλος, ποῦ τὸ σῶμα;
\vs{20}~νῦν δὲ πολλὰ μὲν μέλη, ἓν δὲ σῶμα.
\vs{21}~οὐ δύναται δὲ ὁ ὀφθαλμὸς εἰπεῖν τῇ χειρί· Χρείαν σου οὐκ ἔχω, ἢ πάλιν ἡ κεφαλὴ τοῖς ποσίν· Χρείαν ὑμῶν οὐκ ἔχω·
\vs{22}~ἀλλὰ πολλῷ μᾶλλον τὰ δοκοῦντα μέλη τοῦ σώματος ἀσθενέστερα ὑπάρχειν ἀναγκαῖά ἐστιν,
\vs{23}~καὶ ἃ δοκοῦμεν ἀτιμότερα εἶναι τοῦ σώματος, τούτοις τιμὴν περισσοτέραν περιτίθεμεν, καὶ τὰ ἀσχήμονα ἡμῶν εὐσχημοσύνην περισσοτέραν ἔχει,
\vs{24}~τὰ δὲ εὐσχήμονα ἡμῶν οὐ χρείαν ἔχει. ἀλλὰ ὁ θεὸς συνεκέρασεν τὸ σῶμα, τῷ ὑστεροῦντι περισσοτέραν δοὺς τιμήν,
\vs{25}~ἵνα μὴ ᾖ σχίσμα ἐν τῷ σώματι, ἀλλὰ τὸ αὐτὸ ὑπὲρ ἀλλήλων μεριμνῶσι τὰ μέλη.
\vs{26}~καὶ εἴτε πάσχει ἓν μέλος, συμπάσχει πάντα τὰ μέλη· εἴτε δοξάζεται μέλος, συγχαίρει πάντα τὰ μέλη.
\vs{27}~Ὑμεῖς δέ ἐστε σῶμα Χριστοῦ καὶ μέλη ἐκ μέρους.
\vs{28}~καὶ οὓς μὲν ἔθετο ὁ θεὸς ἐν τῇ ἐκκλησίᾳ πρῶτον ἀποστόλους, δεύτερον προφήτας, τρίτον διδασκάλους, ἔπειτα δυνάμεις, ἔπειτα χαρίσματα ἰαμάτων, ἀντιλήμψεις, κυβερνήσεις, γένη γλωσσῶν.
\vs{29}~μὴ πάντες ἀπόστολοι; μὴ πάντες προφῆται; μὴ πάντες διδάσκαλοι; μὴ πάντες δυνάμεις;
\vs{30}~μὴ πάντες χαρίσματα ἔχουσιν ἰαμάτων; μὴ πάντες γλώσσαις λαλοῦσιν; μὴ πάντες διερμηνεύουσιν;
\vs{31}~ζηλοῦτε δὲ τὰ χαρίσματα τὰ μείζονα. καὶ ἔτι καθ’ ὑπερβολὴν ὁδὸν ὑμῖν δείκνυμι.

\chapheader{13}
\vs{1}~Ἐὰν ταῖς γλώσσαις τῶν ἀνθρώπων λαλῶ καὶ τῶν ἀγγέλων, ἀγάπην δὲ μὴ ἔχω, γέγονα χαλκὸς ἠχῶν ἢ κύμβαλον ἀλαλάζον.
\vs{2}~καὶ ἐὰν ἔχω προφητείαν καὶ εἰδῶ τὰ μυστήρια πάντα καὶ πᾶσαν τὴν γνῶσιν, καὶ ἐὰν ἔχω πᾶσαν τὴν πίστιν ὥστε ὄρη μεθιστάναι, ἀγάπην δὲ μὴ ἔχω, οὐθέν εἰμι.
\vs{3}~καὶ ἐὰν ψωμίσω πάντα τὰ ὑπάρχοντά μου, καὶ ἐὰν παραδῶ τὸ σῶμά μου, ἵνα καυθήσομαι, ἀγάπην δὲ μὴ ἔχω, οὐδὲν ὠφελοῦμαι.
\vs{4}~Ἡ ἀγάπη μακροθυμεῖ, χρηστεύεται ἡ ἀγάπη, οὐ ζηλοῖ ἡ ἀγάπη, οὐ περπερεύεται, οὐ φυσιοῦται,
\vs{5}~οὐκ ἀσχημονεῖ, οὐ ζητεῖ τὰ ἑαυτῆς, οὐ παροξύνεται, οὐ λογίζεται τὸ κακόν,
\vs{6}~οὐ χαίρει ἐπὶ τῇ ἀδικίᾳ, συγχαίρει δὲ τῇ ἀληθείᾳ·
\vs{7}~πάντα στέγει, πάντα πιστεύει, πάντα ἐλπίζει, πάντα ὑπομένει.
\vs{8}~Ἡ ἀγάπη οὐδέποτε πίπτει. εἴτε δὲ προφητεῖαι, καταργηθήσονται· εἴτε γλῶσσαι, παύσονται· εἴτε γνῶσις, καταργηθήσεται.
\vs{9}~ἐκ μέρους γὰρ γινώσκομεν καὶ ἐκ μέρους προφητεύομεν·
\vs{10}~ὅταν δὲ ἔλθῃ τὸ τέλειον, τὸ ἐκ μέρους καταργηθήσεται.
\vs{11}~ὅτε ἤμην νήπιος, ἐλάλουν ὡς νήπιος, ἐφρόνουν ὡς νήπιος, ἐλογιζόμην ὡς νήπιος· ὅτε γέγονα ἀνήρ, κατήργηκα τὰ τοῦ νηπίου.
\vs{12}~βλέπομεν γὰρ ἄρτι δι’ ἐσόπτρου ἐν αἰνίγματι, τότε δὲ πρόσωπον πρὸς πρόσωπον· ἄρτι γινώσκω ἐκ μέρους, τότε δὲ ἐπιγνώσομαι καθὼς καὶ ἐπεγνώσθην.
\vs{13}~νυνὶ δὲ μένει πίστις, ἐλπίς, ἀγάπη· τὰ τρία ταῦτα, μείζων δὲ τούτων ἡ ἀγάπη.

\chapheader{14}
\vs{1}~Διώκετε τὴν ἀγάπην, ζηλοῦτε δὲ τὰ πνευματικά, μᾶλλον δὲ ἵνα προφητεύητε.
\vs{2}~ὁ γὰρ λαλῶν γλώσσῃ οὐκ ἀνθρώποις λαλεῖ ἀλλὰ θεῷ, οὐδεὶς γὰρ ἀκούει, πνεύματι δὲ λαλεῖ μυστήρια·
\vs{3}~ὁ δὲ προφητεύων ἀνθρώποις λαλεῖ οἰκοδομὴν καὶ παράκλησιν καὶ παραμυθίαν.
\vs{4}~ὁ λαλῶν γλώσσῃ ἑαυτὸν οἰκοδομεῖ· ὁ δὲ προφητεύων ἐκκλησίαν οἰκοδομεῖ.
\vs{5}~θέλω δὲ πάντας ὑμᾶς λαλεῖν γλώσσαις, μᾶλλον δὲ ἵνα προφητεύητε· μείζων δὲ ὁ προφητεύων ἢ ὁ λαλῶν γλώσσαις, ἐκτὸς εἰ μὴ διερμηνεύῃ, ἵνα ἡ ἐκκλησία οἰκοδομὴν λάβῃ.
\vs{6}~Νῦν δέ, \adeA{ἀδελφοί}, ἐὰν ἔλθω πρὸς ὑμᾶς γλώσσαις λαλῶν, τί ὑμᾶς ὠφελήσω, ἐὰν μὴ ὑμῖν λαλήσω ἢ ἐν ἀποκαλύψει ἢ ἐν γνώσει ἢ ἐν προφητείᾳ ἢ ἐν διδαχῇ;
\vs{7}~ὅμως τὰ ἄψυχα φωνὴν διδόντα, εἴτε αὐλὸς εἴτε κιθάρα, ἐὰν διαστολὴν τοῖς φθόγγοις μὴ δῷ, πῶς γνωσθήσεται τὸ αὐλούμενον ἢ τὸ κιθαριζόμενον;
\vs{8}~καὶ γὰρ ἐὰν ἄδηλον φωνὴν σάλπιγξ δῷ, τίς παρασκευάσεται εἰς πόλεμον;
\vs{9}~οὕτως καὶ ὑμεῖς διὰ τῆς γλώσσης ἐὰν μὴ εὔσημον λόγον δῶτε, πῶς γνωσθήσεται τὸ λαλούμενον; ἔσεσθε γὰρ εἰς ἀέρα λαλοῦντες.
\vs{10}~τοσαῦτα εἰ τύχοι γένη φωνῶν εἰσιν ἐν κόσμῳ, καὶ οὐδὲν ἄφωνον·
\vs{11}~ἐὰν οὖν μὴ εἰδῶ τὴν δύναμιν τῆς φωνῆς, ἔσομαι τῷ λαλοῦντι βάρβαρος καὶ ὁ λαλῶν ἐν ἐμοὶ βάρβαρος.
\vs{12}~οὕτως καὶ ὑμεῖς, ἐπεὶ ζηλωταί ἐστε πνευμάτων, πρὸς τὴν οἰκοδομὴν τῆς ἐκκλησίας ζητεῖτε ἵνα περισσεύητε.
\vs{13}~Διὸ ὁ λαλῶν γλώσσῃ προσευχέσθω ἵνα διερμηνεύῃ.
\vs{14}~ἐὰν γὰρ προσεύχωμαι γλώσσῃ, τὸ πνεῦμά μου προσεύχεται, ὁ δὲ νοῦς μου ἄκαρπός ἐστιν.
\vs{15}~τί οὖν ἐστιν; προσεύξομαι τῷ πνεύματι, προσεύξομαι δὲ καὶ τῷ νοΐ· ψαλῶ τῷ πνεύματι, ψαλῶ δὲ καὶ τῷ νοΐ·
\vs{16}~ἐπεὶ ἐὰν εὐλογῇς πνεύματι, ὁ ἀναπληρῶν τὸν τόπον τοῦ ἰδιώτου πῶς ἐρεῖ τὸ Ἀμήν ἐπὶ τῇ σῇ εὐχαριστίᾳ; ἐπειδὴ τί λέγεις οὐκ οἶδεν·
\vs{17}~σὺ μὲν γὰρ καλῶς εὐχαριστεῖς, ἀλλ’ ὁ ἕτερος οὐκ οἰκοδομεῖται.
\vs{18}~εὐχαριστῶ τῷ θεῷ, πάντων ὑμῶν μᾶλλον γλώσσαις λαλῶ·
\vs{19}~ἀλλὰ ἐν ἐκκλησίᾳ θέλω πέντε λόγους τῷ νοΐ μου λαλῆσαι, ἵνα καὶ ἄλλους κατηχήσω, ἢ μυρίους λόγους ἐν γλώσσῃ.
\vs{20}~\adeA{Ἀδελφοί}, μὴ παιδία γίνεσθε ταῖς φρεσίν, ἀλλὰ τῇ κακίᾳ νηπιάζετε, ταῖς δὲ φρεσὶν τέλειοι γίνεσθε.
\vs{21}~ἐν τῷ νόμῳ γέγραπται ὅτι Ἐν ἑτερογλώσσοις καὶ ἐν χείλεσιν ἑτέρων λαλήσω τῷ λαῷ τούτῳ, καὶ οὐδ’ οὕτως εἰσακούσονταί μου, λέγει κύριος.
\vs{22}~ὥστε αἱ γλῶσσαι εἰς σημεῖόν εἰσιν οὐ τοῖς πιστεύουσιν ἀλλὰ τοῖς ἀπίστοις, ἡ δὲ προφητεία οὐ τοῖς ἀπίστοις ἀλλὰ τοῖς πιστεύουσιν.
\vs{23}~ἐὰν οὖν συνέλθῃ ἡ ἐκκλησία ὅλη ἐπὶ τὸ αὐτὸ καὶ πάντες λαλῶσιν γλώσσαις, εἰσέλθωσιν δὲ ἰδιῶται ἢ ἄπιστοι, οὐκ ἐροῦσιν ὅτι μαίνεσθε;
\vs{24}~ἐὰν δὲ πάντες προφητεύωσιν, εἰσέλθῃ δέ τις ἄπιστος ἢ ἰδιώτης, ἐλέγχεται ὑπὸ πάντων, ἀνακρίνεται ὑπὸ πάντων,
\vs{25}~τὰ κρυπτὰ τῆς καρδίας αὐτοῦ φανερὰ γίνεται, καὶ οὕτως πεσὼν ἐπὶ πρόσωπον προσκυνήσει τῷ θεῷ, ἀπαγγέλλων ὅτι Ὄντως ὁ θεὸς ἐν ὑμῖν ἐστιν.
\vs{26}~Τί οὖν ἐστιν, \adeA{ἀδελφοί}; ὅταν συνέρχησθε, ἕκαστος ψαλμὸν ἔχει, διδαχὴν ἔχει, ἀποκάλυψιν ἔχει, γλῶσσαν ἔχει, ἑρμηνείαν ἔχει· πάντα πρὸς οἰκοδομὴν γινέσθω.
\vs{27}~εἴτε γλώσσῃ τις λαλεῖ, κατὰ δύο ἢ τὸ πλεῖστον τρεῖς, καὶ ἀνὰ μέρος, καὶ εἷς διερμηνευέτω·
\vs{28}~ἐὰν δὲ μὴ ᾖ διερμηνευτής, σιγάτω ἐν ἐκκλησίᾳ, ἑαυτῷ δὲ λαλείτω καὶ τῷ θεῷ.
\vs{29}~προφῆται δὲ δύο ἢ τρεῖς λαλείτωσαν, καὶ οἱ ἄλλοι διακρινέτωσαν·
\vs{30}~ἐὰν δὲ ἄλλῳ ἀποκαλυφθῇ καθημένῳ, ὁ πρῶτος σιγάτω.
\vs{31}~δύνασθε γὰρ καθ’ ἕνα πάντες προφητεύειν, ἵνα πάντες μανθάνωσιν καὶ πάντες παρακαλῶνται
\vs{32}~(καὶ πνεύματα προφητῶν προφήταις ὑποτάσσεται,
\vs{33}~οὐ γάρ ἐστιν ἀκαταστασίας ὁ θεὸς ἀλλὰ εἰρήνης), ὡς ἐν πάσαις ταῖς ἐκκλησίαις τῶν ἁγίων.
\vs{34}~Αἱ γυναῖκες ἐν ταῖς ἐκκλησίαις σιγάτωσαν, οὐ γὰρ ἐπιτρέπεται αὐταῖς λαλεῖν· ἀλλὰ ὑποτασσέσθωσαν, καθὼς καὶ ὁ νόμος λέγει.
\vs{35}~εἰ δέ τι μαθεῖν θέλουσιν, ἐν οἴκῳ τοὺς ἰδίους ἄνδρας ἐπερωτάτωσαν, αἰσχρὸν γάρ ἐστιν γυναικὶ λαλεῖν ἐν ἐκκλησίᾳ.
\vs{36}~ἢ ἀφ’ ὑμῶν ὁ λόγος τοῦ θεοῦ ἐξῆλθεν, ἢ εἰς ὑμᾶς μόνους κατήντησεν;
\vs{37}~Εἴ τις δοκεῖ προφήτης εἶναι ἢ πνευματικός, ἐπιγινωσκέτω ἃ γράφω ὑμῖν ὅτι κυρίου ἐστίν·
\vs{38}~εἰ δέ τις ἀγνοεῖ, ἀγνοεῖται.
\vs{39}~ὥστε, \adeA{ἀδελφοί} μου, ζηλοῦτε τὸ προφητεύειν, καὶ τὸ λαλεῖν μὴ κωλύετε γλώσσαις·
\vs{40}~πάντα δὲ εὐσχημόνως καὶ κατὰ τάξιν γινέσθω.

\chapheader{15}
\vs{1}~Γνωρίζω δὲ ὑμῖν, \adeA{ἀδελφοί}, τὸ εὐαγγέλιον ὃ εὐηγγελισάμην ὑμῖν, ὃ καὶ παρελάβετε, ἐν ᾧ καὶ ἑστήκατε,
\vs{2}~δι’ οὗ καὶ σῴζεσθε, τίνι λόγῳ εὐηγγελισάμην ὑμῖν, εἰ κατέχετε, ἐκτὸς εἰ μὴ εἰκῇ ἐπιστεύσατε.
\vs{3}~Παρέδωκα γὰρ ὑμῖν ἐν πρώτοις, ὃ καὶ παρέλαβον, ὅτι Χριστὸς ἀπέθανεν ὑπὲρ τῶν ἁμαρτιῶν ἡμῶν κατὰ τὰς γραφάς,
\vs{4}~καὶ ὅτι ἐτάφη, καὶ ὅτι ἐγήγερται τῇ ἡμέρᾳ τῇ τρίτῃ κατὰ τὰς γραφάς,
\vs{5}~καὶ ὅτι ὤφθη Κηφᾷ, εἶτα τοῖς δώδεκα·
\vs{6}~ἔπειτα ὤφθη ἐπάνω πεντακοσίοις \adeA{ἀδελφοῖς} ἐφάπαξ, ἐξ ὧν οἱ πλείονες μένουσιν ἕως ἄρτι, τινὲς δὲ ἐκοιμήθησαν·
\vs{7}~ἔπειτα ὤφθη Ἰακώβῳ, εἶτα τοῖς ἀποστόλοις πᾶσιν·
\vs{8}~ἔσχατον δὲ πάντων ὡσπερεὶ τῷ ἐκτρώματι ὤφθη κἀμοί.
\vs{9}~ἐγὼ γάρ εἰμι ὁ ἐλάχιστος τῶν ἀποστόλων, ὃς οὐκ εἰμὶ ἱκανὸς καλεῖσθαι ἀπόστολος, διότι ἐδίωξα τὴν ἐκκλησίαν τοῦ θεοῦ·
\vs{10}~χάριτι δὲ θεοῦ εἰμι ὅ εἰμι, καὶ ἡ χάρις αὐτοῦ ἡ εἰς ἐμὲ οὐ κενὴ ἐγενήθη, ἀλλὰ περισσότερον αὐτῶν πάντων ἐκοπίασα, οὐκ ἐγὼ δὲ ἀλλὰ ἡ χάρις τοῦ θεοῦ ἡ σὺν ἐμοί.
\vs{11}~εἴτε οὖν ἐγὼ εἴτε ἐκεῖνοι, οὕτως κηρύσσομεν καὶ οὕτως ἐπιστεύσατε.
\vs{12}~Εἰ δὲ Χριστὸς κηρύσσεται ὅτι ἐκ νεκρῶν ἐγήγερται, πῶς λέγουσιν ἐν ὑμῖν τινες ὅτι ἀνάστασις νεκρῶν οὐκ ἔστιν;
\vs{13}~εἰ δὲ ἀνάστασις νεκρῶν οὐκ ἔστιν, οὐδὲ Χριστὸς ἐγήγερται·
\vs{14}~εἰ δὲ Χριστὸς οὐκ ἐγήγερται, κενὸν ἄρα τὸ κήρυγμα ἡμῶν, κενὴ καὶ ἡ πίστις ὑμῶν,
\vs{15}~εὑρισκόμεθα δὲ καὶ ψευδομάρτυρες τοῦ θεοῦ, ὅτι ἐμαρτυρήσαμεν κατὰ τοῦ θεοῦ ὅτι ἤγειρεν τὸν Χριστόν, ὃν οὐκ ἤγειρεν εἴπερ ἄρα νεκροὶ οὐκ ἐγείρονται.
\vs{16}~εἰ γὰρ νεκροὶ οὐκ ἐγείρονται, οὐδὲ Χριστὸς ἐγήγερται·
\vs{17}~εἰ δὲ Χριστὸς οὐκ ἐγήγερται, ματαία ἡ πίστις ὑμῶν, ἔτι ἐστὲ ἐν ταῖς ἁμαρτίαις ὑμῶν.
\vs{18}~ἄρα καὶ οἱ κοιμηθέντες ἐν Χριστῷ ἀπώλοντο.
\vs{19}~εἰ ἐν τῇ ζωῇ ταύτῃ ἐν Χριστῷ ἠλπικότες ἐσμὲν μόνον, ἐλεεινότεροι πάντων ἀνθρώπων ἐσμέν.
\vs{20}~Νυνὶ δὲ Χριστὸς ἐγήγερται ἐκ νεκρῶν, ἀπαρχὴ τῶν κεκοιμημένων.
\vs{21}~ἐπειδὴ γὰρ δι’ ἀνθρώπου θάνατος, καὶ δι’ ἀνθρώπου ἀνάστασις νεκρῶν·
\vs{22}~ὥσπερ γὰρ ἐν τῷ Ἀδὰμ πάντες ἀποθνῄσκουσιν, οὕτως καὶ ἐν τῷ Χριστῷ πάντες ζῳοποιηθήσονται.
\vs{23}~ἕκαστος δὲ ἐν τῷ ἰδίῳ τάγματι· ἀπαρχὴ Χριστός, ἔπειτα οἱ τοῦ Χριστοῦ ἐν τῇ παρουσίᾳ αὐτοῦ·
\vs{24}~εἶτα τὸ τέλος, ὅταν παραδιδῷ τὴν βασιλείαν τῷ θεῷ καὶ πατρί, ὅταν καταργήσῃ πᾶσαν ἀρχὴν καὶ πᾶσαν ἐξουσίαν καὶ δύναμιν,
\vs{25}~δεῖ γὰρ αὐτὸν βασιλεύειν ἄχρι οὗ θῇ πάντας τοὺς ἐχθροὺς ὑπὸ τοὺς πόδας αὐτοῦ.
\vs{26}~ἔσχατος ἐχθρὸς καταργεῖται ὁ θάνατος,
\vs{27}~πάντα γὰρ ὑπέταξεν ὑπὸ τοὺς πόδας αὐτοῦ. ὅταν δὲ εἴπῃ ὅτι πάντα ὑποτέτακται, δῆλον ὅτι ἐκτὸς τοῦ ὑποτάξαντος αὐτῷ τὰ πάντα.
\vs{28}~ὅταν δὲ ὑποταγῇ αὐτῷ τὰ πάντα, τότε αὐτὸς ὁ υἱὸς ὑποταγήσεται τῷ ὑποτάξαντι αὐτῷ τὰ πάντα, ἵνα ᾖ ὁ θεὸς πάντα ἐν πᾶσιν.
\vs{29}~Ἐπεὶ τί ποιήσουσιν οἱ βαπτιζόμενοι ὑπὲρ τῶν νεκρῶν; εἰ ὅλως νεκροὶ οὐκ ἐγείρονται, τί καὶ βαπτίζονται ὑπὲρ αὐτῶν;
\vs{30}~τί καὶ ἡμεῖς κινδυνεύομεν πᾶσαν ὥραν;
\vs{31}~καθ’ ἡμέραν ἀποθνῄσκω, νὴ τὴν ὑμετέραν καύχησιν, ἣν ἔχω ἐν Χριστῷ Ἰησοῦ τῷ κυρίῳ ἡμῶν.
\vs{32}~εἰ κατὰ ἄνθρωπον ἐθηριομάχησα ἐν Ἐφέσῳ, τί μοι τὸ ὄφελος; εἰ νεκροὶ οὐκ ἐγείρονται, Φάγωμεν καὶ πίωμεν, αὔριον γὰρ ἀποθνῄσκομεν.
\vs{33}~μὴ πλανᾶσθε· φθείρουσιν ἤθη χρηστὰ ὁμιλίαι κακαί.
\vs{34}~ἐκνήψατε δικαίως καὶ μὴ ἁμαρτάνετε, ἀγνωσίαν γὰρ θεοῦ τινες ἔχουσιν· πρὸς ἐντροπὴν ὑμῖν λαλῶ.
\vs{35}~Ἀλλὰ ἐρεῖ τις· Πῶς ἐγείρονται οἱ νεκροί, ποίῳ δὲ σώματι ἔρχονται;
\vs{36}~ἄφρων, σὺ ὃ σπείρεις, οὐ ζῳοποιεῖται ἐὰν μὴ ἀποθάνῃ·
\vs{37}~καὶ ὃ σπείρεις, οὐ τὸ σῶμα τὸ γενησόμενον σπείρεις ἀλλὰ γυμνὸν κόκκον εἰ τύχοι σίτου ἤ τινος τῶν λοιπῶν·
\vs{38}~ὁ δὲ θεὸς δίδωσιν αὐτῷ σῶμα καθὼς ἠθέλησεν, καὶ ἑκάστῳ τῶν σπερμάτων ἴδιον σῶμα.
\vs{39}~οὐ πᾶσα σὰρξ ἡ αὐτὴ σάρξ, ἀλλὰ ἄλλη μὲν ἀνθρώπων, ἄλλη δὲ σὰρξ κτηνῶν, ἄλλη δὲ σὰρξ πτηνῶν, ἄλλη δὲ ἰχθύων.
\vs{40}~καὶ σώματα ἐπουράνια, καὶ σώματα ἐπίγεια· ἀλλὰ ἑτέρα μὲν ἡ τῶν ἐπουρανίων δόξα, ἑτέρα δὲ ἡ τῶν ἐπιγείων.
\vs{41}~ἄλλη δόξα ἡλίου, καὶ ἄλλη δόξα σελήνης, καὶ ἄλλη δόξα ἀστέρων, ἀστὴρ γὰρ ἀστέρος διαφέρει ἐν δόξῃ.
\vs{42}~Οὕτως καὶ ἡ ἀνάστασις τῶν νεκρῶν. σπείρεται ἐν φθορᾷ, ἐγείρεται ἐν ἀφθαρσίᾳ·
\vs{43}~σπείρεται ἐν ἀτιμίᾳ, ἐγείρεται ἐν δόξῃ· σπείρεται ἐν ἀσθενείᾳ, ἐγείρεται ἐν δυνάμει·
\vs{44}~σπείρεται σῶμα ψυχικόν, ἐγείρεται σῶμα πνευματικόν. Εἰ ἔστιν σῶμα ψυχικόν, ἔστιν καὶ πνευματικόν.
\vs{45}~οὕτως καὶ γέγραπται· Ἐγένετο ὁ πρῶτος ἄνθρωπος Ἀδὰμ εἰς ψυχὴν ζῶσαν· ὁ ἔσχατος Ἀδὰμ εἰς πνεῦμα ζῳοποιοῦν.
\vs{46}~ἀλλ’ οὐ πρῶτον τὸ πνευματικὸν ἀλλὰ τὸ ψυχικόν, ἔπειτα τὸ πνευματικόν.
\vs{47}~ὁ πρῶτος ἄνθρωπος ἐκ γῆς χοϊκός, ὁ δεύτερος ἄνθρωπος ἐξ οὐρανοῦ.
\vs{48}~οἷος ὁ χοϊκός, τοιοῦτοι καὶ οἱ χοϊκοί, καὶ οἷος ὁ ἐπουράνιος, τοιοῦτοι καὶ οἱ ἐπουράνιοι·
\vs{49}~καὶ καθὼς ἐφορέσαμεν τὴν εἰκόνα τοῦ χοϊκοῦ, φορέσομεν καὶ τὴν εἰκόνα τοῦ ἐπουρανίου.
\vs{50}~Τοῦτο δέ φημι, \adeA{ἀδελφοί}, ὅτι σὰρξ καὶ αἷμα βασιλείαν θεοῦ κληρονομῆσαι οὐ δύναται, οὐδὲ ἡ φθορὰ τὴν ἀφθαρσίαν κληρονομεῖ.
\vs{51}~ἰδοὺ μυστήριον ὑμῖν λέγω· πάντες οὐ κοιμηθησόμεθα πάντες δὲ ἀλλαγησόμεθα,
\vs{52}~ἐν ἀτόμῳ, ἐν ῥιπῇ ὀφθαλμοῦ, ἐν τῇ ἐσχάτῃ σάλπιγγι· σαλπίσει γάρ, καὶ οἱ νεκροὶ ἐγερθήσονται ἄφθαρτοι, καὶ ἡμεῖς ἀλλαγησόμεθα.
\vs{53}~δεῖ γὰρ τὸ φθαρτὸν τοῦτο ἐνδύσασθαι ἀφθαρσίαν καὶ τὸ θνητὸν τοῦτο ἐνδύσασθαι ἀθανασίαν.
\vs{54}~ὅταν δὲ τὸ φθαρτὸν τοῦτο ἐνδύσηται ἀφθαρσίαν καὶ τὸ θνητὸν τοῦτο ἐνδύσηται ἀθανασίαν, τότε γενήσεται ὁ λόγος ὁ γεγραμμένος· Κατεπόθη ὁ θάνατος εἰς νῖκος.
\vs{55}~ποῦ σου, θάνατε, τὸ νῖκος; ποῦ σου, θάνατε, τὸ κέντρον;
\vs{56}~τὸ δὲ κέντρον τοῦ θανάτου ἡ ἁμαρτία, ἡ δὲ δύναμις τῆς ἁμαρτίας ὁ νόμος·
\vs{57}~τῷ δὲ θεῷ χάρις τῷ διδόντι ἡμῖν τὸ νῖκος διὰ τοῦ κυρίου ἡμῶν Ἰησοῦ Χριστοῦ.
\vs{58}~Ὥστε, \adeA{ἀδελφοί} μου ἀγαπητοί, ἑδραῖοι γίνεσθε, ἀμετακίνητοι, περισσεύοντες ἐν τῷ ἔργῳ τοῦ κυρίου πάντοτε, εἰδότες ὅτι ὁ κόπος ὑμῶν οὐκ ἔστιν κενὸς ἐν κυρίῳ.

\chapheader{16}
\vs{1}~Περὶ δὲ τῆς λογείας τῆς εἰς τοὺς ἁγίους, ὥσπερ διέταξα ταῖς ἐκκλησίαις τῆς Γαλατίας, οὕτως καὶ ὑμεῖς ποιήσατε.
\vs{2}~κατὰ μίαν σαββάτου ἕκαστος ὑμῶν παρ’ ἑαυτῷ τιθέτω θησαυρίζων ὅ τι ἐὰν εὐοδῶται, ἵνα μὴ ὅταν ἔλθω τότε λογεῖαι γίνωνται.
\vs{3}~ὅταν δὲ παραγένωμαι, οὓς ἐὰν δοκιμάσητε δι’ ἐπιστολῶν, τούτους πέμψω ἀπενεγκεῖν τὴν χάριν ὑμῶν εἰς Ἰερουσαλήμ·
\vs{4}~ἐὰν δὲ ἄξιον ᾖ τοῦ κἀμὲ πορεύεσθαι, σὺν ἐμοὶ πορεύσονται.
\vs{5}~Ἐλεύσομαι δὲ πρὸς ὑμᾶς ὅταν Μακεδονίαν διέλθω, Μακεδονίαν γὰρ διέρχομαι,
\vs{6}~πρὸς ὑμᾶς δὲ τυχὸν παραμενῶ ἢ καὶ παραχειμάσω, ἵνα ὑμεῖς με προπέμψητε οὗ ἐὰν πορεύωμαι.
\vs{7}~οὐ θέλω γὰρ ὑμᾶς ἄρτι ἐν παρόδῳ ἰδεῖν, ἐλπίζω γὰρ χρόνον τινὰ ἐπιμεῖναι πρὸς ὑμᾶς, ἐὰν ὁ κύριος ἐπιτρέψῃ.
\vs{8}~ἐπιμενῶ δὲ ἐν Ἐφέσῳ ἕως τῆς πεντηκοστῆς·
\vs{9}~θύρα γάρ μοι ἀνέῳγεν μεγάλη καὶ ἐνεργής, καὶ ἀντικείμενοι πολλοί.
\vs{10}~Ἐὰν δὲ ἔλθῃ Τιμόθεος, βλέπετε ἵνα ἀφόβως γένηται πρὸς ὑμᾶς, τὸ γὰρ ἔργον κυρίου ἐργάζεται ὡς κἀγώ·
\vs{11}~μή τις οὖν αὐτὸν ἐξουθενήσῃ. προπέμψατε δὲ αὐτὸν ἐν εἰρήνῃ, ἵνα ἔλθῃ πρός με, ἐκδέχομαι γὰρ αὐτὸν μετὰ τῶν \adeA{ἀδελφῶν}.
\vs{12}~Περὶ δὲ Ἀπολλῶ τοῦ \adeA{ἀδελφοῦ}, πολλὰ παρεκάλεσα αὐτὸν ἵνα ἔλθῃ πρὸς ὑμᾶς μετὰ τῶν \adeA{ἀδελφῶν}· καὶ πάντως οὐκ ἦν θέλημα ἵνα νῦν ἔλθῃ, ἐλεύσεται δὲ ὅταν εὐκαιρήσῃ.
\vs{13}~Γρηγορεῖτε, στήκετε ἐν τῇ πίστει, ἀνδρίζεσθε, κραταιοῦσθε.
\vs{14}~πάντα ὑμῶν ἐν ἀγάπῃ γινέσθω.
\vs{15}~Παρακαλῶ δὲ ὑμᾶς, \adeA{ἀδελφοί}· οἴδατε τὴν οἰκίαν Στεφανᾶ, ὅτι ἐστὶν ἀπαρχὴ τῆς Ἀχαΐας καὶ εἰς διακονίαν τοῖς ἁγίοις ἔταξαν ἑαυτούς·
\vs{16}~ἵνα καὶ ὑμεῖς ὑποτάσσησθε τοῖς τοιούτοις καὶ παντὶ τῷ συνεργοῦντι καὶ κοπιῶντι.
\vs{17}~χαίρω δὲ ἐπὶ τῇ παρουσίᾳ Στεφανᾶ καὶ Φορτουνάτου καὶ Ἀχαϊκοῦ, ὅτι τὸ ὑμέτερον ὑστέρημα οὗτοι ἀνεπλήρωσαν,
\vs{18}~ἀνέπαυσαν γὰρ τὸ ἐμὸν πνεῦμα καὶ τὸ ὑμῶν. ἐπιγινώσκετε οὖν τοὺς τοιούτους.
\vs{19}~Ἀσπάζονται ὑμᾶς αἱ ἐκκλησίαι τῆς Ἀσίας. ἀσπάζεται ὑμᾶς ἐν κυρίῳ πολλὰ Ἀκύλας καὶ Πρίσκα σὺν τῇ κατ’ οἶκον αὐτῶν ἐκκλησίᾳ.
\vs{20}~ἀσπάζονται ὑμᾶς οἱ \adeA{ἀδελφοὶ} πάντες. ἀσπάσασθε ἀλλήλους ἐν φιλήματι ἁγίῳ.
\vs{21}~Ὁ ἀσπασμὸς τῇ ἐμῇ χειρὶ Παύλου.
\vs{22}~εἴ τις οὐ φιλεῖ τὸν κύριον, ἤτω ἀνάθεμα. Μαράνα θά.
\vs{23}~ἡ χάρις τοῦ κυρίου Ἰησοῦ μεθ’ ὑμῶν.
\vs{24}~ἡ ἀγάπη μου μετὰ πάντων ὑμῶν ἐν Χριστῷ Ἰησοῦ.

\booksection{2 Corinthians}{Πρὸς Κορινθίους β΄}

\chapheader{1}
\vs{1}~Παῦλος ἀπόστολος Χριστοῦ Ἰησοῦ διὰ θελήματος θεοῦ καὶ Τιμόθεος ὁ \adeA{ἀδελφὸς} τῇ ἐκκλησίᾳ τοῦ θεοῦ τῇ οὔσῃ ἐν Κορίνθῳ, σὺν τοῖς ἁγίοις πᾶσιν τοῖς οὖσιν ἐν ὅλῃ τῇ Ἀχαΐᾳ·
\vs{2}~χάρις ὑμῖν καὶ εἰρήνη ἀπὸ θεοῦ πατρὸς ἡμῶν καὶ κυρίου Ἰησοῦ Χριστοῦ.
\vs{3}~Εὐλογητὸς ὁ θεὸς καὶ πατὴρ τοῦ κυρίου ἡμῶν Ἰησοῦ Χριστοῦ, ὁ πατὴρ τῶν οἰκτιρμῶν καὶ θεὸς πάσης παρακλήσεως,
\vs{4}~ὁ παρακαλῶν ἡμᾶς ἐπὶ πάσῃ τῇ θλίψει ἡμῶν, εἰς τὸ δύνασθαι ἡμᾶς παρακαλεῖν τοὺς ἐν πάσῃ θλίψει διὰ τῆς παρακλήσεως ἧς παρακαλούμεθα αὐτοὶ ὑπὸ τοῦ θεοῦ.
\vs{5}~ὅτι καθὼς περισσεύει τὰ παθήματα τοῦ Χριστοῦ εἰς ἡμᾶς, οὕτως διὰ τοῦ Χριστοῦ περισσεύει καὶ ἡ παράκλησις ἡμῶν.
\vs{6}~εἴτε δὲ θλιβόμεθα, ὑπὲρ τῆς ὑμῶν παρακλήσεως καὶ σωτηρίας· εἴτε παρακαλούμεθα, ὑπὲρ τῆς ὑμῶν παρακλήσεως τῆς ἐνεργουμένης ἐν ὑπομονῇ τῶν αὐτῶν παθημάτων ὧν καὶ ἡμεῖς πάσχομεν,
\vs{7}~καὶ ἡ ἐλπὶς ἡμῶν βεβαία ὑπὲρ ὑμῶν· εἰδότες ὅτι ὡς κοινωνοί ἐστε τῶν παθημάτων, οὕτως καὶ τῆς παρακλήσεως.
\vs{8}~Οὐ γὰρ θέλομεν ὑμᾶς ἀγνοεῖν, \adeA{ἀδελφοί}, ὑπὲρ τῆς θλίψεως ἡμῶν τῆς γενομένης ἐν τῇ Ἀσίᾳ, ὅτι καθ’ ὑπερβολὴν ὑπὲρ δύναμιν ἐβαρήθημεν, ὥστε ἐξαπορηθῆναι ἡμᾶς καὶ τοῦ ζῆν·
\vs{9}~ἀλλὰ αὐτοὶ ἐν ἑαυτοῖς τὸ ἀπόκριμα τοῦ θανάτου ἐσχήκαμεν, ἵνα μὴ πεποιθότες ὦμεν ἐφ’ ἑαυτοῖς ἀλλ’ ἐπὶ τῷ θεῷ τῷ ἐγείροντι τοὺς νεκρούς·
\vs{10}~ὃς ἐκ τηλικούτου θανάτου ἐρρύσατο ἡμᾶς καὶ ῥύσεται, εἰς ὃν ἠλπίκαμεν ὅτι καὶ ἔτι ῥύσεται,
\vs{11}~συνυπουργούντων καὶ ὑμῶν ὑπὲρ ἡμῶν τῇ δεήσει, ἵνα ἐκ πολλῶν προσώπων τὸ εἰς ἡμᾶς χάρισμα διὰ πολλῶν εὐχαριστηθῇ ὑπὲρ ἡμῶν.
\vs{12}~Ἡ γὰρ καύχησις ἡμῶν αὕτη ἐστίν, τὸ μαρτύριον τῆς συνειδήσεως ἡμῶν, ὅτι ἐν ἁγιότητι καὶ εἰλικρινείᾳ τοῦ θεοῦ, οὐκ ἐν σοφίᾳ σαρκικῇ ἀλλ’ ἐν χάριτι θεοῦ, ἀνεστράφημεν ἐν τῷ κόσμῳ, περισσοτέρως δὲ πρὸς ὑμᾶς·
\vs{13}~οὐ γὰρ ἄλλα γράφομεν ὑμῖν ἀλλ’ ἢ ἃ ἀναγινώσκετε ἢ καὶ ἐπιγινώσκετε, ἐλπίζω δὲ ὅτι ἕως τέλους ἐπιγνώσεσθε,
\vs{14}~καθὼς καὶ ἐπέγνωτε ἡμᾶς ἀπὸ μέρους, ὅτι καύχημα ὑμῶν ἐσμεν καθάπερ καὶ ὑμεῖς ἡμῶν ἐν τῇ ἡμέρᾳ τοῦ κυρίου ἡμῶν Ἰησοῦ.
\vs{15}~Καὶ ταύτῃ τῇ πεποιθήσει ἐβουλόμην πρότερον πρὸς ὑμᾶς ἐλθεῖν, ἵνα δευτέραν χάριν σχῆτε,
\vs{16}~καὶ δι’ ὑμῶν διελθεῖν εἰς Μακεδονίαν, καὶ πάλιν ἀπὸ Μακεδονίας ἐλθεῖν πρὸς ὑμᾶς καὶ ὑφ’ ὑμῶν προπεμφθῆναι εἰς τὴν Ἰουδαίαν.
\vs{17}~τοῦτο οὖν βουλόμενος μήτι ἄρα τῇ ἐλαφρίᾳ ἐχρησάμην; ἢ ἃ βουλεύομαι κατὰ σάρκα βουλεύομαι, ἵνα ᾖ παρ’ ἐμοὶ τὸ Ναὶ ναὶ καὶ τὸ Οὒ οὔ;
\vs{18}~πιστὸς δὲ ὁ θεὸς ὅτι ὁ λόγος ἡμῶν ὁ πρὸς ὑμᾶς οὐκ ἔστιν Ναὶ καὶ Οὔ.
\vs{19}~ὁ τοῦ θεοῦ γὰρ υἱὸς Ἰησοῦς Χριστὸς ὁ ἐν ὑμῖν δι’ ἡμῶν κηρυχθείς, δι’ ἐμοῦ καὶ Σιλουανοῦ καὶ Τιμοθέου, οὐκ ἐγένετο Ναὶ καὶ Οὒ, ἀλλὰ Ναὶ ἐν αὐτῷ γέγονεν·
\vs{20}~ὅσαι γὰρ ἐπαγγελίαι θεοῦ, ἐν αὐτῷ τὸ Ναί· διὸ καὶ δι’ αὐτοῦ τὸ Ἀμὴν τῷ θεῷ πρὸς δόξαν δι’ ἡμῶν.
\vs{21}~ὁ δὲ βεβαιῶν ἡμᾶς σὺν ὑμῖν εἰς Χριστὸν καὶ χρίσας ἡμᾶς θεός,
\vs{22}~ὁ καὶ σφραγισάμενος ἡμᾶς καὶ δοὺς τὸν ἀρραβῶνα τοῦ πνεύματος ἐν ταῖς καρδίαις ἡμῶν.
\vs{23}~Ἐγὼ δὲ μάρτυρα τὸν θεὸν ἐπικαλοῦμαι ἐπὶ τὴν ἐμὴν ψυχήν, ὅτι φειδόμενος ὑμῶν οὐκέτι ἦλθον εἰς Κόρινθον.
\vs{24}~οὐχ ὅτι κυριεύομεν ὑμῶν τῆς πίστεως, ἀλλὰ συνεργοί ἐσμεν τῆς χαρᾶς ὑμῶν, τῇ γὰρ πίστει ἑστήκατε.

\chapheader{2}
\vs{1}~ἔκρινα γὰρ ἐμαυτῷ τοῦτο, τὸ μὴ πάλιν ἐν λύπῃ πρὸς ὑμᾶς ἐλθεῖν·
\vs{2}~εἰ γὰρ ἐγὼ λυπῶ ὑμᾶς, καὶ τίς ὁ εὐφραίνων με εἰ μὴ ὁ λυπούμενος ἐξ ἐμοῦ;
\vs{3}~καὶ ἔγραψα τοῦτο αὐτὸ ἵνα μὴ ἐλθὼν λύπην σχῶ ἀφ’ ὧν ἔδει με χαίρειν, πεποιθὼς ἐπὶ πάντας ὑμᾶς ὅτι ἡ ἐμὴ χαρὰ πάντων ὑμῶν ἐστιν.
\vs{4}~ἐκ γὰρ πολλῆς θλίψεως καὶ συνοχῆς καρδίας ἔγραψα ὑμῖν διὰ πολλῶν δακρύων, οὐχ ἵνα λυπηθῆτε, ἀλλὰ τὴν ἀγάπην ἵνα γνῶτε ἣν ἔχω περισσοτέρως εἰς ὑμᾶς.
\vs{5}~Εἰ δέ τις λελύπηκεν, οὐκ ἐμὲ λελύπηκεν, ἀλλὰ ἀπὸ μέρους ἵνα μὴ ἐπιβαρῶ πάντας ὑμᾶς.
\vs{6}~ἱκανὸν τῷ τοιούτῳ ἡ ἐπιτιμία αὕτη ἡ ὑπὸ τῶν πλειόνων,
\vs{7}~ὥστε τοὐναντίον μᾶλλον ὑμᾶς χαρίσασθαι καὶ παρακαλέσαι, μή πως τῇ περισσοτέρᾳ λύπῃ καταποθῇ ὁ τοιοῦτος.
\vs{8}~διὸ παρακαλῶ ὑμᾶς κυρῶσαι εἰς αὐτὸν ἀγάπην·
\vs{9}~εἰς τοῦτο γὰρ καὶ ἔγραψα ἵνα γνῶ τὴν δοκιμὴν ὑμῶν, εἰ εἰς πάντα ὑπήκοοί ἐστε.
\vs{10}~ᾧ δέ τι χαρίζεσθε, κἀγώ· καὶ γὰρ ἐγὼ ὃ κεχάρισμαι, εἴ τι κεχάρισμαι, δι’ ὑμᾶς ἐν προσώπῳ Χριστοῦ,
\vs{11}~ἵνα μὴ πλεονεκτηθῶμεν ὑπὸ τοῦ Σατανᾶ, οὐ γὰρ αὐτοῦ τὰ νοήματα ἀγνοοῦμεν.
\vs{12}~Ἐλθὼν δὲ εἰς τὴν Τρῳάδα εἰς τὸ εὐαγγέλιον τοῦ Χριστοῦ καὶ θύρας μοι ἀνεῳγμένης ἐν κυρίῳ,
\vs{13}~οὐκ ἔσχηκα ἄνεσιν τῷ πνεύματί μου τῷ μὴ εὑρεῖν με Τίτον τὸν \adeA{ἀδελφόν} μου, ἀλλὰ ἀποταξάμενος αὐτοῖς ἐξῆλθον εἰς Μακεδονίαν.
\vs{14}~Τῷ δὲ θεῷ χάρις τῷ πάντοτε θριαμβεύοντι ἡμᾶς ἐν τῷ Χριστῷ καὶ τὴν ὀσμὴν τῆς γνώσεως αὐτοῦ φανεροῦντι δι’ ἡμῶν ἐν παντὶ τόπῳ·
\vs{15}~ὅτι Χριστοῦ εὐωδία ἐσμὲν τῷ θεῷ ἐν τοῖς σῳζομένοις καὶ ἐν τοῖς ἀπολλυμένοις,
\vs{16}~οἷς μὲν ὀσμὴ ἐκ θανάτου εἰς θάνατον, οἷς δὲ ὀσμὴ ἐκ ζωῆς εἰς ζωήν. καὶ πρὸς ταῦτα τίς ἱκανός;
\vs{17}~οὐ γάρ ἐσμεν ὡς οἱ πολλοὶ καπηλεύοντες τὸν λόγον τοῦ θεοῦ, ἀλλ’ ὡς ἐξ εἰλικρινείας, ἀλλ’ ὡς ἐκ θεοῦ κατέναντι θεοῦ ἐν Χριστῷ λαλοῦμεν.

\chapheader{3}
\vs{1}~Ἀρχόμεθα πάλιν ἑαυτοὺς συνιστάνειν; ἢ μὴ χρῄζομεν ὥς τινες συστατικῶν ἐπιστολῶν πρὸς ὑμᾶς ἢ ἐξ ὑμῶν;
\vs{2}~ἡ ἐπιστολὴ ἡμῶν ὑμεῖς ἐστε, ἐγγεγραμμένη ἐν ταῖς καρδίαις ἡμῶν, γινωσκομένη καὶ ἀναγινωσκομένη ὑπὸ πάντων ἀνθρώπων·
\vs{3}~φανερούμενοι ὅτι ἐστὲ ἐπιστολὴ Χριστοῦ διακονηθεῖσα ὑφ’ ἡμῶν, ἐγγεγραμμένη οὐ μέλανι ἀλλὰ πνεύματι θεοῦ ζῶντος, οὐκ ἐν πλαξὶν λιθίναις ἀλλ’ ἐν πλαξὶν καρδίαις σαρκίναις.
\vs{4}~Πεποίθησιν δὲ τοιαύτην ἔχομεν διὰ τοῦ Χριστοῦ πρὸς τὸν θεόν.
\vs{5}~οὐχ ὅτι ἀφ’ ἑαυτῶν ἱκανοί ἐσμεν λογίσασθαί τι ὡς ἐξ αὑτῶν, ἀλλ’ ἡ ἱκανότης ἡμῶν ἐκ τοῦ θεοῦ,
\vs{6}~ὃς καὶ ἱκάνωσεν ἡμᾶς διακόνους καινῆς διαθήκης, οὐ γράμματος ἀλλὰ πνεύματος, τὸ γὰρ γράμμα ἀποκτέννει, τὸ δὲ πνεῦμα ζῳοποιεῖ.
\vs{7}~Εἰ δὲ ἡ διακονία τοῦ θανάτου ἐν γράμμασιν ἐντετυπωμένη λίθοις ἐγενήθη ἐν δόξῃ, ὥστε μὴ δύνασθαι ἀτενίσαι τοὺς υἱοὺς Ἰσραὴλ εἰς τὸ πρόσωπον Μωϋσέως διὰ τὴν δόξαν τοῦ προσώπου αὐτοῦ τὴν καταργουμένην,
\vs{8}~πῶς οὐχὶ μᾶλλον ἡ διακονία τοῦ πνεύματος ἔσται ἐν δόξῃ;
\vs{9}~εἰ γὰρ τῇ διακονίᾳ τῆς κατακρίσεως δόξα, πολλῷ μᾶλλον περισσεύει ἡ διακονία τῆς δικαιοσύνης δόξῃ.
\vs{10}~καὶ γὰρ οὐ δεδόξασται τὸ δεδοξασμένον ἐν τούτῳ τῷ μέρει εἵνεκεν τῆς ὑπερβαλλούσης δόξης·
\vs{11}~εἰ γὰρ τὸ καταργούμενον διὰ δόξης, πολλῷ μᾶλλον τὸ μένον ἐν δόξῃ.
\vs{12}~Ἔχοντες οὖν τοιαύτην ἐλπίδα πολλῇ παρρησίᾳ χρώμεθα,
\vs{13}~καὶ οὐ καθάπερ Μωϋσῆς ἐτίθει κάλυμμα ἐπὶ τὸ πρόσωπον αὐτοῦ, πρὸς τὸ μὴ ἀτενίσαι τοὺς υἱοὺς Ἰσραὴλ εἰς τὸ τέλος τοῦ καταργουμένου.
\vs{14}~ἀλλὰ ἐπωρώθη τὰ νοήματα αὐτῶν. ἄχρι γὰρ τῆς σήμερον ἡμέρας τὸ αὐτὸ κάλυμμα ἐπὶ τῇ ἀναγνώσει τῆς παλαιᾶς διαθήκης μένει μὴ ἀνακαλυπτόμενον, ὅτι ἐν Χριστῷ καταργεῖται,
\vs{15}~ἀλλ’ ἕως σήμερον ἡνίκα ἂν ἀναγινώσκηται Μωϋσῆς κάλυμμα ἐπὶ τὴν καρδίαν αὐτῶν κεῖται·
\vs{16}~ἡνίκα δὲ ἐὰν ἐπιστρέψῃ πρὸς κύριον, περιαιρεῖται τὸ κάλυμμα.
\vs{17}~ὁ δὲ κύριος τὸ πνεῦμά ἐστιν· οὗ δὲ τὸ πνεῦμα κυρίου, ἐλευθερία.
\vs{18}~ἡμεῖς δὲ πάντες ἀνακεκαλυμμένῳ προσώπῳ τὴν δόξαν κυρίου κατοπτριζόμενοι τὴν αὐτὴν εἰκόνα μεταμορφούμεθα ἀπὸ δόξης εἰς δόξαν, καθάπερ ἀπὸ κυρίου πνεύματος.

\chapheader{4}
\vs{1}~Διὰ τοῦτο, ἔχοντες τὴν διακονίαν ταύτην καθὼς ἠλεήθημεν, οὐκ ἐγκακοῦμεν,
\vs{2}~ἀλλὰ ἀπειπάμεθα τὰ κρυπτὰ τῆς αἰσχύνης, μὴ περιπατοῦντες ἐν πανουργίᾳ μηδὲ δολοῦντες τὸν λόγον τοῦ θεοῦ, ἀλλὰ τῇ φανερώσει τῆς ἀληθείας συνιστάνοντες ἑαυτοὺς πρὸς πᾶσαν συνείδησιν ἀνθρώπων ἐνώπιον τοῦ θεοῦ.
\vs{3}~εἰ δὲ καὶ ἔστιν κεκαλυμμένον τὸ εὐαγγέλιον ἡμῶν, ἐν τοῖς ἀπολλυμένοις ἐστὶν κεκαλυμμένον,
\vs{4}~ἐν οἷς ὁ θεὸς τοῦ αἰῶνος τούτου ἐτύφλωσεν τὰ νοήματα τῶν ἀπίστων εἰς τὸ μὴ αὐγάσαι τὸν φωτισμὸν τοῦ εὐαγγελίου τῆς δόξης τοῦ Χριστοῦ, ὅς ἐστιν εἰκὼν τοῦ θεοῦ.
\vs{5}~οὐ γὰρ ἑαυτοὺς κηρύσσομεν ἀλλὰ Χριστὸν Ἰησοῦν κύριον, ἑαυτοὺς δὲ δούλους ὑμῶν διὰ Ἰησοῦν.
\vs{6}~ὅτι ὁ θεὸς ὁ εἰπών· Ἐκ σκότους φῶς λάμψει, ὃς ἔλαμψεν ἐν ταῖς καρδίαις ἡμῶν πρὸς φωτισμὸν τῆς γνώσεως τῆς δόξης τοῦ θεοῦ ἐν προσώπῳ Χριστοῦ.
\vs{7}~Ἔχομεν δὲ τὸν θησαυρὸν τοῦτον ἐν ὀστρακίνοις σκεύεσιν, ἵνα ἡ ὑπερβολὴ τῆς δυνάμεως ᾖ τοῦ θεοῦ καὶ μὴ ἐξ ἡμῶν·
\vs{8}~ἐν παντὶ θλιβόμενοι ἀλλ’ οὐ στενοχωρούμενοι, ἀπορούμενοι ἀλλ’ οὐκ ἐξαπορούμενοι,
\vs{9}~διωκόμενοι ἀλλ’ οὐκ ἐγκαταλειπόμενοι, καταβαλλόμενοι ἀλλ’ οὐκ ἀπολλύμενοι,
\vs{10}~πάντοτε τὴν νέκρωσιν τοῦ Ἰησοῦ ἐν τῷ σώματι περιφέροντες, ἵνα καὶ ἡ ζωὴ τοῦ Ἰησοῦ ἐν τῷ σώματι ἡμῶν φανερωθῇ·
\vs{11}~ἀεὶ γὰρ ἡμεῖς οἱ ζῶντες εἰς θάνατον παραδιδόμεθα διὰ Ἰησοῦν, ἵνα καὶ ἡ ζωὴ τοῦ Ἰησοῦ φανερωθῇ ἐν τῇ θνητῇ σαρκὶ ἡμῶν.
\vs{12}~ὥστε ὁ θάνατος ἐν ἡμῖν ἐνεργεῖται, ἡ δὲ ζωὴ ἐν ὑμῖν.
\vs{13}~Ἔχοντες δὲ τὸ αὐτὸ πνεῦμα τῆς πίστεως, κατὰ τὸ γεγραμμένον· Ἐπίστευσα, διὸ ἐλάλησα, καὶ ἡμεῖς πιστεύομεν, διὸ καὶ λαλοῦμεν,
\vs{14}~εἰδότες ὅτι ὁ ἐγείρας τὸν Ἰησοῦν καὶ ἡμᾶς σὺν Ἰησοῦ ἐγερεῖ καὶ παραστήσει σὺν ὑμῖν.
\vs{15}~τὰ γὰρ πάντα δι’ ὑμᾶς, ἵνα ἡ χάρις πλεονάσασα διὰ τῶν πλειόνων τὴν εὐχαριστίαν περισσεύσῃ εἰς τὴν δόξαν τοῦ θεοῦ.
\vs{16}~Διὸ οὐκ ἐγκακοῦμεν, ἀλλ’ εἰ καὶ ὁ ἔξω ἡμῶν ἄνθρωπος διαφθείρεται, ἀλλ’ ὁ ἔσω ἡμῶν ἀνακαινοῦται ἡμέρᾳ καὶ ἡμέρᾳ.
\vs{17}~τὸ γὰρ παραυτίκα ἐλαφρὸν τῆς θλίψεως ἡμῶν καθ’ ὑπερβολὴν εἰς ὑπερβολὴν αἰώνιον βάρος δόξης κατεργάζεται ἡμῖν,
\vs{18}~μὴ σκοπούντων ἡμῶν τὰ βλεπόμενα ἀλλὰ τὰ μὴ βλεπόμενα, τὰ γὰρ βλεπόμενα πρόσκαιρα, τὰ δὲ μὴ βλεπόμενα αἰώνια.

\chapheader{5}
\vs{1}~Οἴδαμεν γὰρ ὅτι ἐὰν ἡ ἐπίγειος ἡμῶν οἰκία τοῦ σκήνους καταλυθῇ, οἰκοδομὴν ἐκ θεοῦ ἔχομεν οἰκίαν ἀχειροποίητον αἰώνιον ἐν τοῖς οὐρανοῖς.
\vs{2}~καὶ γὰρ ἐν τούτῳ στενάζομεν, τὸ οἰκητήριον ἡμῶν τὸ ἐξ οὐρανοῦ ἐπενδύσασθαι ἐπιποθοῦντες,
\vs{3}~εἴ γε καὶ ἐνδυσάμενοι οὐ γυμνοὶ εὑρεθησόμεθα.
\vs{4}~καὶ γὰρ οἱ ὄντες ἐν τῷ σκήνει στενάζομεν βαρούμενοι ἐφ’ ᾧ οὐ θέλομεν ἐκδύσασθαι ἀλλ’ ἐπενδύσασθαι, ἵνα καταποθῇ τὸ θνητὸν ὑπὸ τῆς ζωῆς.
\vs{5}~ὁ δὲ κατεργασάμενος ἡμᾶς εἰς αὐτὸ τοῦτο θεός, ὁ δοὺς ἡμῖν τὸν ἀρραβῶνα τοῦ πνεύματος.
\vs{6}~Θαρροῦντες οὖν πάντοτε καὶ εἰδότες ὅτι ἐνδημοῦντες ἐν τῷ σώματι ἐκδημοῦμεν ἀπὸ τοῦ κυρίου,
\vs{7}~διὰ πίστεως γὰρ περιπατοῦμεν, οὐ διὰ εἴδους—
\vs{8}~θαρροῦμεν δὲ καὶ εὐδοκοῦμεν μᾶλλον ἐκδημῆσαι ἐκ τοῦ σώματος καὶ ἐνδημῆσαι πρὸς τὸν κύριον·
\vs{9}~διὸ καὶ φιλοτιμούμεθα, εἴτε ἐνδημοῦντες εἴτε ἐκδημοῦντες, εὐάρεστοι αὐτῷ εἶναι.
\vs{10}~τοὺς γὰρ πάντας ἡμᾶς φανερωθῆναι δεῖ ἔμπροσθεν τοῦ βήματος τοῦ Χριστοῦ, ἵνα κομίσηται ἕκαστος τὰ διὰ τοῦ σώματος πρὸς ἃ ἔπραξεν, εἴτε ἀγαθὸν εἴτε φαῦλον.
\vs{11}~Εἰδότες οὖν τὸν φόβον τοῦ κυρίου ἀνθρώπους πείθομεν, θεῷ δὲ πεφανερώμεθα· ἐλπίζω δὲ καὶ ἐν ταῖς συνειδήσεσιν ὑμῶν πεφανερῶσθαι.
\vs{12}~οὐ πάλιν ἑαυτοὺς συνιστάνομεν ὑμῖν, ἀλλὰ ἀφορμὴν διδόντες ὑμῖν καυχήματος ὑπὲρ ἡμῶν, ἵνα ἔχητε πρὸς τοὺς ἐν προσώπῳ καυχωμένους καὶ μὴ ἐν καρδίᾳ.
\vs{13}~εἴτε γὰρ ἐξέστημεν, θεῷ· εἴτε σωφρονοῦμεν, ὑμῖν.
\vs{14}~ἡ γὰρ ἀγάπη τοῦ Χριστοῦ συνέχει ἡμᾶς, κρίναντας τοῦτο ὅτι εἷς ὑπὲρ πάντων ἀπέθανεν· ἄρα οἱ πάντες ἀπέθανον·
\vs{15}~καὶ ὑπὲρ πάντων ἀπέθανεν ἵνα οἱ ζῶντες μηκέτι ἑαυτοῖς ζῶσιν ἀλλὰ τῷ ὑπὲρ αὐτῶν ἀποθανόντι καὶ ἐγερθέντι.
\vs{16}~Ὥστε ἡμεῖς ἀπὸ τοῦ νῦν οὐδένα οἴδαμεν κατὰ σάρκα· εἰ καὶ ἐγνώκαμεν κατὰ σάρκα Χριστόν, ἀλλὰ νῦν οὐκέτι γινώσκομεν.
\vs{17}~ὥστε εἴ τις ἐν Χριστῷ, καινὴ κτίσις· τὰ ἀρχαῖα παρῆλθεν, ἰδοὺ γέγονεν καινά·
\vs{18}~τὰ δὲ πάντα ἐκ τοῦ θεοῦ τοῦ καταλλάξαντος ἡμᾶς ἑαυτῷ διὰ Χριστοῦ καὶ δόντος ἡμῖν τὴν διακονίαν τῆς καταλλαγῆς,
\vs{19}~ὡς ὅτι θεὸς ἦν ἐν Χριστῷ κόσμον καταλλάσσων ἑαυτῷ, μὴ λογιζόμενος αὐτοῖς τὰ παραπτώματα αὐτῶν, καὶ θέμενος ἐν ἡμῖν τὸν λόγον τῆς καταλλαγῆς.
\vs{20}~ὑπὲρ Χριστοῦ οὖν πρεσβεύομεν ὡς τοῦ θεοῦ παρακαλοῦντος δι’ ἡμῶν· δεόμεθα ὑπὲρ Χριστοῦ, καταλλάγητε τῷ θεῷ.
\vs{21}~τὸν μὴ γνόντα ἁμαρτίαν ὑπὲρ ἡμῶν ἁμαρτίαν ἐποίησεν, ἵνα ἡμεῖς γενώμεθα δικαιοσύνη θεοῦ ἐν αὐτῷ.

\chapheader{6}
\vs{1}~Συνεργοῦντες δὲ καὶ παρακαλοῦμεν μὴ εἰς κενὸν τὴν χάριν τοῦ θεοῦ δέξασθαι ὑμᾶς·
\vs{2}~λέγει γάρ· Καιρῷ δεκτῷ ἐπήκουσά σου καὶ ἐν ἡμέρᾳ σωτηρίας ἐβοήθησά σοι· ἰδοὺ νῦν καιρὸς εὐπρόσδεκτος, ἰδοὺ νῦν ἡμέρα σωτηρίας·
\vs{3}~μηδεμίαν ἐν μηδενὶ διδόντες προσκοπήν, ἵνα μὴ μωμηθῇ ἡ διακονία,
\vs{4}~ἀλλ’ ἐν παντὶ συνιστάνοντες ἑαυτοὺς ὡς θεοῦ διάκονοι· ἐν ὑπομονῇ πολλῇ, ἐν θλίψεσιν, ἐν ἀνάγκαις, ἐν στενοχωρίαις,
\vs{5}~ἐν πληγαῖς, ἐν φυλακαῖς, ἐν ἀκαταστασίαις, ἐν κόποις, ἐν ἀγρυπνίαις, ἐν νηστείαις,
\vs{6}~ἐν ἁγνότητι, ἐν γνώσει, ἐν μακροθυμίᾳ, ἐν χρηστότητι, ἐν πνεύματι ἁγίῳ, ἐν ἀγάπῃ ἀνυποκρίτῳ,
\vs{7}~ἐν λόγῳ ἀληθείας, ἐν δυνάμει θεοῦ· διὰ τῶν ὅπλων τῆς δικαιοσύνης τῶν δεξιῶν καὶ ἀριστερῶν,
\vs{8}~διὰ δόξης καὶ ἀτιμίας, διὰ δυσφημίας καὶ εὐφημίας· ὡς πλάνοι καὶ ἀληθεῖς,
\vs{9}~ὡς ἀγνοούμενοι καὶ ἐπιγινωσκόμενοι, ὡς ἀποθνῄσκοντες καὶ ἰδοὺ ζῶμεν, ὡς παιδευόμενοι καὶ μὴ θανατούμενοι,
\vs{10}~ὡς λυπούμενοι ἀεὶ δὲ χαίροντες, ὡς πτωχοὶ πολλοὺς δὲ πλουτίζοντες, ὡς μηδὲν ἔχοντες καὶ πάντα κατέχοντες.
\vs{11}~Τὸ στόμα ἡμῶν ἀνέῳγεν πρὸς ὑμᾶς, Κορίνθιοι, ἡ καρδία ἡμῶν πεπλάτυνται·
\vs{12}~οὐ στενοχωρεῖσθε ἐν ἡμῖν, στενοχωρεῖσθε δὲ ἐν τοῖς σπλάγχνοις ὑμῶν·
\vs{13}~τὴν δὲ αὐτὴν ἀντιμισθίαν, ὡς τέκνοις λέγω, πλατύνθητε καὶ ὑμεῖς.
\vs{14}~Μὴ γίνεσθε ἑτεροζυγοῦντες ἀπίστοις· τίς γὰρ μετοχὴ δικαιοσύνῃ καὶ ἀνομίᾳ, ἢ τίς κοινωνία φωτὶ πρὸς σκότος;
\vs{15}~τίς δὲ συμφώνησις Χριστοῦ πρὸς Βελιάρ, ἢ τίς μερὶς πιστῷ μετὰ ἀπίστου;
\vs{16}~τίς δὲ συγκατάθεσις ναῷ θεοῦ μετὰ εἰδώλων; ἡμεῖς γὰρ ναὸς θεοῦ ἐσμεν ζῶντος· καθὼς εἶπεν ὁ θεὸς ὅτι Ἐνοικήσω ἐν αὐτοῖς καὶ ἐμπεριπατήσω, καὶ ἔσομαι αὐτῶν θεός, καὶ αὐτοὶ ἔσονταί μου λαός.
\vs{17}~διὸ ἐξέλθατε ἐκ μέσου αὐτῶν, καὶ ἀφορίσθητε, λέγει κύριος, καὶ ἀκαθάρτου μὴ ἅπτεσθε· κἀγὼ εἰσδέξομαι ὑμᾶς·
\vs{18}~καὶ ἔσομαι ὑμῖν εἰς πατέρα, καὶ ὑμεῖς ἔσεσθέ μοι εἰς υἱοὺς καὶ θυγατέρας, λέγει κύριος παντοκράτωρ.

\chapheader{7}
\vs{1}~ταύτας οὖν ἔχοντες τὰς ἐπαγγελίας, ἀγαπητοί, καθαρίσωμεν ἑαυτοὺς ἀπὸ παντὸς μολυσμοῦ σαρκὸς καὶ πνεύματος, ἐπιτελοῦντες ἁγιωσύνην ἐν φόβῳ θεοῦ.
\vs{2}~Χωρήσατε ἡμᾶς· οὐδένα ἠδικήσαμεν, οὐδένα ἐφθείραμεν, οὐδένα ἐπλεονεκτήσαμεν.
\vs{3}~πρὸς κατάκρισιν οὐ λέγω, προείρηκα γὰρ ὅτι ἐν ταῖς καρδίαις ἡμῶν ἐστε εἰς τὸ συναποθανεῖν καὶ συζῆν.
\vs{4}~πολλή μοι παρρησία πρὸς ὑμᾶς, πολλή μοι καύχησις ὑπὲρ ὑμῶν· πεπλήρωμαι τῇ παρακλήσει, ὑπερπερισσεύομαι τῇ χαρᾷ ἐπὶ πάσῃ τῇ θλίψει ἡμῶν.
\vs{5}~Καὶ γὰρ ἐλθόντων ἡμῶν εἰς Μακεδονίαν οὐδεμίαν ἔσχηκεν ἄνεσιν ἡ σὰρξ ἡμῶν, ἀλλ’ ἐν παντὶ θλιβόμενοι— ἔξωθεν μάχαι, ἔσωθεν φόβοι—
\vs{6}~ἀλλ’ ὁ παρακαλῶν τοὺς ταπεινοὺς παρεκάλεσεν ἡμᾶς ὁ θεὸς ἐν τῇ παρουσίᾳ Τίτου·
\vs{7}~οὐ μόνον δὲ ἐν τῇ παρουσίᾳ αὐτοῦ, ἀλλὰ καὶ ἐν τῇ παρακλήσει ᾗ παρεκλήθη ἐφ’ ὑμῖν, ἀναγγέλλων ἡμῖν τὴν ὑμῶν ἐπιπόθησιν, τὸν ὑμῶν ὀδυρμόν, τὸν ὑμῶν ζῆλον ὑπὲρ ἐμοῦ, ὥστε με μᾶλλον χαρῆναι.
\vs{8}~ὅτι εἰ καὶ ἐλύπησα ὑμᾶς ἐν τῇ ἐπιστολῇ, οὐ μεταμέλομαι· εἰ καὶ μετεμελόμην (βλέπω ὅτι ἡ ἐπιστολὴ ἐκείνη εἰ καὶ πρὸς ὥραν ἐλύπησεν ὑμᾶς),
\vs{9}~νῦν χαίρω, οὐχ ὅτι ἐλυπήθητε, ἀλλ’ ὅτι ἐλυπήθητε εἰς μετάνοιαν, ἐλυπήθητε γὰρ κατὰ θεόν, ἵνα ἐν μηδενὶ ζημιωθῆτε ἐξ ἡμῶν.
\vs{10}~ἡ γὰρ κατὰ θεὸν λύπη μετάνοιαν εἰς σωτηρίαν ἀμεταμέλητον ἐργάζεται· ἡ δὲ τοῦ κόσμου λύπη θάνατον κατεργάζεται.
\vs{11}~ἰδοὺ γὰρ αὐτὸ τοῦτο τὸ κατὰ θεὸν λυπηθῆναι πόσην κατειργάσατο ὑμῖν σπουδήν, ἀλλὰ ἀπολογίαν, ἀλλὰ ἀγανάκτησιν, ἀλλὰ φόβον, ἀλλὰ ἐπιπόθησιν, ἀλλὰ ζῆλον, ἀλλὰ ἐκδίκησιν· ἐν παντὶ συνεστήσατε ἑαυτοὺς ἁγνοὺς εἶναι τῷ πράγματι.
\vs{12}~ἄρα εἰ καὶ ἔγραψα ὑμῖν, οὐχ ἕνεκεν τοῦ ἀδικήσαντος, οὐδὲ ἕνεκεν τοῦ ἀδικηθέντος, ἀλλ’ 1ἕνεκεν τοῦ φανερωθῆναι τὴν σπουδὴν ὑμῶν τὴν ὑπὲρ ἡμῶν πρὸς ὑμᾶς ἐνώπιον τοῦ θεοῦ.
\vs{13}~διὰ τοῦτο παρακεκλήμεθα. Ἐπὶ δὲ τῇ παρακλήσει ἡμῶν περισσοτέρως μᾶλλον ἐχάρημεν ἐπὶ τῇ χαρᾷ Τίτου, ὅτι ἀναπέπαυται τὸ πνεῦμα αὐτοῦ ἀπὸ πάντων ὑμῶν·
\vs{14}~ὅτι εἴ τι αὐτῷ ὑπὲρ ὑμῶν κεκαύχημαι, οὐ κατῃσχύνθην, ἀλλ’ ὡς πάντα ἐν ἀληθείᾳ ἐλαλήσαμεν ὑμῖν, οὕτως καὶ ἡ καύχησις ἡμῶν ἡ ἐπὶ Τίτου ἀλήθεια ἐγενήθη.
\vs{15}~καὶ τὰ σπλάγχνα αὐτοῦ περισσοτέρως εἰς ὑμᾶς ἐστιν ἀναμιμνῃσκομένου τὴν πάντων ὑμῶν ὑπακοήν, ὡς μετὰ φόβου καὶ τρόμου ἐδέξασθε αὐτόν.
\vs{16}~χαίρω ὅτι ἐν παντὶ θαρρῶ ἐν ὑμῖν.

\chapheader{8}
\vs{1}~Γνωρίζομεν δὲ ὑμῖν, \adeA{ἀδελφοί}, τὴν χάριν τοῦ θεοῦ τὴν δεδομένην ἐν ταῖς ἐκκλησίαις τῆς Μακεδονίας,
\vs{2}~ὅτι ἐν πολλῇ δοκιμῇ θλίψεως ἡ περισσεία τῆς χαρᾶς αὐτῶν καὶ ἡ κατὰ βάθους πτωχεία αὐτῶν ἐπερίσσευσεν εἰς τὸ πλοῦτος τῆς ἁπλότητος αὐτῶν·
\vs{3}~ὅτι κατὰ δύναμιν, μαρτυρῶ, καὶ παρὰ δύναμιν, αὐθαίρετοι
\vs{4}~μετὰ πολλῆς παρακλήσεως δεόμενοι ἡμῶν, τὴν χάριν καὶ τὴν κοινωνίαν τῆς διακονίας τῆς εἰς τοὺς ἁγίους—
\vs{5}~καὶ οὐ καθὼς ἠλπίσαμεν ἀλλ’ ἑαυτοὺς ἔδωκαν πρῶτον τῷ κυρίῳ καὶ ἡμῖν διὰ θελήματος θεοῦ,
\vs{6}~εἰς τὸ παρακαλέσαι ἡμᾶς Τίτον ἵνα καθὼς προενήρξατο οὕτως καὶ ἐπιτελέσῃ εἰς ὑμᾶς καὶ τὴν χάριν ταύτην·
\vs{7}~ἀλλ’ ὥσπερ ἐν παντὶ περισσεύετε, πίστει καὶ λόγῳ καὶ γνώσει καὶ πάσῃ σπουδῇ καὶ τῇ ἐξ ἡμῶν ἐν ὑμῖν ἀγάπῃ, ἵνα καὶ ἐν ταύτῃ τῇ χάριτι περισσεύητε.
\vs{8}~Οὐ κατ’ ἐπιταγὴν λέγω ἀλλὰ διὰ τῆς ἑτέρων σπουδῆς καὶ τὸ τῆς ὑμετέρας ἀγάπης γνήσιον δοκιμάζων·
\vs{9}~γινώσκετε γὰρ τὴν χάριν τοῦ κυρίου ἡμῶν Ἰησοῦ Χριστοῦ, ὅτι δι’ ὑμᾶς ἐπτώχευσεν πλούσιος ὤν, ἵνα ὑμεῖς τῇ ἐκείνου πτωχείᾳ πλουτήσητε.
\vs{10}~καὶ γνώμην ἐν τούτῳ δίδωμι· τοῦτο γὰρ ὑμῖν συμφέρει, οἵτινες οὐ μόνον τὸ ποιῆσαι ἀλλὰ καὶ τὸ θέλειν προενήρξασθε ἀπὸ πέρυσι·
\vs{11}~νυνὶ δὲ καὶ τὸ ποιῆσαι ἐπιτελέσατε, ὅπως καθάπερ ἡ προθυμία τοῦ θέλειν οὕτως καὶ τὸ ἐπιτελέσαι ἐκ τοῦ ἔχειν.
\vs{12}~εἰ γὰρ ἡ προθυμία πρόκειται, καθὸ ἐὰν ἔχῃ εὐπρόσδεκτος, οὐ καθὸ οὐκ ἔχει.
\vs{13}~οὐ γὰρ ἵνα ἄλλοις ἄνεσις, ὑμῖν θλῖψις· ἀλλ’ ἐξ ἰσότητος
\vs{14}~ἐν τῷ νῦν καιρῷ τὸ ὑμῶν περίσσευμα εἰς τὸ ἐκείνων ὑστέρημα, ἵνα καὶ τὸ ἐκείνων περίσσευμα γένηται εἰς τὸ ὑμῶν ὑστέρημα, ὅπως γένηται ἰσότης·
\vs{15}~καθὼς γέγραπται· Ὁ τὸ πολὺ οὐκ ἐπλεόνασεν, καὶ ὁ τὸ ὀλίγον οὐκ ἠλαττόνησεν.
\vs{16}~Χάρις δὲ τῷ θεῷ τῷ διδόντι τὴν αὐτὴν σπουδὴν ὑπὲρ ὑμῶν ἐν τῇ καρδίᾳ Τίτου,
\vs{17}~ὅτι τὴν μὲν παράκλησιν ἐδέξατο, σπουδαιότερος δὲ ὑπάρχων αὐθαίρετος ἐξῆλθεν πρὸς ὑμᾶς.
\vs{18}~συνεπέμψαμεν δὲ μετ’ αὐτοῦ τὸν \adeA{ἀδελφὸν} οὗ ὁ ἔπαινος ἐν τῷ εὐαγγελίῳ διὰ πασῶν τῶν ἐκκλησιῶν—
\vs{19}~οὐ μόνον δὲ ἀλλὰ καὶ χειροτονηθεὶς ὑπὸ τῶν ἐκκλησιῶν συνέκδημος ἡμῶν σὺν τῇ χάριτι ταύτῃ τῇ διακονουμένῃ ὑφ’ ἡμῶν πρὸς τὴν αὐτοῦ τοῦ κυρίου δόξαν καὶ προθυμίαν ἡμῶν—
\vs{20}~στελλόμενοι τοῦτο μή τις ἡμᾶς μωμήσηται ἐν τῇ ἁδρότητι ταύτῃ τῇ διακονουμένῃ ὑφ’ ἡμῶν,
\vs{21}~προνοοῦμεν γὰρ καλὰ οὐ μόνον ἐνώπιον κυρίου ἀλλὰ καὶ ἐνώπιον ἀνθρώπων.
\vs{22}~συνεπέμψαμεν δὲ αὐτοῖς τὸν \adeA{ἀδελφὸν} ἡμῶν ὃν ἐδοκιμάσαμεν ἐν πολλοῖς πολλάκις σπουδαῖον ὄντα, νυνὶ δὲ πολὺ σπουδαιότερον πεποιθήσει πολλῇ τῇ εἰς ὑμᾶς.
\vs{23}~εἴτε ὑπὲρ Τίτου, κοινωνὸς ἐμὸς καὶ εἰς ὑμᾶς συνεργός· εἴτε \adeA{ἀδελφοὶ} ἡμῶν, ἀπόστολοι ἐκκλησιῶν, δόξα Χριστοῦ.
\vs{24}~τὴν οὖν ἔνδειξιν τῆς ἀγάπης ὑμῶν καὶ ἡμῶν καυχήσεως ὑπὲρ ὑμῶν εἰς αὐτοὺς ἐνδεικνύμενοι εἰς πρόσωπον τῶν ἐκκλησιῶν.

\chapheader{9}
\vs{1}~Περὶ μὲν γὰρ τῆς διακονίας τῆς εἰς τοὺς ἁγίους περισσόν μοί ἐστιν τὸ γράφειν ὑμῖν,
\vs{2}~οἶδα γὰρ τὴν προθυμίαν ὑμῶν ἣν ὑπὲρ ὑμῶν καυχῶμαι Μακεδόσιν ὅτι Ἀχαΐα παρεσκεύασται ἀπὸ πέρυσι, καὶ τὸ ὑμῶν ζῆλος ἠρέθισε τοὺς πλείονας.
\vs{3}~ἔπεμψα δὲ τοὺς \adeA{ἀδελφούς}, ἵνα μὴ τὸ καύχημα ἡμῶν τὸ ὑπὲρ ὑμῶν κενωθῇ ἐν τῷ μέρει τούτῳ, ἵνα καθὼς ἔλεγον παρεσκευασμένοι ἦτε,
\vs{4}~μή πως ἐὰν ἔλθωσιν σὺν ἐμοὶ Μακεδόνες καὶ εὕρωσιν ὑμᾶς ἀπαρασκευάστους καταισχυνθῶμεν ἡμεῖς, ἵνα μὴ λέγωμεν ὑμεῖς, ἐν τῇ ὑποστάσει ταύτῃ.
\vs{5}~ἀναγκαῖον οὖν ἡγησάμην παρακαλέσαι τοὺς \adeA{ἀδελφοὺς} ἵνα προέλθωσιν εἰς ὑμᾶς καὶ προκαταρτίσωσι τὴν προεπηγγελμένην εὐλογίαν ὑμῶν, ταύτην ἑτοίμην εἶναι οὕτως ὡς εὐλογίαν καὶ μὴ ὡς πλεονεξίαν.
\vs{6}~Τοῦτο δέ, ὁ σπείρων φειδομένως φειδομένως καὶ θερίσει, καὶ ὁ σπείρων ἐπ’ εὐλογίαις ἐπ’ εὐλογίαις καὶ θερίσει.
\vs{7}~ἕκαστος καθὼς προῄρηται τῇ καρδίᾳ, μὴ ἐκ λύπης ἢ ἐξ ἀνάγκης· ἱλαρὸν γὰρ δότην ἀγαπᾷ ὁ θεός.
\vs{8}~δυνατεῖ δὲ ὁ θεὸς πᾶσαν χάριν περισσεῦσαι εἰς ὑμᾶς, ἵνα ἐν παντὶ πάντοτε πᾶσαν αὐτάρκειαν ἔχοντες περισσεύητε εἰς πᾶν ἔργον ἀγαθόν·
\vs{9}~(καθὼς γέγραπται· Ἐσκόρπισεν, ἔδωκεν τοῖς πένησιν, ἡ δικαιοσύνη αὐτοῦ μένει εἰς τὸν αἰῶνα·
\vs{10}~ὁ δὲ ἐπιχορηγῶν σπόρον τῷ σπείροντι καὶ ἄρτον εἰς βρῶσιν χορηγήσει καὶ πληθυνεῖ τὸν σπόρον ὑμῶν καὶ αὐξήσει τὰ γενήματα τῆς δικαιοσύνης ὑμῶν·)
\vs{11}~ἐν παντὶ πλουτιζόμενοι εἰς πᾶσαν ἁπλότητα, ἥτις κατεργάζεται δι’ ἡμῶν εὐχαριστίαν τῷ θεῷ—
\vs{12}~ὅτι ἡ διακονία τῆς λειτουργίας ταύτης οὐ μόνον ἐστὶν προσαναπληροῦσα τὰ ὑστερήματα τῶν ἁγίων, ἀλλὰ καὶ περισσεύουσα διὰ πολλῶν εὐχαριστιῶν τῷ θεῷ—
\vs{13}~διὰ τῆς δοκιμῆς τῆς διακονίας ταύτης δοξάζοντες τὸν θεὸν ἐπὶ τῇ ὑποταγῇ τῆς ὁμολογίας ὑμῶν εἰς τὸ εὐαγγέλιον τοῦ Χριστοῦ καὶ ἁπλότητι τῆς κοινωνίας εἰς αὐτοὺς καὶ εἰς πάντας,
\vs{14}~καὶ αὐτῶν δεήσει ὑπὲρ ὑμῶν ἐπιποθούντων ὑμᾶς διὰ τὴν ὑπερβάλλουσαν χάριν τοῦ θεοῦ ἐφ’ ὑμῖν.
\vs{15}~χάρις τῷ θεῷ ἐπὶ τῇ ἀνεκδιηγήτῳ αὐτοῦ δωρεᾷ.

\chapheader{10}
\vs{1}~Αὐτὸς δὲ ἐγὼ Παῦλος παρακαλῶ ὑμᾶς διὰ τῆς πραΰτητος καὶ ἐπιεικείας τοῦ Χριστοῦ, ὃς κατὰ πρόσωπον μὲν ταπεινὸς ἐν ὑμῖν, ἀπὼν δὲ θαρρῶ εἰς ὑμᾶς·
\vs{2}~δέομαι δὲ τὸ μὴ παρὼν θαρρῆσαι τῇ πεποιθήσει ᾗ λογίζομαι τολμῆσαι ἐπί τινας τοὺς λογιζομένους ἡμᾶς ὡς κατὰ σάρκα περιπατοῦντας.
\vs{3}~ἐν σαρκὶ γὰρ περιπατοῦντες οὐ κατὰ σάρκα στρατευόμεθα—
\vs{4}~τὰ γὰρ ὅπλα τῆς στρατείας ἡμῶν οὐ σαρκικὰ ἀλλὰ δυνατὰ τῷ θεῷ πρὸς καθαίρεσιν ὀχυρωμάτων— λογισμοὺς καθαιροῦντες
\vs{5}~καὶ πᾶν ὕψωμα ἐπαιρόμενον κατὰ τῆς γνώσεως τοῦ θεοῦ, καὶ αἰχμαλωτίζοντες πᾶν νόημα εἰς τὴν ὑπακοὴν τοῦ Χριστοῦ,
\vs{6}~καὶ ἐν ἑτοίμῳ ἔχοντες ἐκδικῆσαι πᾶσαν παρακοήν, ὅταν πληρωθῇ ὑμῶν ἡ ὑπακοή.
\vs{7}~Τὰ κατὰ πρόσωπον βλέπετε. εἴ τις πέποιθεν ἑαυτῷ Χριστοῦ εἶναι, τοῦτο λογιζέσθω πάλιν ἐφ’ ἑαυτοῦ ὅτι καθὼς αὐτὸς Χριστοῦ οὕτως καὶ ἡμεῖς.
\vs{8}~ἐάν τε γὰρ περισσότερόν τι καυχήσωμαι περὶ τῆς ἐξουσίας ἡμῶν, ἧς ἔδωκεν ὁ κύριος εἰς οἰκοδομὴν καὶ οὐκ εἰς καθαίρεσιν ὑμῶν, οὐκ αἰσχυνθήσομαι,
\vs{9}~ἵνα μὴ δόξω ὡς ἂν ἐκφοβεῖν ὑμᾶς διὰ τῶν ἐπιστολῶν·
\vs{10}~ὅτι Αἱ ἐπιστολαὶ μέν, φησίν, βαρεῖαι καὶ ἰσχυραί, ἡ δὲ παρουσία τοῦ σώματος ἀσθενὴς καὶ ὁ λόγος ἐξουθενημένος.
\vs{11}~τοῦτο λογιζέσθω ὁ τοιοῦτος, ὅτι οἷοί ἐσμεν τῷ λόγῳ δι’ ἐπιστολῶν ἀπόντες, τοιοῦτοι καὶ παρόντες τῷ ἔργῳ.
\vs{12}~Οὐ γὰρ τολμῶμεν ἐγκρῖναι ἢ συγκρῖναι ἑαυτούς τισιν τῶν ἑαυτοὺς συνιστανόντων· ἀλλὰ αὐτοὶ ἐν ἑαυτοῖς ἑαυτοὺς μετροῦντες καὶ συγκρίνοντες ἑαυτοὺς ἑαυτοῖς οὐ συνιᾶσιν.
\vs{13}~ἡμεῖς δὲ οὐκ εἰς τὰ ἄμετρα καυχησόμεθα, ἀλλὰ κατὰ τὸ μέτρον τοῦ κανόνος οὗ ἐμέρισεν ἡμῖν ὁ θεὸς μέτρου, ἐφικέσθαι ἄχρι καὶ ὑμῶν—
\vs{14}~οὐ γὰρ ὡς μὴ ἐφικνούμενοι εἰς ὑμᾶς ὑπερεκτείνομεν ἑαυτούς, ἄχρι γὰρ καὶ ὑμῶν ἐφθάσαμεν ἐν τῷ εὐαγγελίῳ τοῦ Χριστοῦ—
\vs{15}~οὐκ εἰς τὰ ἄμετρα καυχώμενοι ἐν ἀλλοτρίοις κόποις, ἐλπίδα δὲ ἔχοντες αὐξανομένης τῆς πίστεως ὑμῶν ἐν ὑμῖν μεγαλυνθῆναι κατὰ τὸν κανόνα ἡμῶν εἰς περισσείαν,
\vs{16}~εἰς τὰ ὑπερέκεινα ὑμῶν εὐαγγελίσασθαι, οὐκ ἐν ἀλλοτρίῳ κανόνι εἰς τὰ ἕτοιμα καυχήσασθαι.
\vs{17}~Ὁ δὲ καυχώμενος ἐν κυρίῳ καυχάσθω·
\vs{18}~οὐ γὰρ ὁ ἑαυτὸν συνιστάνων, ἐκεῖνός ἐστιν δόκιμος, ἀλλὰ ὃν ὁ κύριος συνίστησιν.

\chapheader{11}
\vs{1}~Ὄφελον ἀνείχεσθέ μου μικρόν τι ἀφροσύνης· ἀλλὰ καὶ ἀνέχεσθέ μου.
\vs{2}~ζηλῶ γὰρ ὑμᾶς θεοῦ ζήλῳ, ἡρμοσάμην γὰρ ὑμᾶς ἑνὶ ἀνδρὶ παρθένον ἁγνὴν παραστῆσαι τῷ Χριστῷ·
\vs{3}~φοβοῦμαι δὲ μή πως, ὡς ὁ ὄφις ἐξηπάτησεν Εὕαν ἐν τῇ πανουργίᾳ αὐτοῦ, φθαρῇ τὰ νοήματα ὑμῶν ἀπὸ τῆς ἁπλότητος καὶ τῆς ἁγνότητος τῆς εἰς τὸν Χριστόν.
\vs{4}~εἰ μὲν γὰρ ὁ ἐρχόμενος ἄλλον Ἰησοῦν κηρύσσει ὃν οὐκ ἐκηρύξαμεν, ἢ πνεῦμα ἕτερον λαμβάνετε ὃ οὐκ ἐλάβετε, ἢ εὐαγγέλιον ἕτερον ὃ οὐκ ἐδέξασθε, καλῶς ἀνέχεσθε.
\vs{5}~λογίζομαι γὰρ μηδὲν ὑστερηκέναι τῶν ὑπερλίαν ἀποστόλων·
\vs{6}~εἰ δὲ καὶ ἰδιώτης τῷ λόγῳ, ἀλλ’ οὐ τῇ γνώσει, ἀλλ’ ἐν παντὶ φανερώσαντες ἐν πᾶσιν εἰς ὑμᾶς.
\vs{7}~Ἢ ἁμαρτίαν ἐποίησα ἐμαυτὸν ταπεινῶν ἵνα ὑμεῖς ὑψωθῆτε, ὅτι δωρεὰν τὸ τοῦ θεοῦ εὐαγγέλιον εὐηγγελισάμην ὑμῖν;
\vs{8}~ἄλλας ἐκκλησίας ἐσύλησα λαβὼν ὀψώνιον πρὸς τὴν ὑμῶν διακονίαν,
\vs{9}~καὶ παρὼν πρὸς ὑμᾶς καὶ ὑστερηθεὶς οὐ κατενάρκησα οὐθενός· τὸ γὰρ ὑστέρημά μου προσανεπλήρωσαν οἱ \adeA{ἀδελφοὶ} ἐλθόντες ἀπὸ Μακεδονίας· καὶ ἐν παντὶ ἀβαρῆ ἐμαυτὸν ὑμῖν ἐτήρησα καὶ τηρήσω.
\vs{10}~ἔστιν ἀλήθεια Χριστοῦ ἐν ἐμοὶ ὅτι ἡ καύχησις αὕτη οὐ φραγήσεται εἰς ἐμὲ ἐν τοῖς κλίμασι τῆς Ἀχαΐας.
\vs{11}~διὰ τί; ὅτι οὐκ ἀγαπῶ ὑμᾶς; ὁ θεὸς οἶδεν.
\vs{12}~Ὃ δὲ ποιῶ καὶ ποιήσω, ἵνα ἐκκόψω τὴν ἀφορμὴν τῶν θελόντων ἀφορμήν, ἵνα ἐν ᾧ καυχῶνται εὑρεθῶσιν καθὼς καὶ ἡμεῖς.
\vs{13}~οἱ γὰρ τοιοῦτοι ψευδαπόστολοι, ἐργάται δόλιοι, μετασχηματιζόμενοι εἰς ἀποστόλους Χριστοῦ·
\vs{14}~καὶ οὐ θαῦμα, αὐτὸς γὰρ ὁ Σατανᾶς μετασχηματίζεται εἰς ἄγγελον φωτός·
\vs{15}~οὐ μέγα οὖν εἰ καὶ οἱ διάκονοι αὐτοῦ μετασχηματίζονται ὡς διάκονοι δικαιοσύνης, ὧν τὸ τέλος ἔσται κατὰ τὰ ἔργα αὐτῶν.
\vs{16}~Πάλιν λέγω, μή τίς με δόξῃ ἄφρονα εἶναι— εἰ δὲ μή γε, κἂν ὡς ἄφρονα δέξασθέ με, ἵνα κἀγὼ μικρόν τι καυχήσωμαι·
\vs{17}~ὃ λαλῶ οὐ κατὰ κύριον λαλῶ, ἀλλ’ ὡς ἐν ἀφροσύνῃ, ἐν ταύτῃ τῇ ὑποστάσει τῆς καυχήσεως.
\vs{18}~ἐπεὶ πολλοὶ καυχῶνται κατὰ σάρκα, κἀγὼ καυχήσομαι.
\vs{19}~ἡδέως γὰρ ἀνέχεσθε τῶν ἀφρόνων φρόνιμοι ὄντες·
\vs{20}~ἀνέχεσθε γὰρ εἴ τις ὑμᾶς καταδουλοῖ, εἴ τις κατεσθίει, εἴ τις λαμβάνει, εἴ τις ἐπαίρεται, εἴ τις εἰς πρόσωπον ὑμᾶς δέρει.
\vs{21}~κατὰ ἀτιμίαν λέγω, ὡς ὅτι ἡμεῖς ἠσθενήκαμεν· ἐν ᾧ δ’ ἄν τις τολμᾷ, ἐν ἀφροσύνῃ λέγω, τολμῶ κἀγώ.
\vs{22}~Ἑβραῖοί εἰσιν; κἀγώ. Ἰσραηλῖταί εἰσιν; κἀγώ. σπέρμα Ἀβραάμ εἰσιν; κἀγώ.
\vs{23}~διάκονοι Χριστοῦ εἰσιν; παραφρονῶν λαλῶ, ὑπὲρ ἐγώ· ἐν κόποις περισσοτέρως, ἐν φυλακαῖς περισσοτέρως, ἐν πληγαῖς ὑπερβαλλόντως, ἐν θανάτοις πολλάκις·
\vs{24}~ὑπὸ Ἰουδαίων πεντάκις τεσσεράκοντα παρὰ μίαν ἔλαβον,
\vs{25}~τρὶς ἐραβδίσθην, ἅπαξ ἐλιθάσθην, τρὶς ἐναυάγησα, νυχθήμερον ἐν τῷ βυθῷ πεποίηκα·
\vs{26}~ὁδοιπορίαις πολλάκις, κινδύνοις ποταμῶν, κινδύνοις λῃστῶν, κινδύνοις ἐκ γένους, κινδύνοις ἐξ ἐθνῶν, κινδύνοις ἐν πόλει, κινδύνοις ἐν ἐρημίᾳ, κινδύνοις ἐν θαλάσσῃ, κινδύνοις ἐν \adeA{ψευδαδέλφοις},
\vs{27}~κόπῳ καὶ μόχθῳ, ἐν ἀγρυπνίαις πολλάκις, ἐν λιμῷ καὶ δίψει, ἐν νηστείαις πολλάκις, ἐν ψύχει καὶ γυμνότητι·
\vs{28}~χωρὶς τῶν παρεκτὸς ἡ ἐπίστασίς μοι ἡ καθ’ ἡμέραν, ἡ μέριμνα πασῶν τῶν ἐκκλησιῶν.
\vs{29}~τίς ἀσθενεῖ, καὶ οὐκ ἀσθενῶ; τίς σκανδαλίζεται καὶ οὐκ ἐγὼ πυροῦμαι;
\vs{30}~Εἰ καυχᾶσθαι δεῖ, τὰ τῆς ἀσθενείας μου καυχήσομαι.
\vs{31}~ὁ θεὸς καὶ πατὴρ τοῦ κυρίου Ἰησοῦ οἶδεν, ὁ ὢν εὐλογητὸς εἰς τοὺς αἰῶνας, ὅτι οὐ ψεύδομαι.
\vs{32}~ἐν Δαμασκῷ ὁ ἐθνάρχης Ἁρέτα τοῦ βασιλέως ἐφρούρει τὴν πόλιν Δαμασκηνῶν πιάσαι με,
\vs{33}~καὶ διὰ θυρίδος ἐν σαργάνῃ ἐχαλάσθην διὰ τοῦ τείχους καὶ ἐξέφυγον τὰς χεῖρας αὐτοῦ.

\chapheader{12}
\vs{1}~Καυχᾶσθαι δεῖ· οὐ συμφέρον μέν, ἐλεύσομαι δὲ εἰς ὀπτασίας καὶ ἀποκαλύψεις κυρίου.
\vs{2}~οἶδα ἄνθρωπον ἐν Χριστῷ πρὸ ἐτῶν δεκατεσσάρων— εἴτε ἐν σώματι οὐκ οἶδα, εἴτε ἐκτὸς τοῦ σώματος οὐκ οἶδα, ὁ θεὸς οἶδεν— ἁρπαγέντα τὸν τοιοῦτον ἕως τρίτου οὐρανοῦ.
\vs{3}~καὶ οἶδα τὸν τοιοῦτον ἄνθρωπον— εἴτε ἐν σώματι εἴτε χωρὶς τοῦ σώματος οὐκ οἶδα, ὁ θεὸς οἶδεν—
\vs{4}~ὅτι ἡρπάγη εἰς τὸν παράδεισον καὶ ἤκουσεν ἄρρητα ῥήματα ἃ οὐκ ἐξὸν ἀνθρώπῳ λαλῆσαι.
\vs{5}~ὑπὲρ τοῦ τοιούτου καυχήσομαι, ὑπὲρ δὲ ἐμαυτοῦ οὐ καυχήσομαι εἰ μὴ ἐν ταῖς ἀσθενείαις.
\vs{6}~ἐὰν γὰρ θελήσω καυχήσασθαι, οὐκ ἔσομαι ἄφρων, ἀλήθειαν γὰρ ἐρῶ· φείδομαι δέ, μή τις εἰς ἐμὲ λογίσηται ὑπὲρ ὃ βλέπει με ἢ ἀκούει τι ἐξ ἐμοῦ,
\vs{7}~καὶ τῇ ὑπερβολῇ τῶν ἀποκαλύψεων. διὸ ἵνα μὴ ὑπεραίρωμαι, ἐδόθη μοι σκόλοψ τῇ σαρκί, ἄγγελος Σατανᾶ, ἵνα με κολαφίζῃ, ἵνα μὴ ὑπεραίρωμαι.
\vs{8}~ὑπὲρ τούτου τρὶς τὸν κύριον παρεκάλεσα ἵνα ἀποστῇ ἀπ’ ἐμοῦ·
\vs{9}~καὶ εἴρηκέν μοι· Ἀρκεῖ σοι ἡ χάρις μου· ἡ γὰρ δύναμις ἐν ἀσθενείᾳ τελεῖται. ἥδιστα οὖν μᾶλλον καυχήσομαι ἐν ταῖς ἀσθενείαις μου, ἵνα ἐπισκηνώσῃ ἐπ’ ἐμὲ ἡ δύναμις τοῦ Χριστοῦ.
\vs{10}~διὸ εὐδοκῶ ἐν ἀσθενείαις, ἐν ὕβρεσιν, ἐν ἀνάγκαις, ἐν διωγμοῖς καὶ στενοχωρίαις, ὑπὲρ Χριστοῦ· ὅταν γὰρ ἀσθενῶ, τότε δυνατός εἰμι.
\vs{11}~Γέγονα ἄφρων· ὑμεῖς με ἠναγκάσατε· ἐγὼ γὰρ ὤφειλον ὑφ’ ὑμῶν συνίστασθαι. οὐδὲν γὰρ ὑστέρησα τῶν ὑπερλίαν ἀποστόλων, εἰ καὶ οὐδέν εἰμι·
\vs{12}~τὰ μὲν σημεῖα τοῦ ἀποστόλου κατειργάσθη ἐν ὑμῖν ἐν πάσῃ ὑπομονῇ, σημείοις τε καὶ τέρασιν καὶ δυνάμεσιν.
\vs{13}~τί γάρ ἐστιν ὃ ἡσσώθητε ὑπὲρ τὰς λοιπὰς ἐκκλησίας, εἰ μὴ ὅτι αὐτὸς ἐγὼ οὐ κατενάρκησα ὑμῶν; χαρίσασθέ μοι τὴν ἀδικίαν ταύτην.
\vs{14}~Ἰδοὺ τρίτον τοῦτο ἑτοίμως ἔχω ἐλθεῖν πρὸς ὑμᾶς, καὶ οὐ καταναρκήσω· οὐ γὰρ ζητῶ τὰ ὑμῶν ἀλλὰ ὑμᾶς, οὐ γὰρ ὀφείλει τὰ τέκνα τοῖς γονεῦσιν θησαυρίζειν, ἀλλὰ οἱ γονεῖς τοῖς τέκνοις.
\vs{15}~ἐγὼ δὲ ἥδιστα δαπανήσω καὶ ἐκδαπανηθήσομαι ὑπὲρ τῶν ψυχῶν ὑμῶν. εἰ περισσοτέρως ὑμᾶς ἀγαπῶν, ἧσσον ἀγαπῶμαι;
\vs{16}~ἔστω δέ, ἐγὼ οὐ κατεβάρησα ὑμᾶς· ἀλλὰ ὑπάρχων πανοῦργος δόλῳ ὑμᾶς ἔλαβον.
\vs{17}~μή τινα ὧν ἀπέσταλκα πρὸς ὑμᾶς, δι’ αὐτοῦ ἐπλεονέκτησα ὑμᾶς;
\vs{18}~παρεκάλεσα Τίτον καὶ συναπέστειλα τὸν \adeA{ἀδελφόν}· μήτι ἐπλεονέκτησεν ὑμᾶς Τίτος; οὐ τῷ αὐτῷ πνεύματι περιεπατήσαμεν; οὐ τοῖς αὐτοῖς ἴχνεσιν;
\vs{19}~Πάλαι δοκεῖτε ὅτι ὑμῖν ἀπολογούμεθα; κατέναντι θεοῦ ἐν Χριστῷ λαλοῦμεν. τὰ δὲ πάντα, ἀγαπητοί, ὑπὲρ τῆς ὑμῶν οἰκοδομῆς,
\vs{20}~φοβοῦμαι γὰρ μή πως ἐλθὼν οὐχ οἵους θέλω εὕρω ὑμᾶς, κἀγὼ εὑρεθῶ ὑμῖν οἷον οὐ θέλετε, μή πως ἔρις, ζῆλος, θυμοί, ἐριθεῖαι, καταλαλιαί, ψιθυρισμοί, φυσιώσεις, ἀκαταστασίαι·
\vs{21}~μὴ πάλιν ἐλθόντος μου ταπεινώσῃ με ὁ θεός μου πρὸς ὑμᾶς, καὶ πενθήσω πολλοὺς τῶν προημαρτηκότων καὶ μὴ μετανοησάντων ἐπὶ τῇ ἀκαθαρσίᾳ καὶ πορνείᾳ καὶ ἀσελγείᾳ ᾗ ἔπραξαν.

\chapheader{13}
\vs{1}~Τρίτον τοῦτο ἔρχομαι πρὸς ὑμᾶς· ἐπὶ στόματος δύο μαρτύρων καὶ τριῶν σταθήσεται πᾶν ῥῆμα.
\vs{2}~προείρηκα καὶ προλέγω ὡς παρὼν τὸ δεύτερον καὶ ἀπὼν νῦν, τοῖς προημαρτηκόσιν καὶ τοῖς λοιποῖς πᾶσιν, ὅτι ἐὰν ἔλθω εἰς τὸ πάλιν οὐ φείσομαι,
\vs{3}~ἐπεὶ δοκιμὴν ζητεῖτε τοῦ ἐν ἐμοὶ λαλοῦντος Χριστοῦ· ὃς εἰς ὑμᾶς οὐκ ἀσθενεῖ ἀλλὰ δυνατεῖ ἐν ὑμῖν,
\vs{4}~καὶ γὰρ ἐσταυρώθη ἐξ ἀσθενείας, ἀλλὰ ζῇ ἐκ δυνάμεως θεοῦ. καὶ γὰρ ἡμεῖς ἀσθενοῦμεν ἐν αὐτῷ, ἀλλὰ ζήσομεν σὺν αὐτῷ ἐκ δυνάμεως θεοῦ εἰς ὑμᾶς.
\vs{5}~Ἑαυτοὺς πειράζετε εἰ ἐστὲ ἐν τῇ πίστει, ἑαυτοὺς δοκιμάζετε· ἢ οὐκ ἐπιγινώσκετε ἑαυτοὺς ὅτι Ἰησοῦς Χριστὸς ἐν ὑμῖν; εἰ μήτι ἀδόκιμοί ἐστε.
\vs{6}~ἐλπίζω δὲ ὅτι γνώσεσθε ὅτι ἡμεῖς οὐκ ἐσμὲν ἀδόκιμοι.
\vs{7}~εὐχόμεθα δὲ πρὸς τὸν θεὸν μὴ ποιῆσαι ὑμᾶς κακὸν μηδέν, οὐχ ἵνα ἡμεῖς δόκιμοι φανῶμεν, ἀλλ’ ἵνα ὑμεῖς τὸ καλὸν ποιῆτε, ἡμεῖς δὲ ὡς ἀδόκιμοι ὦμεν.
\vs{8}~οὐ γὰρ δυνάμεθά τι κατὰ τῆς ἀληθείας, ἀλλὰ ὑπὲρ τῆς ἀληθείας.
\vs{9}~χαίρομεν γὰρ ὅταν ἡμεῖς ἀσθενῶμεν, ὑμεῖς δὲ δυνατοὶ ἦτε· τοῦτο καὶ εὐχόμεθα, τὴν ὑμῶν κατάρτισιν.
\vs{10}~διὰ τοῦτο ταῦτα ἀπὼν γράφω, ἵνα παρὼν μὴ ἀποτόμως χρήσωμαι κατὰ τὴν ἐξουσίαν ἣν ὁ κύριος ἔδωκέν μοι, εἰς οἰκοδομὴν καὶ οὐκ εἰς καθαίρεσιν.
\vs{11}~Λοιπόν, \adeA{ἀδελφοί}, χαίρετε, καταρτίζεσθε, παρακαλεῖσθε, τὸ αὐτὸ φρονεῖτε, εἰρηνεύετε, καὶ ὁ θεὸς τῆς ἀγάπης καὶ εἰρήνης ἔσται μεθ’ ὑμῶν.
\vs{12}~ἀσπάσασθε ἀλλήλους ἐν ἁγίῳ φιλήματι. ἀσπάζονται ὑμᾶς οἱ ἅγιοι πάντες.
\vs{13}~ἡ χάρις τοῦ κυρίου Ἰησοῦ Χριστοῦ καὶ ἡ ἀγάπη τοῦ θεοῦ καὶ ἡ κοινωνία τοῦ ἁγίου πνεύματος μετὰ πάντων ὑμῶν.

\booksection{Galatians}{Πρὸς Γαλάτας}

\chapheader{1}
\vs{1}~Παῦλος ἀπόστολος, οὐκ ἀπ’ ἀνθρώπων οὐδὲ δι’ ἀνθρώπου ἀλλὰ διὰ Ἰησοῦ Χριστοῦ καὶ θεοῦ πατρὸς τοῦ ἐγείραντος αὐτὸν ἐκ νεκρῶν,
\vs{2}~καὶ οἱ σὺν ἐμοὶ πάντες \adeA{ἀδελφοί}, ταῖς ἐκκλησίαις τῆς Γαλατίας·
\vs{3}~χάρις ὑμῖν καὶ εἰρήνη ἀπὸ θεοῦ πατρὸς καὶ κυρίου ἡμῶν Ἰησοῦ Χριστοῦ,
\vs{4}~τοῦ δόντος ἑαυτὸν ὑπὲρ τῶν ἁμαρτιῶν ἡμῶν ὅπως ἐξέληται ἡμᾶς ἐκ τοῦ αἰῶνος τοῦ ἐνεστῶτος πονηροῦ κατὰ τὸ θέλημα τοῦ θεοῦ καὶ πατρὸς ἡμῶν,
\vs{5}~ᾧ ἡ δόξα εἰς τοὺς αἰῶνας τῶν αἰώνων· ἀμήν.
\vs{6}~Θαυμάζω ὅτι οὕτως ταχέως μετατίθεσθε ἀπὸ τοῦ καλέσαντος ὑμᾶς ἐν χάριτι Χριστοῦ εἰς ἕτερον εὐαγγέλιον,
\vs{7}~ὃ οὐκ ἔστιν ἄλλο· εἰ μή τινές εἰσιν οἱ ταράσσοντες ὑμᾶς καὶ θέλοντες μεταστρέψαι τὸ εὐαγγέλιον τοῦ Χριστοῦ.
\vs{8}~ἀλλὰ καὶ ἐὰν ἡμεῖς ἢ ἄγγελος ἐξ οὐρανοῦ εὐαγγελίζηται ὑμῖν παρ’ ὃ εὐηγγελισάμεθα ὑμῖν, ἀνάθεμα ἔστω.
\vs{9}~ὡς προειρήκαμεν, καὶ ἄρτι πάλιν λέγω, εἴ τις ὑμᾶς εὐαγγελίζεται παρ’ ὃ παρελάβετε, ἀνάθεμα ἔστω.
\vs{10}~Ἄρτι γὰρ ἀνθρώπους πείθω ἢ τὸν θεόν; ἢ ζητῶ ἀνθρώποις ἀρέσκειν; εἰ ἔτι ἀνθρώποις ἤρεσκον, Χριστοῦ δοῦλος οὐκ ἂν ἤμην.
\vs{11}~Γνωρίζω γὰρ ὑμῖν, \adeA{ἀδελφοί}, τὸ εὐαγγέλιον τὸ εὐαγγελισθὲν ὑπ’ ἐμοῦ ὅτι οὐκ ἔστιν κατὰ ἄνθρωπον·
\vs{12}~οὐδὲ γὰρ ἐγὼ παρὰ ἀνθρώπου παρέλαβον αὐτό, οὔτε ἐδιδάχθην, ἀλλὰ δι’ ἀποκαλύψεως Ἰησοῦ Χριστοῦ.
\vs{13}~Ἠκούσατε γὰρ τὴν ἐμὴν ἀναστροφήν ποτε ἐν τῷ Ἰουδαϊσμῷ, ὅτι καθ’ ὑπερβολὴν ἐδίωκον τὴν ἐκκλησίαν τοῦ θεοῦ καὶ ἐπόρθουν αὐτήν,
\vs{14}~καὶ προέκοπτον ἐν τῷ Ἰουδαϊσμῷ ὑπὲρ πολλοὺς συνηλικιώτας ἐν τῷ γένει μου, περισσοτέρως ζηλωτὴς ὑπάρχων τῶν πατρικῶν μου παραδόσεων.
\vs{15}~ὅτε δὲ εὐδόκησεν ὁ ἀφορίσας με ἐκ κοιλίας μητρός μου καὶ καλέσας διὰ τῆς χάριτος αὐτοῦ
\vs{16}~ἀποκαλύψαι τὸν υἱὸν αὐτοῦ ἐν ἐμοὶ ἵνα εὐαγγελίζωμαι αὐτὸν ἐν τοῖς ἔθνεσιν, εὐθέως οὐ προσανεθέμην σαρκὶ καὶ αἵματι,
\vs{17}~οὐδὲ ἀνῆλθον εἰς Ἱεροσόλυμα πρὸς τοὺς πρὸ ἐμοῦ ἀποστόλους, ἀλλὰ ἀπῆλθον εἰς Ἀραβίαν, καὶ πάλιν ὑπέστρεψα εἰς Δαμασκόν.
\vs{18}~Ἔπειτα μετὰ ἔτη τρία ἀνῆλθον εἰς Ἱεροσόλυμα ἱστορῆσαι Κηφᾶν, καὶ ἐπέμεινα πρὸς αὐτὸν ἡμέρας δεκαπέντε·
\vs{19}~ἕτερον δὲ τῶν ἀποστόλων οὐκ εἶδον, εἰ μὴ Ἰάκωβον τὸν \adeC{ἀδελφὸν} τοῦ κυρίου.
\vs{20}~ἃ δὲ γράφω ὑμῖν, ἰδοὺ ἐνώπιον τοῦ θεοῦ ὅτι οὐ ψεύδομαι.
\vs{21}~ἔπειτα ἦλθον εἰς τὰ κλίματα τῆς Συρίας καὶ τῆς Κιλικίας.
\vs{22}~ἤμην δὲ ἀγνοούμενος τῷ προσώπῳ ταῖς ἐκκλησίαις τῆς Ἰουδαίας ταῖς ἐν Χριστῷ,
\vs{23}~μόνον δὲ ἀκούοντες ἦσαν ὅτι Ὁ διώκων ἡμᾶς ποτε νῦν εὐαγγελίζεται τὴν πίστιν ἥν ποτε ἐπόρθει,
\vs{24}~καὶ ἐδόξαζον ἐν ἐμοὶ τὸν θεόν.

\chapheader{2}
\vs{1}~Ἔπειτα διὰ δεκατεσσάρων ἐτῶν πάλιν ἀνέβην εἰς Ἱεροσόλυμα μετὰ Βαρναβᾶ συμπαραλαβὼν καὶ Τίτον·
\vs{2}~ἀνέβην δὲ κατὰ ἀποκάλυψιν· καὶ ἀνεθέμην αὐτοῖς τὸ εὐαγγέλιον ὃ κηρύσσω ἐν τοῖς ἔθνεσιν, κατ’ ἰδίαν δὲ τοῖς δοκοῦσιν, μή πως εἰς κενὸν τρέχω ἢ ἔδραμον.
\vs{3}~ἀλλ’ οὐδὲ Τίτος ὁ σὺν ἐμοί, Ἕλλην ὤν, ἠναγκάσθη περιτμηθῆναι·
\vs{4}~διὰ δὲ τοὺς παρεισάκτους \adeA{ψευδαδέλφους}, οἵτινες παρεισῆλθον κατασκοπῆσαι τὴν ἐλευθερίαν ἡμῶν ἣν ἔχομεν ἐν Χριστῷ Ἰησοῦ, ἵνα ἡμᾶς καταδουλώσουσιν—
\vs{5}~οἷς οὐδὲ πρὸς ὥραν εἴξαμεν τῇ ὑποταγῇ, ἵνα ἡ ἀλήθεια τοῦ εὐαγγελίου διαμείνῃ πρὸς ὑμᾶς.
\vs{6}~ἀπὸ δὲ τῶν δοκούντων εἶναί τι— ὁποῖοί ποτε ἦσαν οὐδέν μοι διαφέρει· πρόσωπον θεὸς ἀνθρώπου οὐ λαμβάνει— ἐμοὶ γὰρ οἱ δοκοῦντες οὐδὲν προσανέθεντο,
\vs{7}~ἀλλὰ τοὐναντίον ἰδόντες ὅτι πεπίστευμαι τὸ εὐαγγέλιον τῆς ἀκροβυστίας καθὼς Πέτρος τῆς περιτομῆς,
\vs{8}~ὁ γὰρ ἐνεργήσας Πέτρῳ εἰς ἀποστολὴν τῆς περιτομῆς ἐνήργησεν καὶ ἐμοὶ εἰς τὰ ἔθνη,
\vs{9}~καὶ γνόντες τὴν χάριν τὴν δοθεῖσάν μοι, Ἰάκωβος καὶ Κηφᾶς καὶ Ἰωάννης, οἱ δοκοῦντες στῦλοι εἶναι, δεξιὰς ἔδωκαν ἐμοὶ καὶ Βαρναβᾷ κοινωνίας, ἵνα ἡμεῖς εἰς τὰ ἔθνη, αὐτοὶ δὲ εἰς τὴν περιτομήν·
\vs{10}~μόνον τῶν πτωχῶν ἵνα μνημονεύωμεν, ὃ καὶ ἐσπούδασα αὐτὸ τοῦτο ποιῆσαι.
\vs{11}~Ὅτε δὲ ἦλθεν Κηφᾶς εἰς Ἀντιόχειαν, κατὰ πρόσωπον αὐτῷ ἀντέστην, ὅτι κατεγνωσμένος ἦν·
\vs{12}~πρὸ τοῦ γὰρ ἐλθεῖν τινας ἀπὸ Ἰακώβου μετὰ τῶν ἐθνῶν συνήσθιεν· ὅτε δὲ ἦλθον, ὑπέστελλεν καὶ ἀφώριζεν ἑαυτόν, φοβούμενος τοὺς ἐκ περιτομῆς.
\vs{13}~καὶ συνυπεκρίθησαν αὐτῷ καὶ οἱ λοιποὶ Ἰουδαῖοι, ὥστε καὶ Βαρναβᾶς συναπήχθη αὐτῶν τῇ ὑποκρίσει.
\vs{14}~ἀλλ’ ὅτε εἶδον ὅτι οὐκ ὀρθοποδοῦσιν πρὸς τὴν ἀλήθειαν τοῦ εὐαγγελίου, εἶπον τῷ Κηφᾷ ἔμπροσθεν πάντων· Εἰ σὺ Ἰουδαῖος ὑπάρχων ἐθνικῶς καὶ οὐκ Ἰουδαϊκῶς ζῇς, πῶς τὰ ἔθνη ἀναγκάζεις Ἰουδαΐζειν;
\vs{15}~Ἡμεῖς φύσει Ἰουδαῖοι καὶ οὐκ ἐξ ἐθνῶν ἁμαρτωλοί,
\vs{16}~εἰδότες δὲ ὅτι οὐ δικαιοῦται ἄνθρωπος ἐξ ἔργων νόμου ἐὰν μὴ διὰ πίστεως Ἰησοῦ Χριστοῦ, καὶ ἡμεῖς εἰς Χριστὸν Ἰησοῦν ἐπιστεύσαμεν, ἵνα δικαιωθῶμεν ἐκ πίστεως Χριστοῦ καὶ οὐκ ἐξ ἔργων νόμου, ὅτι ἐξ ἔργων νόμου οὐ δικαιωθήσεται πᾶσα σάρξ.
\vs{17}~εἰ δὲ ζητοῦντες δικαιωθῆναι ἐν Χριστῷ εὑρέθημεν καὶ αὐτοὶ ἁμαρτωλοί, ἆρα Χριστὸς ἁμαρτίας διάκονος; μὴ γένοιτο·
\vs{18}~εἰ γὰρ ἃ κατέλυσα ταῦτα πάλιν οἰκοδομῶ, παραβάτην ἐμαυτὸν συνιστάνω.
\vs{19}~ἐγὼ γὰρ διὰ νόμου νόμῳ ἀπέθανον ἵνα θεῷ ζήσω· Χριστῷ συνεσταύρωμαι·
\vs{20}~ζῶ δὲ οὐκέτι ἐγώ, ζῇ δὲ ἐν ἐμοὶ Χριστός· ὃ δὲ νῦν ζῶ ἐν σαρκί, ἐν πίστει ζῶ τῇ τοῦ υἱοῦ τοῦ θεοῦ τοῦ ἀγαπήσαντός με καὶ παραδόντος ἑαυτὸν ὑπὲρ ἐμοῦ.
\vs{21}~οὐκ ἀθετῶ τὴν χάριν τοῦ θεοῦ· εἰ γὰρ διὰ νόμου δικαιοσύνη, ἄρα Χριστὸς δωρεὰν ἀπέθανεν.

\chapheader{3}
\vs{1}~Ὦ ἀνόητοι Γαλάται, τίς ὑμᾶς ἐβάσκανεν, οἷς κατ’ ὀφθαλμοὺς Ἰησοῦς Χριστὸς προεγράφη ἐσταυρωμένος;
\vs{2}~τοῦτο μόνον θέλω μαθεῖν ἀφ’ ὑμῶν, ἐξ ἔργων νόμου τὸ πνεῦμα ἐλάβετε ἢ ἐξ ἀκοῆς πίστεως;
\vs{3}~οὕτως ἀνόητοί ἐστε; ἐναρξάμενοι πνεύματι νῦν σαρκὶ ἐπιτελεῖσθε;
\vs{4}~τοσαῦτα ἐπάθετε εἰκῇ; εἴ γε καὶ εἰκῇ.
\vs{5}~ὁ οὖν ἐπιχορηγῶν ὑμῖν τὸ πνεῦμα καὶ ἐνεργῶν δυνάμεις ἐν ὑμῖν ἐξ ἔργων νόμου ἢ ἐξ ἀκοῆς πίστεως;
\vs{6}~καθὼς Ἀβραὰμ ἐπίστευσεν τῷ θεῷ, καὶ ἐλογίσθη αὐτῷ εἰς δικαιοσύνην.
\vs{7}~Γινώσκετε ἄρα ὅτι οἱ ἐκ πίστεως, οὗτοι υἱοί εἰσιν Ἀβραάμ.
\vs{8}~προϊδοῦσα δὲ ἡ γραφὴ ὅτι ἐκ πίστεως δικαιοῖ τὰ ἔθνη ὁ θεὸς προευηγγελίσατο τῷ Ἀβραὰμ ὅτι Ἐνευλογηθήσονται ἐν σοὶ πάντα τὰ ἔθνη.
\vs{9}~ὥστε οἱ ἐκ πίστεως εὐλογοῦνται σὺν τῷ πιστῷ Ἀβραάμ.
\vs{10}~Ὅσοι γὰρ ἐξ ἔργων νόμου εἰσὶν ὑπὸ κατάραν εἰσίν, γέγραπται γὰρ ὅτι Ἐπικατάρατος πᾶς ὃς οὐκ ἐμμένει πᾶσιν τοῖς γεγραμμένοις ἐν τῷ βιβλίῳ τοῦ νόμου τοῦ ποιῆσαι αὐτά.
\vs{11}~ὅτι δὲ ἐν νόμῳ οὐδεὶς δικαιοῦται παρὰ τῷ θεῷ δῆλον, ὅτι Ὁ δίκαιος ἐκ πίστεως ζήσεται,
\vs{12}~ὁ δὲ νόμος οὐκ ἔστιν ἐκ πίστεως, ἀλλ’· Ὁ ποιήσας αὐτὰ ζήσεται ἐν αὐτοῖς.
\vs{13}~Χριστὸς ἡμᾶς ἐξηγόρασεν ἐκ τῆς κατάρας τοῦ νόμου γενόμενος ὑπὲρ ἡμῶν κατάρα, ὅτι γέγραπται· Ἐπικατάρατος πᾶς ὁ κρεμάμενος ἐπὶ ξύλου,
\vs{14}~ἵνα εἰς τὰ ἔθνη ἡ εὐλογία τοῦ Ἀβραὰμ γένηται ἐν Χριστῷ Ἰησοῦ, ἵνα τὴν ἐπαγγελίαν τοῦ πνεύματος λάβωμεν διὰ τῆς πίστεως.
\vs{15}~\adeA{Ἀδελφοί}, κατὰ ἄνθρωπον λέγω· ὅμως ἀνθρώπου κεκυρωμένην διαθήκην οὐδεὶς ἀθετεῖ ἢ ἐπιδιατάσσεται.
\vs{16}~τῷ δὲ Ἀβραὰμ ἐρρέθησαν αἱ ἐπαγγελίαι καὶ τῷ σπέρματι αὐτοῦ· οὐ λέγει· Καὶ τοῖς σπέρμασιν, ὡς ἐπὶ πολλῶν, ἀλλ’ ὡς ἐφ’ ἑνός· Καὶ τῷ σπέρματί σου, ὅς ἐστιν Χριστός.
\vs{17}~τοῦτο δὲ λέγω· διαθήκην προκεκυρωμένην ὑπὸ τοῦ θεοῦ ὁ μετὰ τετρακόσια καὶ τριάκοντα ἔτη γεγονὼς νόμος οὐκ ἀκυροῖ, εἰς τὸ καταργῆσαι τὴν ἐπαγγελίαν.
\vs{18}~εἰ γὰρ ἐκ νόμου ἡ κληρονομία, οὐκέτι ἐξ ἐπαγγελίας· τῷ δὲ Ἀβραὰμ δι’ ἐπαγγελίας κεχάρισται ὁ θεός.
\vs{19}~Τί οὖν ὁ νόμος; τῶν παραβάσεων χάριν προσετέθη, ἄχρις οὗ ἔλθῃ τὸ σπέρμα ᾧ ἐπήγγελται, διαταγεὶς δι’ ἀγγέλων ἐν χειρὶ μεσίτου·
\vs{20}~ὁ δὲ μεσίτης ἑνὸς οὐκ ἔστιν, ὁ δὲ θεὸς εἷς ἐστιν.
\vs{21}~Ὁ οὖν νόμος κατὰ τῶν ἐπαγγελιῶν τοῦ θεοῦ; μὴ γένοιτο· εἰ γὰρ ἐδόθη νόμος ὁ δυνάμενος ζῳοποιῆσαι, ὄντως ἐκ νόμου ἂν ἦν ἡ δικαιοσύνη.
\vs{22}~ἀλλὰ συνέκλεισεν ἡ γραφὴ τὰ πάντα ὑπὸ ἁμαρτίαν ἵνα ἡ ἐπαγγελία ἐκ πίστεως Ἰησοῦ Χριστοῦ δοθῇ τοῖς πιστεύουσιν.
\vs{23}~Πρὸ τοῦ δὲ ἐλθεῖν τὴν πίστιν ὑπὸ νόμον ἐφρουρούμεθα συγκλειόμενοι εἰς τὴν μέλλουσαν πίστιν ἀποκαλυφθῆναι.
\vs{24}~ὥστε ὁ νόμος παιδαγωγὸς ἡμῶν γέγονεν εἰς Χριστόν, ἵνα ἐκ πίστεως δικαιωθῶμεν·
\vs{25}~ἐλθούσης δὲ τῆς πίστεως οὐκέτι ὑπὸ παιδαγωγόν ἐσμεν.
\vs{26}~πάντες γὰρ υἱοὶ θεοῦ ἐστε διὰ τῆς πίστεως ἐν Χριστῷ Ἰησοῦ.
\vs{27}~ὅσοι γὰρ εἰς Χριστὸν ἐβαπτίσθητε, Χριστὸν ἐνεδύσασθε·
\vs{28}~οὐκ ἔνι Ἰουδαῖος οὐδὲ Ἕλλην, οὐκ ἔνι δοῦλος οὐδὲ ἐλεύθερος, οὐκ ἔνι ἄρσεν καὶ θῆλυ· πάντες γὰρ ὑμεῖς εἷς ἐστε ἐν Χριστῷ Ἰησοῦ.
\vs{29}~εἰ δὲ ὑμεῖς Χριστοῦ, ἄρα τοῦ Ἀβραὰμ σπέρμα ἐστέ, κατ’ ἐπαγγελίαν κληρονόμοι.

\chapheader{4}
\vs{1}~Λέγω δέ, ἐφ’ ὅσον χρόνον ὁ κληρονόμος νήπιός ἐστιν, οὐδὲν διαφέρει δούλου κύριος πάντων ὤν,
\vs{2}~ἀλλὰ ὑπὸ ἐπιτρόπους ἐστὶ καὶ οἰκονόμους ἄχρι τῆς προθεσμίας τοῦ πατρός.
\vs{3}~οὕτως καὶ ἡμεῖς, ὅτε ἦμεν νήπιοι, ὑπὸ τὰ στοιχεῖα τοῦ κόσμου ἤμεθα δεδουλωμένοι·
\vs{4}~ὅτε δὲ ἦλθεν τὸ πλήρωμα τοῦ χρόνου, ἐξαπέστειλεν ὁ θεὸς τὸν υἱὸν αὐτοῦ, γενόμενον ἐκ γυναικός, γενόμενον ὑπὸ νόμον,
\vs{5}~ἵνα τοὺς ὑπὸ νόμον ἐξαγοράσῃ, ἵνα τὴν υἱοθεσίαν ἀπολάβωμεν.
\vs{6}~ὅτι δέ ἐστε υἱοί, ἐξαπέστειλεν ὁ θεὸς τὸ πνεῦμα τοῦ υἱοῦ αὐτοῦ εἰς τὰς καρδίας ἡμῶν, κρᾶζον· Αββα ὁ πατήρ.
\vs{7}~ὥστε οὐκέτι εἶ δοῦλος ἀλλὰ υἱός· εἰ δὲ υἱός, καὶ κληρονόμος διὰ θεοῦ.
\vs{8}~Ἀλλὰ τότε μὲν οὐκ εἰδότες θεὸν ἐδουλεύσατε τοῖς φύσει μὴ οὖσι θεοῖς·
\vs{9}~νῦν δὲ γνόντες θεόν, μᾶλλον δὲ γνωσθέντες ὑπὸ θεοῦ, πῶς ἐπιστρέφετε πάλιν ἐπὶ τὰ ἀσθενῆ καὶ πτωχὰ στοιχεῖα, οἷς πάλιν ἄνωθεν δουλεύειν θέλετε;
\vs{10}~ἡμέρας παρατηρεῖσθε καὶ μῆνας καὶ καιροὺς καὶ ἐνιαυτούς.
\vs{11}~φοβοῦμαι ὑμᾶς μή πως εἰκῇ κεκοπίακα εἰς ὑμᾶς.
\vs{12}~Γίνεσθε ὡς ἐγώ, ὅτι κἀγὼ ὡς ὑμεῖς, \adeA{ἀδελφοί}, δέομαι ὑμῶν. οὐδέν με ἠδικήσατε·
\vs{13}~οἴδατε δὲ ὅτι δι’ ἀσθένειαν τῆς σαρκὸς εὐηγγελισάμην ὑμῖν τὸ πρότερον,
\vs{14}~καὶ τὸν πειρασμὸν ὑμῶν ἐν τῇ σαρκί μου οὐκ ἐξουθενήσατε οὐδὲ ἐξεπτύσατε, ἀλλὰ ὡς ἄγγελον θεοῦ ἐδέξασθέ με, ὡς Χριστὸν Ἰησοῦν.
\vs{15}~ποῦ οὖν ὁ μακαρισμὸς ὑμῶν; μαρτυρῶ γὰρ ὑμῖν ὅτι εἰ δυνατὸν τοὺς ὀφθαλμοὺς ὑμῶν ἐξορύξαντες ἐδώκατέ μοι.
\vs{16}~ὥστε ἐχθρὸς ὑμῶν γέγονα ἀληθεύων ὑμῖν;
\vs{17}~ζηλοῦσιν ὑμᾶς οὐ καλῶς, ἀλλὰ ἐκκλεῖσαι ὑμᾶς θέλουσιν, ἵνα αὐτοὺς ζηλοῦτε.
\vs{18}~καλὸν δὲ ζηλοῦσθαι ἐν καλῷ πάντοτε, καὶ μὴ μόνον ἐν τῷ παρεῖναί με πρὸς ὑμᾶς,
\vs{19}~τέκνα μου, οὓς πάλιν ὠδίνω μέχρις οὗ μορφωθῇ Χριστὸς ἐν ὑμῖν·
\vs{20}~ἤθελον δὲ παρεῖναι πρὸς ὑμᾶς ἄρτι, καὶ ἀλλάξαι τὴν φωνήν μου, ὅτι ἀποροῦμαι ἐν ὑμῖν.
\vs{21}~Λέγετέ μοι, οἱ ὑπὸ νόμον θέλοντες εἶναι, τὸν νόμον οὐκ ἀκούετε;
\vs{22}~γέγραπται γὰρ ὅτι Ἀβραὰμ δύο υἱοὺς ἔσχεν, ἕνα ἐκ τῆς παιδίσκης καὶ ἕνα ἐκ τῆς ἐλευθέρας·
\vs{23}~ἀλλ’ ὁ μὲν ἐκ τῆς παιδίσκης κατὰ σάρκα γεγέννηται, ὁ δὲ ἐκ τῆς ἐλευθέρας δι’ ἐπαγγελίας.
\vs{24}~ἅτινά ἐστιν ἀλληγορούμενα· αὗται γάρ εἰσιν δύο διαθῆκαι, μία μὲν ἀπὸ ὄρους Σινᾶ, εἰς δουλείαν γεννῶσα, ἥτις ἐστὶν Ἁγάρ,
\vs{25}~τὸ δὲ Ἁγὰρ Σινᾶ ὄρος ἐστὶν ἐν τῇ Ἀραβίᾳ, συστοιχεῖ δὲ τῇ νῦν Ἰερουσαλήμ, δουλεύει γὰρ μετὰ τῶν τέκνων αὐτῆς·
\vs{26}~ἡ δὲ ἄνω Ἰερουσαλὴμ ἐλευθέρα ἐστίν, ἥτις ἐστὶν μήτηρ ἡμῶν·
\vs{27}~γέγραπται γάρ· Εὐφράνθητι, στεῖρα ἡ οὐ τίκτουσα, ῥῆξον καὶ βόησον, ἡ οὐκ ὠδίνουσα· ὅτι πολλὰ τὰ τέκνα τῆς ἐρήμου μᾶλλον ἢ τῆς ἐχούσης τὸν ἄνδρα.
\vs{28}~ὑμεῖς δέ, \adeA{ἀδελφοί}, κατὰ Ἰσαὰκ ἐπαγγελίας τέκνα ἐστέ·
\vs{29}~ἀλλ’ ὥσπερ τότε ὁ κατὰ σάρκα γεννηθεὶς ἐδίωκε τὸν κατὰ πνεῦμα, οὕτως καὶ νῦν.
\vs{30}~ἀλλὰ τί λέγει ἡ γραφή; Ἔκβαλε τὴν παιδίσκην καὶ τὸν υἱὸν αὐτῆς, οὐ γὰρ μὴ κληρονομήσει ὁ υἱὸς τῆς παιδίσκης μετὰ τοῦ υἱοῦ τῆς ἐλευθέρας.
\vs{31}~διό, \adeA{ἀδελφοί}, οὐκ ἐσμὲν παιδίσκης τέκνα ἀλλὰ τῆς ἐλευθέρας.

\chapheader{5}
\vs{1}~τῇ ἐλευθερίᾳ ἡμᾶς Χριστὸς ἠλευθέρωσεν· στήκετε οὖν καὶ μὴ πάλιν ζυγῷ δουλείας ἐνέχεσθε.
\vs{2}~Ἴδε ἐγὼ Παῦλος λέγω ὑμῖν ὅτι ἐὰν περιτέμνησθε Χριστὸς ὑμᾶς οὐδὲν ὠφελήσει.
\vs{3}~μαρτύρομαι δὲ πάλιν παντὶ ἀνθρώπῳ περιτεμνομένῳ ὅτι ὀφειλέτης ἐστὶν ὅλον τὸν νόμον ποιῆσαι.
\vs{4}~κατηργήθητε ἀπὸ Χριστοῦ οἵτινες ἐν νόμῳ δικαιοῦσθε, τῆς χάριτος ἐξεπέσατε.
\vs{5}~ἡμεῖς γὰρ πνεύματι ἐκ πίστεως ἐλπίδα δικαιοσύνης ἀπεκδεχόμεθα.
\vs{6}~ἐν γὰρ Χριστῷ Ἰησοῦ οὔτε περιτομή τι ἰσχύει οὔτε ἀκροβυστία, ἀλλὰ πίστις δι’ ἀγάπης ἐνεργουμένη.
\vs{7}~Ἐτρέχετε καλῶς· τίς ὑμᾶς ἐνέκοψεν τῇ ἀληθείᾳ μὴ πείθεσθαι;
\vs{8}~ἡ πεισμονὴ οὐκ ἐκ τοῦ καλοῦντος ὑμᾶς.
\vs{9}~μικρὰ ζύμη ὅλον τὸ φύραμα ζυμοῖ.
\vs{10}~ἐγὼ πέποιθα εἰς ὑμᾶς ἐν κυρίῳ ὅτι οὐδὲν ἄλλο φρονήσετε· ὁ δὲ ταράσσων ὑμᾶς βαστάσει τὸ κρίμα, ὅστις ἐὰν ᾖ.
\vs{11}~ἐγὼ δέ, \adeA{ἀδελφοί}, εἰ περιτομὴν ἔτι κηρύσσω, τί ἔτι διώκομαι; ἄρα κατήργηται τὸ σκάνδαλον τοῦ σταυροῦ.
\vs{12}~ὄφελον καὶ ἀποκόψονται οἱ ἀναστατοῦντες ὑμᾶς.
\vs{13}~Ὑμεῖς γὰρ ἐπ’ ἐλευθερίᾳ ἐκλήθητε, \adeA{ἀδελφοί}· μόνον μὴ τὴν ἐλευθερίαν εἰς ἀφορμὴν τῇ σαρκί, ἀλλὰ διὰ τῆς ἀγάπης δουλεύετε ἀλλήλοις·
\vs{14}~ὁ γὰρ πᾶς νόμος ἐν ἑνὶ λόγῳ πεπλήρωται, ἐν τῷ· Ἀγαπήσεις τὸν πλησίον σου ὡς σεαυτόν.
\vs{15}~εἰ δὲ ἀλλήλους δάκνετε καὶ κατεσθίετε, βλέπετε μὴ ὑπ’ ἀλλήλων ἀναλωθῆτε.
\vs{16}~Λέγω δέ, πνεύματι περιπατεῖτε καὶ ἐπιθυμίαν σαρκὸς οὐ μὴ τελέσητε.
\vs{17}~ἡ γὰρ σὰρξ ἐπιθυμεῖ κατὰ τοῦ πνεύματος, τὸ δὲ πνεῦμα κατὰ τῆς σαρκός, ταῦτα γὰρ ἀλλήλοις ἀντίκειται, ἵνα μὴ ἃ ἐὰν θέλητε ταῦτα ποιῆτε.
\vs{18}~εἰ δὲ πνεύματι ἄγεσθε, οὐκ ἐστὲ ὑπὸ νόμον.
\vs{19}~φανερὰ δέ ἐστιν τὰ ἔργα τῆς σαρκός, ἅτινά ἐστιν πορνεία, ἀκαθαρσία, ἀσέλγεια,
\vs{20}~εἰδωλολατρία, φαρμακεία, ἔχθραι, ἔρις, ζῆλος, θυμοί, ἐριθεῖαι, διχοστασίαι, αἱρέσεις,
\vs{21}~φθόνοι, μέθαι, κῶμοι, καὶ τὰ ὅμοια τούτοις, ἃ προλέγω ὑμῖν καθὼς προεῖπον ὅτι οἱ τὰ τοιαῦτα πράσσοντες βασιλείαν θεοῦ οὐ κληρονομήσουσιν.
\vs{22}~Ὁ δὲ καρπὸς τοῦ πνεύματός ἐστιν ἀγάπη, χαρά, εἰρήνη, μακροθυμία, χρηστότης, ἀγαθωσύνη, πίστις,
\vs{23}~πραΰτης, ἐγκράτεια· κατὰ τῶν τοιούτων οὐκ ἔστιν νόμος.
\vs{24}~οἱ δὲ τοῦ Χριστοῦ τὴν σάρκα ἐσταύρωσαν σὺν τοῖς παθήμασιν καὶ ταῖς ἐπιθυμίαις.
\vs{25}~εἰ ζῶμεν πνεύματι, πνεύματι καὶ στοιχῶμεν.
\vs{26}~μὴ γινώμεθα κενόδοξοι, ἀλλήλους προκαλούμενοι, ἀλλήλοις φθονοῦντες.

\chapheader{6}
\vs{1}~\adeA{Ἀδελφοί}, ἐὰν καὶ προλημφθῇ ἄνθρωπος ἔν τινι παραπτώματι, ὑμεῖς οἱ πνευματικοὶ καταρτίζετε τὸν τοιοῦτον ἐν πνεύματι πραΰτητος, σκοπῶν σεαυτόν, μὴ καὶ σὺ πειρασθῇς.
\vs{2}~ἀλλήλων τὰ βάρη βαστάζετε, καὶ οὕτως ἀναπληρώσετε τὸν νόμον τοῦ Χριστοῦ.
\vs{3}~εἰ γὰρ δοκεῖ τις εἶναί τι μηδὲν ὤν, φρεναπατᾷ ἑαυτόν·
\vs{4}~τὸ δὲ ἔργον ἑαυτοῦ δοκιμαζέτω ἕκαστος, καὶ τότε εἰς ἑαυτὸν μόνον τὸ καύχημα ἕξει καὶ οὐκ εἰς τὸν ἕτερον,
\vs{5}~ἕκαστος γὰρ τὸ ἴδιον φορτίον βαστάσει.
\vs{6}~Κοινωνείτω δὲ ὁ κατηχούμενος τὸν λόγον τῷ κατηχοῦντι ἐν πᾶσιν ἀγαθοῖς.
\vs{7}~μὴ πλανᾶσθε, θεὸς οὐ μυκτηρίζεται· ὃ γὰρ ἐὰν σπείρῃ ἄνθρωπος, τοῦτο καὶ θερίσει·
\vs{8}~ὅτι ὁ σπείρων εἰς τὴν σάρκα ἑαυτοῦ ἐκ τῆς σαρκὸς θερίσει φθοράν, ὁ δὲ σπείρων εἰς τὸ πνεῦμα ἐκ τοῦ πνεύματος θερίσει ζωὴν αἰώνιον.
\vs{9}~τὸ δὲ καλὸν ποιοῦντες μὴ ἐγκακῶμεν, καιρῷ γὰρ ἰδίῳ θερίσομεν μὴ ἐκλυόμενοι.
\vs{10}~ἄρα οὖν ὡς καιρὸν ἔχομεν, ἐργαζώμεθα τὸ ἀγαθὸν πρὸς πάντας, μάλιστα δὲ πρὸς τοὺς οἰκείους τῆς πίστεως.
\vs{11}~Ἴδετε πηλίκοις ὑμῖν γράμμασιν ἔγραψα τῇ ἐμῇ χειρί.
\vs{12}~ὅσοι θέλουσιν εὐπροσωπῆσαι ἐν σαρκί, οὗτοι ἀναγκάζουσιν ὑμᾶς περιτέμνεσθαι, μόνον ἵνα τῷ σταυρῷ τοῦ Χριστοῦ μὴ διώκωνται·
\vs{13}~οὐδὲ γὰρ οἱ περιτεμνόμενοι αὐτοὶ νόμον φυλάσσουσιν, ἀλλὰ θέλουσιν ὑμᾶς περιτέμνεσθαι ἵνα ἐν τῇ ὑμετέρᾳ σαρκὶ καυχήσωνται.
\vs{14}~ἐμοὶ δὲ μὴ γένοιτο καυχᾶσθαι εἰ μὴ ἐν τῷ σταυρῷ τοῦ κυρίου ἡμῶν Ἰησοῦ Χριστοῦ, δι’ οὗ ἐμοὶ κόσμος ἐσταύρωται κἀγὼ κόσμῳ.
\vs{15}~οὔτε γὰρ περιτομή τί ἐστιν οὔτε ἀκροβυστία, ἀλλὰ καινὴ κτίσις.
\vs{16}~καὶ ὅσοι τῷ κανόνι τούτῳ στοιχήσουσιν, εἰρήνη ἐπ’ αὐτοὺς καὶ ἔλεος, καὶ ἐπὶ τὸν Ἰσραὴλ τοῦ θεοῦ.
\vs{17}~Τοῦ λοιποῦ κόπους μοι μηδεὶς παρεχέτω, ἐγὼ γὰρ τὰ στίγματα τοῦ Ἰησοῦ ἐν τῷ σώματί μου βαστάζω.
\vs{18}~Ἡ χάρις τοῦ κυρίου ἡμῶν Ἰησοῦ Χριστοῦ μετὰ τοῦ πνεύματος ὑμῶν, \adeA{ἀδελφοί}· ἀμήν.

\booksection{Philippians}{Πρὸς Φιλιππησίους}

\chapheader{1}
\vs{1}~Παῦλος καὶ Τιμόθεος δοῦλοι Χριστοῦ Ἰησοῦ πᾶσιν τοῖς ἁγίοις ἐν Χριστῷ Ἰησοῦ τοῖς οὖσιν ἐν Φιλίπποις σὺν ἐπισκόποις καὶ διακόνοις·
\vs{2}~χάρις ὑμῖν καὶ εἰρήνη ἀπὸ θεοῦ πατρὸς ἡμῶν καὶ κυρίου Ἰησοῦ Χριστοῦ.
\vs{3}~Εὐχαριστῶ τῷ θεῷ μου ἐπὶ πάσῃ τῇ μνείᾳ ὑμῶν
\vs{4}~πάντοτε ἐν πάσῃ δεήσει μου ὑπὲρ πάντων ὑμῶν, μετὰ χαρᾶς τὴν δέησιν ποιούμενος,
\vs{5}~ἐπὶ τῇ κοινωνίᾳ ὑμῶν εἰς τὸ εὐαγγέλιον ἀπὸ τῆς πρώτης ἡμέρας ἄχρι τοῦ νῦν,
\vs{6}~πεποιθὼς αὐτὸ τοῦτο ὅτι ὁ ἐναρξάμενος ἐν ὑμῖν ἔργον ἀγαθὸν ἐπιτελέσει ἄχρι ἡμέρας Χριστοῦ Ἰησοῦ·
\vs{7}~καθώς ἐστιν δίκαιον ἐμοὶ τοῦτο φρονεῖν ὑπὲρ πάντων ὑμῶν, διὰ τὸ ἔχειν με ἐν τῇ καρδίᾳ ὑμᾶς, ἔν τε τοῖς δεσμοῖς μου καὶ ἐν τῇ ἀπολογίᾳ καὶ βεβαιώσει τοῦ εὐαγγελίου συγκοινωνούς μου τῆς χάριτος πάντας ὑμᾶς ὄντας·
\vs{8}~μάρτυς γάρ μου ὁ θεός, ὡς ἐπιποθῶ πάντας ὑμᾶς ἐν σπλάγχνοις Χριστοῦ Ἰησοῦ.
\vs{9}~καὶ τοῦτο προσεύχομαι ἵνα ἡ ἀγάπη ὑμῶν ἔτι μᾶλλον καὶ μᾶλλον περισσεύῃ ἐν ἐπιγνώσει καὶ πάσῃ αἰσθήσει,
\vs{10}~εἰς τὸ δοκιμάζειν ὑμᾶς τὰ διαφέροντα, ἵνα ἦτε εἰλικρινεῖς καὶ ἀπρόσκοποι εἰς ἡμέραν Χριστοῦ,
\vs{11}~πεπληρωμένοι καρπὸν δικαιοσύνης τὸν διὰ Ἰησοῦ Χριστοῦ εἰς δόξαν καὶ ἔπαινον θεοῦ.
\vs{12}~Γινώσκειν δὲ ὑμᾶς βούλομαι, \adeA{ἀδελφοί}, ὅτι τὰ κατ’ ἐμὲ μᾶλλον εἰς προκοπὴν τοῦ εὐαγγελίου ἐλήλυθεν,
\vs{13}~ὥστε τοὺς δεσμούς μου φανεροὺς ἐν Χριστῷ γενέσθαι ἐν ὅλῳ τῷ πραιτωρίῳ καὶ τοῖς λοιποῖς πᾶσιν,
\vs{14}~καὶ τοὺς πλείονας τῶν \adeA{ἀδελφῶν} ἐν κυρίῳ πεποιθότας τοῖς δεσμοῖς μου περισσοτέρως τολμᾶν ἀφόβως τὸν λόγον λαλεῖν.
\vs{15}~Τινὲς μὲν καὶ διὰ φθόνον καὶ ἔριν, τινὲς δὲ καὶ δι’ εὐδοκίαν τὸν Χριστὸν κηρύσσουσιν·
\vs{16}~οἱ μὲν ἐξ ἀγάπης, εἰδότες ὅτι εἰς ἀπολογίαν τοῦ εὐαγγελίου κεῖμαι,
\vs{17}~οἱ δὲ ἐξ ἐριθείας τὸν Χριστὸν καταγγέλλουσιν, οὐχ ἁγνῶς, οἰόμενοι θλῖψιν ἐγείρειν τοῖς δεσμοῖς μου.
\vs{18}~τί γάρ; πλὴν ὅτι παντὶ τρόπῳ, εἴτε προφάσει εἴτε ἀληθείᾳ, Χριστὸς καταγγέλλεται, καὶ ἐν τούτῳ χαίρω· ἀλλὰ καὶ χαρήσομαι,
\vs{19}~οἶδα γὰρ ὅτι τοῦτό μοι ἀποβήσεται εἰς σωτηρίαν διὰ τῆς ὑμῶν δεήσεως καὶ ἐπιχορηγίας τοῦ πνεύματος Ἰησοῦ Χριστοῦ,
\vs{20}~κατὰ τὴν ἀποκαραδοκίαν καὶ ἐλπίδα μου ὅτι ἐν οὐδενὶ αἰσχυνθήσομαι, ἀλλ’ ἐν πάσῃ παρρησίᾳ ὡς πάντοτε καὶ νῦν μεγαλυνθήσεται Χριστὸς ἐν τῷ σώματί μου, εἴτε διὰ ζωῆς εἴτε διὰ θανάτου.
\vs{21}~ἐμοὶ γὰρ τὸ ζῆν Χριστὸς καὶ τὸ ἀποθανεῖν κέρδος.
\vs{22}~εἰ δὲ τὸ ζῆν ἐν σαρκί, τοῦτό μοι καρπὸς ἔργου— καὶ τί αἱρήσομαι οὐ γνωρίζω·
\vs{23}~συνέχομαι δὲ ἐκ τῶν δύο, τὴν ἐπιθυμίαν ἔχων εἰς τὸ ἀναλῦσαι καὶ σὺν Χριστῷ εἶναι, πολλῷ γὰρ μᾶλλον κρεῖσσον,
\vs{24}~τὸ δὲ ἐπιμένειν ἐν τῇ σαρκὶ ἀναγκαιότερον δι’ ὑμᾶς.
\vs{25}~καὶ τοῦτο πεποιθὼς οἶδα ὅτι μενῶ καὶ παραμενῶ πᾶσιν ὑμῖν εἰς τὴν ὑμῶν προκοπὴν καὶ χαρὰν τῆς πίστεως,
\vs{26}~ἵνα τὸ καύχημα ὑμῶν περισσεύῃ ἐν Χριστῷ Ἰησοῦ ἐν ἐμοὶ διὰ τῆς ἐμῆς παρουσίας πάλιν πρὸς ὑμᾶς.
\vs{27}~Μόνον ἀξίως τοῦ εὐαγγελίου τοῦ Χριστοῦ πολιτεύεσθε, ἵνα εἴτε ἐλθὼν καὶ ἰδὼν ὑμᾶς εἴτε ἀπὼν ἀκούω τὰ περὶ ὑμῶν, ὅτι στήκετε ἐν ἑνὶ πνεύματι, μιᾷ ψυχῇ συναθλοῦντες τῇ πίστει τοῦ εὐαγγελίου,
\vs{28}~καὶ μὴ πτυρόμενοι ἐν μηδενὶ ὑπὸ τῶν ἀντικειμένων (ἥτις ἐστὶν αὐτοῖς ἔνδειξις ἀπωλείας, ὑμῶν δὲ σωτηρίας, καὶ τοῦτο ἀπὸ θεοῦ,
\vs{29}~ὅτι ὑμῖν ἐχαρίσθη τὸ ὑπὲρ Χριστοῦ, οὐ μόνον τὸ εἰς αὐτὸν πιστεύειν ἀλλὰ καὶ τὸ ὑπὲρ αὐτοῦ πάσχειν),
\vs{30}~τὸν αὐτὸν ἀγῶνα ἔχοντες οἷον εἴδετε ἐν ἐμοὶ καὶ νῦν ἀκούετε ἐν ἐμοί.

\chapheader{2}
\vs{1}~Εἴ τις οὖν παράκλησις ἐν Χριστῷ, εἴ τι παραμύθιον ἀγάπης, εἴ τις κοινωνία πνεύματος, εἴ τις σπλάγχνα καὶ οἰκτιρμοί,
\vs{2}~πληρώσατέ μου τὴν χαρὰν ἵνα τὸ αὐτὸ φρονῆτε, τὴν αὐτὴν ἀγάπην ἔχοντες, σύμψυχοι, τὸ ἓν φρονοῦντες,
\vs{3}~μηδὲν κατ’ ἐριθείαν μηδὲ κατὰ κενοδοξίαν, ἀλλὰ τῇ ταπεινοφροσύνῃ ἀλλήλους ἡγούμενοι ὑπερέχοντας ἑαυτῶν,
\vs{4}~μὴ τὰ ἑαυτῶν ἕκαστοι σκοποῦντες, ἀλλὰ καὶ τὰ ἑτέρων ἕκαστοι.
\vs{5}~τοῦτο φρονεῖτε ἐν ὑμῖν ὃ καὶ ἐν Χριστῷ Ἰησοῦ,
\vs{6}~ὃς ἐν μορφῇ θεοῦ ὑπάρχων οὐχ ἁρπαγμὸν ἡγήσατο τὸ εἶναι ἴσα θεῷ,
\vs{7}~ἀλλὰ ἑαυτὸν ἐκένωσεν μορφὴν δούλου λαβών, ἐν ὁμοιώματι ἀνθρώπων γενόμενος· καὶ σχήματι εὑρεθεὶς ὡς ἄνθρωπος
\vs{8}~ἐταπείνωσεν ἑαυτὸν γενόμενος ὑπήκοος μέχρι θανάτου, θανάτου δὲ σταυροῦ·
\vs{9}~διὸ καὶ ὁ θεὸς αὐτὸν ὑπερύψωσεν, καὶ ἐχαρίσατο αὐτῷ τὸ ὄνομα τὸ ὑπὲρ πᾶν ὄνομα,
\vs{10}~ἵνα ἐν τῷ ὀνόματι Ἰησοῦ πᾶν γόνυ κάμψῃ ἐπουρανίων καὶ ἐπιγείων καὶ καταχθονίων,
\vs{11}~καὶ πᾶσα γλῶσσα ἐξομολογήσηται ὅτι κύριος Ἰησοῦς Χριστὸς εἰς δόξαν θεοῦ πατρός.
\vs{12}~Ὥστε, ἀγαπητοί μου, καθὼς πάντοτε ὑπηκούσατε, μὴ ὡς ἐν τῇ παρουσίᾳ μου μόνον ἀλλὰ νῦν πολλῷ μᾶλλον ἐν τῇ ἀπουσίᾳ μου, μετὰ φόβου καὶ τρόμου τὴν ἑαυτῶν σωτηρίαν κατεργάζεσθε,
\vs{13}~θεὸς γάρ ἐστιν ὁ ἐνεργῶν ἐν ὑμῖν καὶ τὸ θέλειν καὶ τὸ ἐνεργεῖν ὑπὲρ τῆς εὐδοκίας.
\vs{14}~Πάντα ποιεῖτε χωρὶς γογγυσμῶν καὶ διαλογισμῶν,
\vs{15}~ἵνα γένησθε ἄμεμπτοι καὶ ἀκέραιοι, τέκνα θεοῦ ἄμωμα μέσον γενεᾶς σκολιᾶς καὶ διεστραμμένης, ἐν οἷς φαίνεσθε ὡς φωστῆρες ἐν κόσμῳ
\vs{16}~λόγον ζωῆς ἐπέχοντες, εἰς καύχημα ἐμοὶ εἰς ἡμέραν Χριστοῦ, ὅτι οὐκ εἰς κενὸν ἔδραμον οὐδὲ εἰς κενὸν ἐκοπίασα.
\vs{17}~ἀλλὰ εἰ καὶ σπένδομαι ἐπὶ τῇ θυσίᾳ καὶ λειτουργίᾳ τῆς πίστεως ὑμῶν, χαίρω καὶ συγχαίρω πᾶσιν ὑμῖν·
\vs{18}~τὸ δὲ αὐτὸ καὶ ὑμεῖς χαίρετε καὶ συγχαίρετέ μοι.
\vs{19}~Ἐλπίζω δὲ ἐν κυρίῳ Ἰησοῦ Τιμόθεον ταχέως πέμψαι ὑμῖν, ἵνα κἀγὼ εὐψυχῶ γνοὺς τὰ περὶ ὑμῶν.
\vs{20}~οὐδένα γὰρ ἔχω ἰσόψυχον ὅστις γνησίως τὰ περὶ ὑμῶν μεριμνήσει,
\vs{21}~οἱ πάντες γὰρ τὰ ἑαυτῶν ζητοῦσιν, οὐ τὰ Ἰησοῦ Χριστοῦ.
\vs{22}~τὴν δὲ δοκιμὴν αὐτοῦ γινώσκετε, ὅτι ὡς πατρὶ τέκνον σὺν ἐμοὶ ἐδούλευσεν εἰς τὸ εὐαγγέλιον.
\vs{23}~τοῦτον μὲν οὖν ἐλπίζω πέμψαι ὡς ἂν ἀφίδω τὰ περὶ ἐμὲ ἐξαυτῆς·
\vs{24}~πέποιθα δὲ ἐν κυρίῳ ὅτι καὶ αὐτὸς ταχέως ἐλεύσομαι.
\vs{25}~Ἀναγκαῖον δὲ ἡγησάμην Ἐπαφρόδιτον τὸν \adeA{ἀδελφὸν} καὶ συνεργὸν καὶ συστρατιώτην μου, ὑμῶν δὲ ἀπόστολον καὶ λειτουργὸν τῆς χρείας μου, πέμψαι πρὸς ὑμᾶς,
\vs{26}~ἐπειδὴ ἐπιποθῶν ἦν πάντας ὑμᾶς, καὶ ἀδημονῶν διότι ἠκούσατε ὅτι ἠσθένησεν.
\vs{27}~καὶ γὰρ ἠσθένησεν παραπλήσιον θανάτῳ· ἀλλὰ ὁ θεὸς ἠλέησεν αὐτόν, οὐκ αὐτὸν δὲ μόνον ἀλλὰ καὶ ἐμέ, ἵνα μὴ λύπην ἐπὶ λύπην σχῶ.
\vs{28}~σπουδαιοτέρως οὖν ἔπεμψα αὐτὸν ἵνα ἰδόντες αὐτὸν πάλιν χαρῆτε κἀγὼ ἀλυπότερος ὦ.
\vs{29}~προσδέχεσθε οὖν αὐτὸν ἐν κυρίῳ μετὰ πάσης χαρᾶς, καὶ τοὺς τοιούτους ἐντίμους ἔχετε,
\vs{30}~ὅτι διὰ τὸ ἔργον Χριστοῦ μέχρι θανάτου ἤγγισεν, παραβολευσάμενος τῇ ψυχῇ ἵνα ἀναπληρώσῃ τὸ ὑμῶν ὑστέρημα τῆς πρός με λειτουργίας.

\chapheader{3}
\vs{1}~Τὸ λοιπόν, \adeA{ἀδελφοί} μου, χαίρετε ἐν κυρίῳ. τὰ αὐτὰ γράφειν ὑμῖν ἐμοὶ μὲν οὐκ ὀκνηρόν, ὑμῖν δὲ ἀσφαλές.
\vs{2}~Βλέπετε τοὺς κύνας, βλέπετε τοὺς κακοὺς ἐργάτας, βλέπετε τὴν κατατομήν.
\vs{3}~ἡμεῖς γάρ ἐσμεν ἡ περιτομή, οἱ πνεύματι θεοῦ λατρεύοντες καὶ καυχώμενοι ἐν Χριστῷ Ἰησοῦ καὶ οὐκ ἐν σαρκὶ πεποιθότες,
\vs{4}~καίπερ ἐγὼ ἔχων πεποίθησιν καὶ ἐν σαρκί. Εἴ τις δοκεῖ ἄλλος πεποιθέναι ἐν σαρκί, ἐγὼ μᾶλλον·
\vs{5}~περιτομῇ ὀκταήμερος, ἐκ γένους Ἰσραήλ, φυλῆς Βενιαμίν, Ἑβραῖος ἐξ Ἑβραίων, κατὰ νόμον Φαρισαῖος,
\vs{6}~κατὰ ζῆλος διώκων τὴν ἐκκλησίαν, κατὰ δικαιοσύνην τὴν ἐν νόμῳ γενόμενος ἄμεμπτος.
\vs{7}~Ἀλλὰ ἅτινα ἦν μοι κέρδη, ταῦτα ἥγημαι διὰ τὸν Χριστὸν ζημίαν.
\vs{8}~ἀλλὰ μενοῦνγε καὶ ἡγοῦμαι πάντα ζημίαν εἶναι διὰ τὸ ὑπερέχον τῆς γνώσεως Χριστοῦ Ἰησοῦ τοῦ κυρίου μου δι’ ὃν τὰ πάντα ἐζημιώθην, καὶ ἡγοῦμαι σκύβαλα ἵνα Χριστὸν κερδήσω
\vs{9}~καὶ εὑρεθῶ ἐν αὐτῷ, μὴ ἔχων ἐμὴν δικαιοσύνην τὴν ἐκ νόμου ἀλλὰ τὴν διὰ πίστεως Χριστοῦ, τὴν ἐκ θεοῦ δικαιοσύνην ἐπὶ τῇ πίστει,
\vs{10}~τοῦ γνῶναι αὐτὸν καὶ τὴν δύναμιν τῆς ἀναστάσεως αὐτοῦ καὶ κοινωνίαν παθημάτων αὐτοῦ, συμμορφιζόμενος τῷ θανάτῳ αὐτοῦ,
\vs{11}~εἴ πως καταντήσω εἰς τὴν ἐξανάστασιν τὴν ἐκ νεκρῶν.
\vs{12}~Οὐχ ὅτι ἤδη ἔλαβον ἢ ἤδη τετελείωμαι, διώκω δὲ εἰ καὶ καταλάβω, ἐφ’ ᾧ καὶ κατελήμφθην ὑπὸ Χριστοῦ.
\vs{13}~\adeA{ἀδελφοί}, ἐγὼ ἐμαυτὸν οὐ λογίζομαι κατειληφέναι· ἓν δέ, τὰ μὲν ὀπίσω ἐπιλανθανόμενος τοῖς δὲ ἔμπροσθεν ἐπεκτεινόμενος,
\vs{14}~κατὰ σκοπὸν διώκω εἰς τὸ βραβεῖον τῆς ἄνω κλήσεως τοῦ θεοῦ ἐν Χριστῷ Ἰησοῦ.
\vs{15}~ὅσοι οὖν τέλειοι, τοῦτο φρονῶμεν· καὶ εἴ τι ἑτέρως φρονεῖτε, καὶ τοῦτο ὁ θεὸς ὑμῖν ἀποκαλύψει·
\vs{16}~πλὴν εἰς ὃ ἐφθάσαμεν, τῷ αὐτῷ στοιχεῖν.
\vs{17}~Συμμιμηταί μου γίνεσθε, \adeA{ἀδελφοί}, καὶ σκοπεῖτε τοὺς οὕτω περιπατοῦντας καθὼς ἔχετε τύπον ἡμᾶς·
\vs{18}~πολλοὶ γὰρ περιπατοῦσιν οὓς πολλάκις ἔλεγον ὑμῖν, νῦν δὲ καὶ κλαίων λέγω, τοὺς ἐχθροὺς τοῦ σταυροῦ τοῦ Χριστοῦ,
\vs{19}~ὧν τὸ τέλος ἀπώλεια, ὧν ὁ θεὸς ἡ κοιλία καὶ ἡ δόξα ἐν τῇ αἰσχύνῃ αὐτῶν, οἱ τὰ ἐπίγεια φρονοῦντες.
\vs{20}~ἡμῶν γὰρ τὸ πολίτευμα ἐν οὐρανοῖς ὑπάρχει, ἐξ οὗ καὶ σωτῆρα ἀπεκδεχόμεθα κύριον Ἰησοῦν Χριστόν,
\vs{21}~ὃς μετασχηματίσει τὸ σῶμα τῆς ταπεινώσεως ἡμῶν σύμμορφον τῷ σώματι τῆς δόξης αὐτοῦ κατὰ τὴν ἐνέργειαν τοῦ δύνασθαι αὐτὸν καὶ ὑποτάξαι αὑτῷ τὰ πάντα.

\chapheader{4}
\vs{1}~ὥστε, \adeA{ἀδελφοί} μου ἀγαπητοὶ καὶ ἐπιπόθητοι, χαρὰ καὶ στέφανός μου, οὕτως στήκετε ἐν κυρίῳ, ἀγαπητοί.
\vs{2}~Εὐοδίαν παρακαλῶ καὶ Συντύχην παρακαλῶ τὸ αὐτὸ φρονεῖν ἐν κυρίῳ.
\vs{3}~ναὶ ἐρωτῶ καὶ σέ, γνήσιε σύζυγε, συλλαμβάνου αὐταῖς, αἵτινες ἐν τῷ εὐαγγελίῳ συνήθλησάν μοι μετὰ καὶ Κλήμεντος καὶ τῶν λοιπῶν συνεργῶν μου, ὧν τὰ ὀνόματα ἐν βίβλῳ ζωῆς.
\vs{4}~Χαίρετε ἐν κυρίῳ πάντοτε· πάλιν ἐρῶ, χαίρετε.
\vs{5}~τὸ ἐπιεικὲς ὑμῶν γνωσθήτω πᾶσιν ἀνθρώποις. ὁ κύριος ἐγγύς·
\vs{6}~μηδὲν μεριμνᾶτε, ἀλλ’ ἐν παντὶ τῇ προσευχῇ καὶ τῇ δεήσει μετ’ εὐχαριστίας τὰ αἰτήματα ὑμῶν γνωριζέσθω πρὸς τὸν θεόν·
\vs{7}~καὶ ἡ εἰρήνη τοῦ θεοῦ ἡ ὑπερέχουσα πάντα νοῦν φρουρήσει τὰς καρδίας ὑμῶν καὶ τὰ νοήματα ὑμῶν ἐν Χριστῷ Ἰησοῦ.
\vs{8}~Τὸ λοιπόν, \adeA{ἀδελφοί}, ὅσα ἐστὶν ἀληθῆ, ὅσα σεμνά, ὅσα δίκαια, ὅσα ἁγνά, ὅσα προσφιλῆ, ὅσα εὔφημα, εἴ τις ἀρετὴ καὶ εἴ τις ἔπαινος, ταῦτα λογίζεσθε·
\vs{9}~ἃ καὶ ἐμάθετε καὶ παρελάβετε καὶ ἠκούσατε καὶ εἴδετε ἐν ἐμοί, ταῦτα πράσσετε· καὶ ὁ θεὸς τῆς εἰρήνης ἔσται μεθ’ ὑμῶν.
\vs{10}~Ἐχάρην δὲ ἐν κυρίῳ μεγάλως ὅτι ἤδη ποτὲ ἀνεθάλετε τὸ ὑπὲρ ἐμοῦ φρονεῖν, ἐφ’ ᾧ καὶ ἐφρονεῖτε ἠκαιρεῖσθε δέ.
\vs{11}~οὐχ ὅτι καθ’ ὑστέρησιν λέγω, ἐγὼ γὰρ ἔμαθον ἐν οἷς εἰμι αὐτάρκης εἶναι·
\vs{12}~οἶδα καὶ ταπεινοῦσθαι, οἶδα καὶ περισσεύειν· ἐν παντὶ καὶ ἐν πᾶσιν μεμύημαι, καὶ χορτάζεσθαι καὶ πεινᾶν, καὶ περισσεύειν καὶ ὑστερεῖσθαι·
\vs{13}~πάντα ἰσχύω ἐν τῷ ἐνδυναμοῦντί με.
\vs{14}~πλὴν καλῶς ἐποιήσατε συγκοινωνήσαντές μου τῇ θλίψει.
\vs{15}~Οἴδατε δὲ καὶ ὑμεῖς, Φιλιππήσιοι, ὅτι ἐν ἀρχῇ τοῦ εὐαγγελίου, ὅτε ἐξῆλθον ἀπὸ Μακεδονίας, οὐδεμία μοι ἐκκλησία ἐκοινώνησεν εἰς λόγον δόσεως καὶ λήμψεως εἰ μὴ ὑμεῖς μόνοι,
\vs{16}~ὅτι καὶ ἐν Θεσσαλονίκῃ καὶ ἅπαξ καὶ δὶς εἰς τὴν χρείαν μοι ἐπέμψατε.
\vs{17}~οὐχ ὅτι ἐπιζητῶ τὸ δόμα, ἀλλὰ ἐπιζητῶ τὸν καρπὸν τὸν πλεονάζοντα εἰς λόγον ὑμῶν.
\vs{18}~ἀπέχω δὲ πάντα καὶ περισσεύω· πεπλήρωμαι δεξάμενος παρὰ Ἐπαφροδίτου τὰ παρ’ ὑμῶν, ὀσμὴν εὐωδίας, θυσίαν δεκτήν, εὐάρεστον τῷ θεῷ.
\vs{19}~ὁ δὲ θεός μου πληρώσει πᾶσαν χρείαν ὑμῶν κατὰ τὸ πλοῦτος αὐτοῦ ἐν δόξῃ ἐν Χριστῷ Ἰησοῦ.
\vs{20}~τῷ δὲ θεῷ καὶ πατρὶ ἡμῶν ἡ δόξα εἰς τοὺς αἰῶνας τῶν αἰώνων· ἀμήν.
\vs{21}~Ἀσπάσασθε πάντα ἅγιον ἐν Χριστῷ Ἰησοῦ. ἀσπάζονται ὑμᾶς οἱ σὺν ἐμοὶ \adeA{ἀδελφοί}.
\vs{22}~ἀσπάζονται ὑμᾶς πάντες οἱ ἅγιοι, μάλιστα δὲ οἱ ἐκ τῆς Καίσαρος οἰκίας.
\vs{23}~ἡ χάρις τοῦ κυρίου Ἰησοῦ Χριστοῦ μετὰ τοῦ πνεύματος ὑμῶν.

\booksection{1 Thessalonians}{Πρὸς Θεσσαλονικεῖς α΄}

\chapheader{1}
\vs{1}~Παῦλος καὶ Σιλουανὸς καὶ Τιμόθεος τῇ ἐκκλησίᾳ Θεσσαλονικέων ἐν θεῷ πατρὶ καὶ κυρίῳ Ἰησοῦ Χριστῷ· χάρις ὑμῖν καὶ εἰρήνη.
\vs{2}~Εὐχαριστοῦμεν τῷ θεῷ πάντοτε περὶ πάντων ὑμῶν μνείαν ποιούμενοι ἐπὶ τῶν προσευχῶν ἡμῶν, ἀδιαλείπτως
\vs{3}~μνημονεύοντες ὑμῶν τοῦ ἔργου τῆς πίστεως καὶ τοῦ κόπου τῆς ἀγάπης καὶ τῆς ὑπομονῆς τῆς ἐλπίδος τοῦ κυρίου ἡμῶν Ἰησοῦ Χριστοῦ ἔμπροσθεν τοῦ θεοῦ καὶ πατρὸς ἡμῶν,
\vs{4}~εἰδότες, \adeA{ἀδελφοὶ} ἠγαπημένοι ὑπὸ θεοῦ, τὴν ἐκλογὴν ὑμῶν,
\vs{5}~ὅτι τὸ εὐαγγέλιον ἡμῶν οὐκ ἐγενήθη εἰς ὑμᾶς ἐν λόγῳ μόνον ἀλλὰ καὶ ἐν δυνάμει καὶ ἐν πνεύματι ἁγίῳ καὶ πληροφορίᾳ πολλῇ, καθὼς οἴδατε οἷοι ἐγενήθημεν ἐν ὑμῖν δι’ ὑμᾶς·
\vs{6}~καὶ ὑμεῖς μιμηταὶ ἡμῶν ἐγενήθητε καὶ τοῦ κυρίου, δεξάμενοι τὸν λόγον ἐν θλίψει πολλῇ μετὰ χαρᾶς πνεύματος ἁγίου,
\vs{7}~ὥστε γενέσθαι ὑμᾶς τύπον πᾶσιν τοῖς πιστεύουσιν ἐν τῇ Μακεδονίᾳ καὶ ἐν τῇ Ἀχαΐᾳ.
\vs{8}~ἀφ’ ὑμῶν γὰρ ἐξήχηται ὁ λόγος τοῦ κυρίου οὐ μόνον ἐν τῇ Μακεδονίᾳ καὶ Ἀχαΐᾳ, ἀλλ’ ἐν παντὶ τόπῳ ἡ πίστις ὑμῶν ἡ πρὸς τὸν θεὸν ἐξελήλυθεν, ὥστε μὴ χρείαν ἔχειν ἡμᾶς λαλεῖν τι·
\vs{9}~αὐτοὶ γὰρ περὶ ἡμῶν ἀπαγγέλλουσιν ὁποίαν εἴσοδον ἔσχομεν πρὸς ὑμᾶς, καὶ πῶς ἐπεστρέψατε πρὸς τὸν θεὸν ἀπὸ τῶν εἰδώλων δουλεύειν θεῷ ζῶντι καὶ ἀληθινῷ,
\vs{10}~καὶ ἀναμένειν τὸν υἱὸν αὐτοῦ ἐκ τῶν οὐρανῶν, ὃν ἤγειρεν ἐκ τῶν νεκρῶν, Ἰησοῦν τὸν ῥυόμενον ἡμᾶς ἐκ τῆς ὀργῆς τῆς ἐρχομένης.

\chapheader{2}
\vs{1}~Αὐτοὶ γὰρ οἴδατε, \adeA{ἀδελφοί}, τὴν εἴσοδον ἡμῶν τὴν πρὸς ὑμᾶς ὅτι οὐ κενὴ γέγονεν,
\vs{2}~ἀλλὰ προπαθόντες καὶ ὑβρισθέντες καθὼς οἴδατε ἐν Φιλίπποις ἐπαρρησιασάμεθα ἐν τῷ θεῷ ἡμῶν λαλῆσαι πρὸς ὑμᾶς τὸ εὐαγγέλιον τοῦ θεοῦ ἐν πολλῷ ἀγῶνι.
\vs{3}~ἡ γὰρ παράκλησις ἡμῶν οὐκ ἐκ πλάνης οὐδὲ ἐξ ἀκαθαρσίας οὐδὲ ἐν δόλῳ,
\vs{4}~ἀλλὰ καθὼς δεδοκιμάσμεθα ὑπὸ τοῦ θεοῦ πιστευθῆναι τὸ εὐαγγέλιον οὕτως λαλοῦμεν, οὐχ ὡς ἀνθρώποις ἀρέσκοντες ἀλλὰ θεῷ τῷ δοκιμάζοντι τὰς καρδίας ἡμῶν.
\vs{5}~οὔτε γάρ ποτε ἐν λόγῳ κολακείας ἐγενήθημεν, καθὼς οἴδατε, οὔτε ἐν προφάσει πλεονεξίας, θεὸς μάρτυς,
\vs{6}~οὔτε ζητοῦντες ἐξ ἀνθρώπων δόξαν, οὔτε ἀφ’ ὑμῶν οὔτε ἀπ’ ἄλλων,
\vs{7}~δυνάμενοι ἐν βάρει εἶναι ὡς Χριστοῦ ἀπόστολοι· ἀλλὰ ἐγενήθημεν ἤπιοι ἐν μέσῳ ὑμῶν, ὡς ἐὰν τροφὸς θάλπῃ τὰ ἑαυτῆς τέκνα·
\vs{8}~οὕτως ὁμειρόμενοι ὑμῶν εὐδοκοῦμεν μεταδοῦναι ὑμῖν οὐ μόνον τὸ εὐαγγέλιον τοῦ θεοῦ ἀλλὰ καὶ τὰς ἑαυτῶν ψυχάς, διότι ἀγαπητοὶ ἡμῖν ἐγενήθητε.
\vs{9}~Μνημονεύετε γάρ, \adeA{ἀδελφοί}, τὸν κόπον ἡμῶν καὶ τὸν μόχθον· νυκτὸς καὶ ἡμέρας ἐργαζόμενοι πρὸς τὸ μὴ ἐπιβαρῆσαί τινα ὑμῶν ἐκηρύξαμεν εἰς ὑμᾶς τὸ εὐαγγέλιον τοῦ θεοῦ.
\vs{10}~ὑμεῖς μάρτυρες καὶ ὁ θεός, ὡς ὁσίως καὶ δικαίως καὶ ἀμέμπτως ὑμῖν τοῖς πιστεύουσιν ἐγενήθημεν,
\vs{11}~καθάπερ οἴδατε ὡς ἕνα ἕκαστον ὑμῶν ὡς πατὴρ τέκνα ἑαυτοῦ
\vs{12}~παρακαλοῦντες ὑμᾶς καὶ παραμυθούμενοι καὶ μαρτυρόμενοι, εἰς τὸ περιπατεῖν ὑμᾶς ἀξίως τοῦ θεοῦ τοῦ καλοῦντος ὑμᾶς εἰς τὴν ἑαυτοῦ βασιλείαν καὶ δόξαν.
\vs{13}~Καὶ διὰ τοῦτο καὶ ἡμεῖς εὐχαριστοῦμεν τῷ θεῷ ἀδιαλείπτως, ὅτι παραλαβόντες λόγον ἀκοῆς παρ’ ἡμῶν τοῦ θεοῦ ἐδέξασθε οὐ λόγον ἀνθρώπων ἀλλὰ καθὼς ἀληθῶς ἐστὶν λόγον θεοῦ, ὃς καὶ ἐνεργεῖται ἐν ὑμῖν τοῖς πιστεύουσιν.
\vs{14}~ὑμεῖς γὰρ μιμηταὶ ἐγενήθητε, \adeA{ἀδελφοί}, τῶν ἐκκλησιῶν τοῦ θεοῦ τῶν οὐσῶν ἐν τῇ Ἰουδαίᾳ ἐν Χριστῷ Ἰησοῦ, ὅτι τὰ αὐτὰ ἐπάθετε καὶ ὑμεῖς ὑπὸ τῶν ἰδίων συμφυλετῶν καθὼς καὶ αὐτοὶ ὑπὸ τῶν Ἰουδαίων
\vs{15}~τῶν καὶ τὸν κύριον ἀποκτεινάντων Ἰησοῦν καὶ τοὺς προφήτας καὶ ἡμᾶς ἐκδιωξάντων, καὶ θεῷ μὴ ἀρεσκόντων, καὶ πᾶσιν ἀνθρώποις ἐναντίων,
\vs{16}~κωλυόντων ἡμᾶς τοῖς ἔθνεσιν λαλῆσαι ἵνα σωθῶσιν, εἰς τὸ ἀναπληρῶσαι αὐτῶν τὰς ἁμαρτίας πάντοτε. ἔφθασεν δὲ ἐπ’ αὐτοὺς ἡ ὀργὴ εἰς τέλος.
\vs{17}~Ἡμεῖς δέ, \adeA{ἀδελφοί}, ἀπορφανισθέντες ἀφ’ ὑμῶν πρὸς καιρὸν ὥρας, προσώπῳ οὐ καρδίᾳ, περισσοτέρως ἐσπουδάσαμεν τὸ πρόσωπον ὑμῶν ἰδεῖν ἐν πολλῇ ἐπιθυμίᾳ.
\vs{18}~διότι ἠθελήσαμεν ἐλθεῖν πρὸς ὑμᾶς, ἐγὼ μὲν Παῦλος καὶ ἅπαξ καὶ δίς, καὶ ἐνέκοψεν ἡμᾶς ὁ Σατανᾶς.
\vs{19}~τίς γὰρ ἡμῶν ἐλπὶς ἢ χαρὰ ἢ στέφανος καυχήσεως— ἢ οὐχὶ καὶ ὑμεῖς— ἔμπροσθεν τοῦ κυρίου ἡμῶν Ἰησοῦ ἐν τῇ αὐτοῦ παρουσίᾳ;
\vs{20}~ὑμεῖς γάρ ἐστε ἡ δόξα ἡμῶν καὶ ἡ χαρά.

\chapheader{3}
\vs{1}~Διὸ μηκέτι στέγοντες εὐδοκήσαμεν καταλειφθῆναι ἐν Ἀθήναις μόνοι,
\vs{2}~καὶ ἐπέμψαμεν Τιμόθεον, τὸν \adeA{ἀδελφὸν} ἡμῶν καὶ συνεργὸν τοῦ θεοῦ ἐν τῷ εὐαγγελίῳ τοῦ Χριστοῦ, εἰς τὸ στηρίξαι ὑμᾶς καὶ παρακαλέσαι ὑπὲρ τῆς πίστεως ὑμῶν
\vs{3}~τὸ μηδένα σαίνεσθαι ἐν ταῖς θλίψεσιν ταύταις. αὐτοὶ γὰρ οἴδατε ὅτι εἰς τοῦτο κείμεθα·
\vs{4}~καὶ γὰρ ὅτε πρὸς ὑμᾶς ἦμεν, προελέγομεν ὑμῖν ὅτι μέλλομεν θλίβεσθαι, καθὼς καὶ ἐγένετο καὶ οἴδατε.
\vs{5}~διὰ τοῦτο κἀγὼ μηκέτι στέγων ἔπεμψα εἰς τὸ γνῶναι τὴν πίστιν ὑμῶν, μή πως ἐπείρασεν ὑμᾶς ὁ πειράζων καὶ εἰς κενὸν γένηται ὁ κόπος ἡμῶν.
\vs{6}~Ἄρτι δὲ ἐλθόντος Τιμοθέου πρὸς ἡμᾶς ἀφ’ ὑμῶν καὶ εὐαγγελισαμένου ἡμῖν τὴν πίστιν καὶ τὴν ἀγάπην ὑμῶν, καὶ ὅτι ἔχετε μνείαν ἡμῶν ἀγαθὴν πάντοτε ἐπιποθοῦντες ἡμᾶς ἰδεῖν καθάπερ καὶ ἡμεῖς ὑμᾶς,
\vs{7}~διὰ τοῦτο παρεκλήθημεν, \adeA{ἀδελφοί}, ἐφ’ ὑμῖν ἐπὶ πάσῃ τῇ ἀνάγκῃ καὶ θλίψει ἡμῶν διὰ τῆς ὑμῶν πίστεως,
\vs{8}~ὅτι νῦν ζῶμεν ἐὰν ὑμεῖς στήκετε ἐν κυρίῳ.
\vs{9}~τίνα γὰρ εὐχαριστίαν δυνάμεθα τῷ θεῷ ἀνταποδοῦναι περὶ ὑμῶν ἐπὶ πάσῃ τῇ χαρᾷ ᾗ χαίρομεν δι’ ὑμᾶς ἔμπροσθεν τοῦ θεοῦ ἡμῶν,
\vs{10}~νυκτὸς καὶ ἡμέρας ὑπερεκπερισσοῦ δεόμενοι εἰς τὸ ἰδεῖν ὑμῶν τὸ πρόσωπον καὶ καταρτίσαι τὰ ὑστερήματα τῆς πίστεως ὑμῶν;
\vs{11}~Αὐτὸς δὲ ὁ θεὸς καὶ πατὴρ ἡμῶν καὶ ὁ κύριος ἡμῶν Ἰησοῦς κατευθύναι τὴν ὁδὸν ἡμῶν πρὸς ὑμᾶς·
\vs{12}~ὑμᾶς δὲ ὁ κύριος πλεονάσαι καὶ περισσεύσαι τῇ ἀγάπῃ εἰς ἀλλήλους καὶ εἰς πάντας, καθάπερ καὶ ἡμεῖς εἰς ὑμᾶς,
\vs{13}~εἰς τὸ στηρίξαι ὑμῶν τὰς καρδίας ἀμέμπτους ἐν ἁγιωσύνῃ ἔμπροσθεν τοῦ θεοῦ καὶ πατρὸς ἡμῶν ἐν τῇ παρουσίᾳ τοῦ κυρίου ἡμῶν Ἰησοῦ μετὰ πάντων τῶν ἁγίων αὐτοῦ.

\chapheader{4}
\vs{1}~Λοιπὸν οὖν, \adeA{ἀδελφοί}, ἐρωτῶμεν ὑμᾶς καὶ παρακαλοῦμεν ἐν κυρίῳ Ἰησοῦ, ἵνα καθὼς παρελάβετε παρ’ ἡμῶν τὸ πῶς δεῖ ὑμᾶς περιπατεῖν καὶ ἀρέσκειν θεῷ, καθὼς καὶ περιπατεῖτε, ἵνα περισσεύητε μᾶλλον.
\vs{2}~οἴδατε γὰρ τίνας παραγγελίας ἐδώκαμεν ὑμῖν διὰ τοῦ κυρίου Ἰησοῦ.
\vs{3}~τοῦτο γάρ ἐστιν θέλημα τοῦ θεοῦ, ὁ ἁγιασμὸς ὑμῶν, ἀπέχεσθαι ὑμᾶς ἀπὸ τῆς πορνείας,
\vs{4}~εἰδέναι ἕκαστον ὑμῶν τὸ ἑαυτοῦ σκεῦος κτᾶσθαι ἐν ἁγιασμῷ καὶ τιμῇ,
\vs{5}~μὴ ἐν πάθει ἐπιθυμίας καθάπερ καὶ τὰ ἔθνη τὰ μὴ εἰδότα τὸν θεόν,
\vs{6}~τὸ μὴ ὑπερβαίνειν καὶ πλεονεκτεῖν ἐν τῷ πράγματι τὸν \adeA{ἀδελφὸν} αὐτοῦ, διότι ἔκδικος κύριος περὶ πάντων τούτων, καθὼς καὶ προείπαμεν ὑμῖν καὶ διεμαρτυράμεθα.
\vs{7}~οὐ γὰρ ἐκάλεσεν ἡμᾶς ὁ θεὸς ἐπὶ ἀκαθαρσίᾳ ἀλλ’ ἐν ἁγιασμῷ.
\vs{8}~τοιγαροῦν ὁ ἀθετῶν οὐκ ἄνθρωπον ἀθετεῖ ἀλλὰ τὸν θεὸν τὸν καὶ διδόντα τὸ πνεῦμα αὐτοῦ τὸ ἅγιον εἰς ὑμᾶς.
\vs{9}~Περὶ δὲ τῆς \adeA{φιλαδελφίας} οὐ χρείαν ἔχετε γράφειν ὑμῖν, αὐτοὶ γὰρ ὑμεῖς θεοδίδακτοί ἐστε εἰς τὸ ἀγαπᾶν ἀλλήλους·
\vs{10}~καὶ γὰρ ποιεῖτε αὐτὸ εἰς πάντας τοὺς \adeA{ἀδελφοὺς} τοὺς ἐν ὅλῃ τῇ Μακεδονίᾳ. παρακαλοῦμεν δὲ ὑμᾶς, \adeA{ἀδελφοί}, περισσεύειν μᾶλλον,
\vs{11}~καὶ φιλοτιμεῖσθαι ἡσυχάζειν καὶ πράσσειν τὰ ἴδια καὶ ἐργάζεσθαι ταῖς χερσὶν ὑμῶν, καθὼς ὑμῖν παρηγγείλαμεν,
\vs{12}~ἵνα περιπατῆτε εὐσχημόνως πρὸς τοὺς ἔξω καὶ μηδενὸς χρείαν ἔχητε.
\vs{13}~Οὐ θέλομεν δὲ ὑμᾶς ἀγνοεῖν, \adeA{ἀδελφοί}, περὶ τῶν κοιμωμένων, ἵνα μὴ λυπῆσθε καθὼς καὶ οἱ λοιποὶ οἱ μὴ ἔχοντες ἐλπίδα.
\vs{14}~εἰ γὰρ πιστεύομεν ὅτι Ἰησοῦς ἀπέθανεν καὶ ἀνέστη, οὕτως καὶ ὁ θεὸς τοὺς κοιμηθέντας διὰ τοῦ Ἰησοῦ ἄξει σὺν αὐτῷ.
\vs{15}~τοῦτο γὰρ ὑμῖν λέγομεν ἐν λόγῳ κυρίου, ὅτι ἡμεῖς οἱ ζῶντες οἱ περιλειπόμενοι εἰς τὴν παρουσίαν τοῦ κυρίου οὐ μὴ φθάσωμεν τοὺς κοιμηθέντας·
\vs{16}~ὅτι αὐτὸς ὁ κύριος ἐν κελεύσματι, ἐν φωνῇ ἀρχαγγέλου καὶ ἐν σάλπιγγι θεοῦ, καταβήσεται ἀπ’ οὐρανοῦ, καὶ οἱ νεκροὶ ἐν Χριστῷ ἀναστήσονται πρῶτον,
\vs{17}~ἔπειτα ἡμεῖς οἱ ζῶντες οἱ περιλειπόμενοι ἅμα σὺν αὐτοῖς ἁρπαγησόμεθα ἐν νεφέλαις εἰς ἀπάντησιν τοῦ κυρίου εἰς ἀέρα· καὶ οὕτως πάντοτε σὺν κυρίῳ ἐσόμεθα.
\vs{18}~ὥστε παρακαλεῖτε ἀλλήλους ἐν τοῖς λόγοις τούτοις.

\chapheader{5}
\vs{1}~Περὶ δὲ τῶν χρόνων καὶ τῶν καιρῶν, \adeA{ἀδελφοί}, οὐ χρείαν ἔχετε ὑμῖν γράφεσθαι,
\vs{2}~αὐτοὶ γὰρ ἀκριβῶς οἴδατε ὅτι ἡμέρα κυρίου ὡς κλέπτης ἐν νυκτὶ οὕτως ἔρχεται.
\vs{3}~ὅταν λέγωσιν· Εἰρήνη καὶ ἀσφάλεια, τότε αἰφνίδιος αὐτοῖς ἐφίσταται ὄλεθρος ὥσπερ ἡ ὠδὶν τῇ ἐν γαστρὶ ἐχούσῃ, καὶ οὐ μὴ ἐκφύγωσιν.
\vs{4}~ὑμεῖς δέ, \adeA{ἀδελφοί}, οὐκ ἐστὲ ἐν σκότει, ἵνα ἡ ἡμέρα ὑμᾶς ὡς κλέπτης καταλάβῃ,
\vs{5}~πάντες γὰρ ὑμεῖς υἱοὶ φωτός ἐστε καὶ υἱοὶ ἡμέρας. οὐκ ἐσμὲν νυκτὸς οὐδὲ σκότους·
\vs{6}~ἄρα οὖν μὴ καθεύδωμεν ὡς οἱ λοιποί, ἀλλὰ γρηγορῶμεν καὶ νήφωμεν.
\vs{7}~οἱ γὰρ καθεύδοντες νυκτὸς καθεύδουσιν, καὶ οἱ μεθυσκόμενοι νυκτὸς μεθύουσιν·
\vs{8}~ἡμεῖς δὲ ἡμέρας ὄντες νήφωμεν, ἐνδυσάμενοι θώρακα πίστεως καὶ ἀγάπης καὶ περικεφαλαίαν ἐλπίδα σωτηρίας·
\vs{9}~ὅτι οὐκ ἔθετο ἡμᾶς ὁ θεὸς εἰς ὀργὴν ἀλλὰ εἰς περιποίησιν σωτηρίας διὰ τοῦ κυρίου ἡμῶν Ἰησοῦ Χριστοῦ,
\vs{10}~τοῦ ἀποθανόντος περὶ ἡμῶν ἵνα εἴτε γρηγορῶμεν εἴτε καθεύδωμεν ἅμα σὺν αὐτῷ ζήσωμεν.
\vs{11}~διὸ παρακαλεῖτε ἀλλήλους καὶ οἰκοδομεῖτε εἷς τὸν ἕνα, καθὼς καὶ ποιεῖτε.
\vs{12}~Ἐρωτῶμεν δὲ ὑμᾶς, \adeA{ἀδελφοί}, εἰδέναι τοὺς κοπιῶντας ἐν ὑμῖν καὶ προϊσταμένους ὑμῶν ἐν κυρίῳ καὶ νουθετοῦντας ὑμᾶς,
\vs{13}~καὶ ἡγεῖσθαι αὐτοὺς ὑπερεκπερισσοῦ ἐν ἀγάπῃ διὰ τὸ ἔργον αὐτῶν. εἰρηνεύετε ἐν ἑαυτοῖς.
\vs{14}~παρακαλοῦμεν δὲ ὑμᾶς, \adeA{ἀδελφοί}, νουθετεῖτε τοὺς ἀτάκτους, παραμυθεῖσθε τοὺς ὀλιγοψύχους, ἀντέχεσθε τῶν ἀσθενῶν, μακροθυμεῖτε πρὸς πάντας.
\vs{15}~ὁρᾶτε μή τις κακὸν ἀντὶ κακοῦ τινι ἀποδῷ, ἀλλὰ πάντοτε τὸ ἀγαθὸν διώκετε εἰς ἀλλήλους καὶ εἰς πάντας.
\vs{16}~πάντοτε χαίρετε,
\vs{17}~ἀδιαλείπτως προσεύχεσθε,
\vs{18}~ἐν παντὶ εὐχαριστεῖτε· τοῦτο γὰρ θέλημα θεοῦ ἐν Χριστῷ Ἰησοῦ εἰς ὑμᾶς.
\vs{19}~τὸ πνεῦμα μὴ σβέννυτε,
\vs{20}~προφητείας μὴ ἐξουθενεῖτε·
\vs{21}~πάντα δὲ δοκιμάζετε, τὸ καλὸν κατέχετε,
\vs{22}~ἀπὸ παντὸς εἴδους πονηροῦ ἀπέχεσθε.
\vs{23}~Αὐτὸς δὲ ὁ θεὸς τῆς εἰρήνης ἁγιάσαι ὑμᾶς ὁλοτελεῖς, καὶ ὁλόκληρον ὑμῶν τὸ πνεῦμα καὶ ἡ ψυχὴ καὶ τὸ σῶμα ἀμέμπτως ἐν τῇ παρουσίᾳ τοῦ κυρίου ἡμῶν Ἰησοῦ Χριστοῦ τηρηθείη.
\vs{24}~πιστὸς ὁ καλῶν ὑμᾶς, ὃς καὶ ποιήσει.
\vs{25}~\adeA{Ἀδελφοί}, προσεύχεσθε περὶ ἡμῶν.
\vs{26}~ἀσπάσασθε τοὺς \adeA{ἀδελφοὺς} πάντας ἐν φιλήματι ἁγίῳ.
\vs{27}~ἐνορκίζω ὑμᾶς τὸν κύριον ἀναγνωσθῆναι τὴν ἐπιστολὴν πᾶσιν τοῖς \adeA{ἀδελφοῖς}.
\vs{28}~ἡ χάρις τοῦ κυρίου ἡμῶν Ἰησοῦ Χριστοῦ μεθ’ ὑμῶν.

\booksection{Philemon}{Πρὸς Φιλήμονα}

\chapheader{1}
\vs{1}~Παῦλος δέσμιος Χριστοῦ Ἰησοῦ καὶ Τιμόθεος ὁ \adeA{ἀδελφὸς} Φιλήμονι τῷ ἀγαπητῷ καὶ συνεργῷ ἡμῶν
\vs{2}~καὶ Ἀπφίᾳ τῇ \adeA{ἀδελφῇ} καὶ Ἀρχίππῳ τῷ συστρατιώτῃ ἡμῶν καὶ τῇ κατ’ οἶκόν σου ἐκκλησίᾳ·
\vs{3}~χάρις ὑμῖν καὶ εἰρήνη ἀπὸ θεοῦ πατρὸς ἡμῶν καὶ κυρίου Ἰησοῦ Χριστοῦ.
\vs{4}~Εὐχαριστῶ τῷ θεῷ μου πάντοτε μνείαν σου ποιούμενος ἐπὶ τῶν προσευχῶν μου,
\vs{5}~ἀκούων σου τὴν ἀγάπην καὶ τὴν πίστιν ἣν ἔχεις πρὸς τὸν κύριον Ἰησοῦν καὶ εἰς πάντας τοὺς ἁγίους,
\vs{6}~ὅπως ἡ κοινωνία τῆς πίστεώς σου ἐνεργὴς γένηται ἐν ἐπιγνώσει παντὸς ἀγαθοῦ τοῦ ἐν ἡμῖν εἰς Χριστόν·
\vs{7}~χαρὰν γὰρ πολλὴν ἔσχον καὶ παράκλησιν ἐπὶ τῇ ἀγάπῃ σου, ὅτι τὰ σπλάγχνα τῶν ἁγίων ἀναπέπαυται διὰ σοῦ, \adeA{ἀδελφέ}.
\vs{8}~Διό, πολλὴν ἐν Χριστῷ παρρησίαν ἔχων ἐπιτάσσειν σοι τὸ ἀνῆκον,
\vs{9}~διὰ τὴν ἀγάπην μᾶλλον παρακαλῶ, τοιοῦτος ὢν ὡς Παῦλος πρεσβύτης νυνὶ δὲ καὶ δέσμιος Χριστοῦ Ἰησοῦ—
\vs{10}~παρακαλῶ σε περὶ τοῦ ἐμοῦ τέκνου, ὃν ἐγέννησα ἐν τοῖς δεσμοῖς Ὀνήσιμον,
\vs{11}~τόν ποτέ σοι ἄχρηστον νυνὶ δὲ σοὶ καὶ ἐμοὶ εὔχρηστον,
\vs{12}~ὃν ἀνέπεμψά σοι αὐτόν, τοῦτ’ ἔστιν τὰ ἐμὰ σπλάγχνα·
\vs{13}~ὃν ἐγὼ ἐβουλόμην πρὸς ἐμαυτὸν κατέχειν, ἵνα ὑπὲρ σοῦ μοι διακονῇ ἐν τοῖς δεσμοῖς τοῦ εὐαγγελίου,
\vs{14}~χωρὶς δὲ τῆς σῆς γνώμης οὐδὲν ἠθέλησα ποιῆσαι, ἵνα μὴ ὡς κατὰ ἀνάγκην τὸ ἀγαθόν σου ᾖ ἀλλὰ κατὰ ἑκούσιον·
\vs{15}~τάχα γὰρ διὰ τοῦτο ἐχωρίσθη πρὸς ὥραν ἵνα αἰώνιον αὐτὸν ἀπέχῃς,
\vs{16}~οὐκέτι ὡς δοῦλον ἀλλὰ ὑπὲρ δοῦλον, \adeA{ἀδελφὸν} ἀγαπητόν, μάλιστα ἐμοί, πόσῳ δὲ μᾶλλον σοὶ καὶ ἐν σαρκὶ καὶ ἐν κυρίῳ.
\vs{17}~Εἰ οὖν με ἔχεις κοινωνόν, προσλαβοῦ αὐτὸν ὡς ἐμέ.
\vs{18}~εἰ δέ τι ἠδίκησέν σε ἢ ὀφείλει, τοῦτο ἐμοὶ ἐλλόγα·
\vs{19}~ἐγὼ Παῦλος ἔγραψα τῇ ἐμῇ χειρί, ἐγὼ ἀποτίσω· ἵνα μὴ λέγω σοι ὅτι καὶ σεαυτόν μοι προσοφείλεις.
\vs{20}~ναί, \adeA{ἀδελφέ}, ἐγώ σου ὀναίμην ἐν κυρίῳ· ἀνάπαυσόν μου τὰ σπλάγχνα ἐν Χριστῷ.
\vs{21}~πεποιθὼς τῇ ὑπακοῇ σου ἔγραψά σοι, εἰδὼς ὅτι καὶ ὑπὲρ ἃ λέγω ποιήσεις.
\vs{22}~ἅμα δὲ καὶ ἑτοίμαζέ μοι ξενίαν, ἐλπίζω γὰρ ὅτι διὰ τῶν προσευχῶν ὑμῶν χαρισθήσομαι ὑμῖν.
\vs{23}~Ἀσπάζεταί σε Ἐπαφρᾶς ὁ συναιχμάλωτός μου ἐν Χριστῷ Ἰησοῦ,
\vs{24}~Μᾶρκος, Ἀρίσταρχος, Δημᾶς, Λουκᾶς, οἱ συνεργοί μου.
\vs{25}~Ἡ χάρις τοῦ κυρίου Ἰησοῦ Χριστοῦ μετὰ τοῦ πνεύματος ὑμῶν.
